\documentclass[9pt,twocolumn]{article}

\usepackage[a4paper,top=1.35cm,bottom=1.45cm,left=1.45cm,right=1.45cm,columnsep=0.58cm]{geometry}
\usepackage[T1]{fontenc}
\usepackage[utf8]{inputenc}
\usepackage{mathptmx}
\usepackage[scaled=0.92]{helvet}
\usepackage{amsmath,amssymb,mathtools}
\usepackage{booktabs,array,tabularx}
\usepackage{microtype}
\usepackage{enumitem}
\usepackage{xcolor}
\usepackage{titlesec}
\usepackage{fancyhdr}
\usepackage{graphicx}
\usepackage{caption}
\usepackage{balance}
\usepackage{url}
\usepackage{hyperref} % v5: loaded last

\definecolor{RRHMNavy}{HTML}{17263A}
\definecolor{RRHMSlate}{HTML}{5A6877}
\definecolor{RRHMBlue}{HTML}{315D86}
\definecolor{RRHMLight}{HTML}{E9EEF3}

\hypersetup{
  colorlinks=true,
  linkcolor=RRHMNavy,
  citecolor=RRHMBlue,
  urlcolor=RRHMBlue,
  pdfauthor={Yaoharee Lahtee},
  pdftitle={Why We Phobia? The Regulatory Recoverability Horizon Model of Specific Phobia -- Adversarial Review Edition v21 (Readout-Anchored, Battery-Complete, Coupled, Fenced)},
  pdfsubject={Medical Hypothesis / Translational Systems Neuroscience},
  pdfkeywords={specific phobia, threat imminence, controllability, exposure therapy, recoverability, psychophysiology, computational psychiatry, readout, falsification}
}

\setlength{\parindent}{1em}
\setlength{\parskip}{0.15em}
\setlength{\columnsep}{0.58cm}
\setlength{\emergencystretch}{1.5em}
\setlist{nosep,leftmargin=1.3em}
\raggedbottom
\vfuzz=14pt
\sloppy
\allowdisplaybreaks

\titleformat{\section}
  {\sffamily\bfseries\color{RRHMNavy}\large}
  {\thesection}{0.45em}{}
\titleformat{\subsection}
  {\sffamily\bfseries\color{RRHMNavy}\normalsize}
  {\thesubsection}{0.42em}{}
\titleformat{\subsubsection}
  {\sffamily\bfseries\color{RRHMSlate}\small}
  {\thesubsubsection}{0.40em}{}
\titlespacing*{\section}{0pt}{0.95em}{0.34em}
\titlespacing*{\subsection}{0pt}{0.65em}{0.22em}
\titlespacing*{\subsubsection}{0pt}{0.50em}{0.18em}

\pagestyle{fancy}
\fancyhf{}
\fancyhead[L]{\sffamily\bfseries\scriptsize\color{RRHMNavy} PREPRINT \textbullet\ NOT PEER REVIEWED}
\fancyhead[R]{\sffamily\scriptsize\color{RRHMSlate} WHY WE PHOBIA? \textbullet\ RRHM \textbullet\ MEDICAL EDITION v21}
\fancyfoot[C]{\sffamily\scriptsize\color{RRHMSlate}\thepage}
\renewcommand{\headrulewidth}{0.3pt}
\renewcommand{\footrulewidth}{0pt}

\captionsetup{font=small,labelfont={bf,sf,color=RRHMNavy}}
\renewcommand{\arraystretch}{1.08}

% ---- v5 macros: multi-letter symbols set upright (v4 rendered e*R*R*H) ----
\newcommand{\eRRH}{\mathrm{eRRH}}
\newcommand{\dRRH}{\mathrm{dRRH}}
\newcommand{\eRRHc}{\mathrm{eRRH}_{C}}
\newcommand{\dRRHc}{\mathrm{dRRH}_{C}}
\newcommand{\eRRHhat}{\widehat{\mathrm{eRRH}}}
\newcommand{\dRRHhat}{\widehat{\mathrm{dRRH}}}
\newcommand{\RCE}{\mathrm{RCE}}
\newcommand{\PRHC}{\mathrm{PRHC}}
\newcommand{\Tdt}{T_{\mathrm{DT}}}
\newcommand{\tcorr}{\widehat{\tau}^{E}_{\mathrm{corr}}}
\newcommand{\tier}[1]{{\sffamily\scriptsize\color{RRHMSlate}[#1]}}
\newcommand{\Vset}{\mathcal{V}}
\newcommand{\Hset}{\mathcal{H}}

\title{\vspace{-1.1em}
{\sffamily\bfseries\fontsize{19}{21}\selectfont\color{RRHMNavy} Why We Phobia?}\\[0.35em]
{\sffamily\bfseries\fontsize{12.0}{14}\selectfont The Regulatory Recoverability Horizon Model of Specific Phobia}\\[0.18em]
{\sffamily\fontsize{9.2}{11}\selectfont\color{RRHMSlate} Adversarial Review Edition v21 (Readout-Anchored, Battery-Complete, Coupled, Fenced): A Falsifiable Translational Systems-Neuroscience Hypothesis}}
\author{\sffamily\bfseries Yaoharee Lahtee\\[0.2em]
{\sffamily\mdseries\small Independent Researcher, Open Civil Science Initiative, Bangkok--Pattani, Thailand}\\
{\sffamily\mdseries\small ORCID: \href{https://orcid.org/0009-0005-3861-0626}{0009-0005-3861-0626}}}
\date{\sffamily\small 2 September 2026}

\begin{document}
\maketitle
\thispagestyle{fancy}

\begin{center}
\small\sffamily\color{RRHMSlate}
Article type: Medical Hypothesis / Translational Systems-Neuroscience Concept Paper\\
Adversarial Falsification Edition (v21, readout-anchored, battery-complete, coupled, fenced)
\end{center}

\vspace{0.15em}
\begin{center}
\fcolorbox{RRHMSlate}{RRHMLight}{%
\parbox{0.94\columnwidth}{\footnotesize\sffamily\color{RRHMNavy}
\textbf{Preprint.} This manuscript has not been peer reviewed. It reports no
new human or animal data. Version 20, 2 September 2026. Correspondence:
\href{mailto:arayawedding@gmail.com}{arayawedding@gmail.com}. Distributed under
CC~BY~4.0. Nothing in this document is clinical guidance; see \S\ref{sec:limits}.}}
\end{center}

\vspace{0.35em}
\begin{center}
\parbox{0.95\columnwidth}{\footnotesize\sffamily\color{RRHMNavy}
\textbf{\color{RRHMSlate}Key points}\\[0.25em]
\textbf{Question.} Does a predicted deadline---the time remaining in which
continued engagement can still be made regulatable---govern when a person
switches to defensive termination, independently of threat, escape
availability, control, and the speed of the internal clock?\\[0.2em]
\textbf{Proposal.} Specific phobia involves a persistent, cue-locked
\emph{additive} bias of that estimate, separable from the escape, imminence,
control and timing channels.\\[0.2em]
\textbf{What is new here.} Of the eleven falsification gates this hypothesis
commits to, ten are simulable and all ten have now been run before any
participant was recruited (K1 is a procedural rule, not a simulation target).
Six were run in v6--v8: four of those failed, the design and one construct's
algebraic form were changed, and the repaired battery discriminates. The
remaining four (K3, K6, K8, K9) were run at v9: one usable as written, two
repaired, one uninformative in its timing-only form. This edition (v20) then
passed the whole corpus through an independent six-dimension adversarial
review (v20); the confirmed findings --- including withdrawal of one open-data
result as an artifact --- are dispositioned in the version history. This
edition (v21) adds a second round of preregistered open-data contacts
(C25--C28, four further datasets, criteria frozen before access), reported as
a challenge ledger: lane-specific individual structure proved longitudinally
stable across six months while same-day cross-lane convergence stayed weak;
one frozen secondary prediction failed, and the interpretation it carried ---
that the rating lane is the privileged stable carrier --- is withdrawn in the
text. The data changed the reading; the reading was not bent to the data.}
\end{center}

\begin{abstract}
\textbf{Background:} Specific phobia is defined clinically by persistent, disproportionate fear or anxiety toward a circumscribed cue, accompanied by avoidance or endurance with intense distress and functional impairment. Conditioning, extinction, threat-imminence neuroscience, escape-decision models, controllability, interoception, active inference, and exposure therapy explain major parts of the phenotype, yet three problems remain incompletely unified: physiologically heterogeneous phobias converge on similar defensive termination; only some ordinary fears become persistent disorders; and remission is variably durable.

\textbf{Concept:} The Regulatory Recoverability Horizon Model (RRHM) proposes that the critical latent variable is not fear intensity itself but the estimated remaining interval during which continued engagement can still be preserved or restored by effective corrective action. This engagement-preserving Regulatory Recoverability Horizon (eRRH) is separated from escape availability (dRRH). Specific phobia is hypothesized to involve cue-locked premature contraction of estimated eRRH relative to causal eRRH, termed Premature Regulatory Horizon Contraction (PRHC).

\textbf{Anchoring:} This edition states the model as an instance of a general record law---an estimate is a translation of a state, never the state itself, and the translation gap never closes to zero---so that PRHC is a cue-locked \emph{bias} of a horizon channel rather than the mere presence of error. The anchoring supplies a minimal generative model (the anchor--drive model), an identifiability gate that forces a correction-latency module, a distinction between access-based and assumption-based recalibration that yields a new precommitted prediction, and a typing rule under which every kill criterion must be able to fail.

\textbf{Extension (Dr):} One added term---coupling to an external calm anchor, $\dot d=T-(v+g_0+k)d$---derives four clinical facts of the companion phenomenon in agoraphobia from the same equation: horizon extension, rescue at $v=0$, bereavement onset as deletion of a shared term (a generic, non-cue-locked contraction, forcing a double dissociation against specific phobia's cue-locked bias), and why companion fading works while abrupt removal does not. An empirical fence assembled around the extension (S10) forced the companion term into three separable parts --- co-regulation $k$ (presence, channel-borne, readable in physiological recovery slopes), capacity $\Delta g_{ext}$, and a declared observer term --- and extracted a new problem: the portable anchor (coupling should order by channel bandwidth, not by symbolic meaning; the ladder from in-person to talisman is its test). The extension carries its own gates (K12: capacity versus matched felt safety, shown two-sided by simulation with a declared effect floor; C18: recovery slopes), is fenced outside the core scope, asserts nothing about rival theories whose primary texts remain unread, and names anticipatory dread and the social-observer channel as open terms it does not possess.

\textbf{Adversarial objective:} The theory is treated as guilty until discriminating evidence is produced. This edition specifies failure modes, seven rival models including generic interval timing and escape-decision models, anti-circular measurement rules, subtype counterexamples with their strongest non-RRHM accounts, transdiagnostic challenges, and eleven numeric kill gates (K1--K11) each carrying a positive and a negative control. Ten of the eleven can be tested by simulation before any participant is recruited (K1 is a procedural ordering rule, not a simulation target), and all ten have now been run on synthetic cohorts. Six were run first, calibrated to published escape-decision benchmarks: four of the six failed their own controls, and the manuscript was revised accordingly---the pathological bias is committed to its additive form, mid-trial horizon probes and an independent correction-latency measure are made mandatory, and two thresholds are restated; the repaired battery was re-run, and the gate that had been void in both directions now discriminates cleanly. The remaining four (K3, K6, K8, K9) were then run on an independent re-implementation of the generative model: K8 is usable as written; K3 is void on observed correlations at any feasible sample size and is repaired by reliability disattenuation, becoming two-sided at $N\ge220$; K9 requires the same hierarchical two-standard-error rule that repaired K5, and its fixed $\Delta$ELPD threshold must be preregistered jointly with $N$; K6 is uninformative at feasible sample sizes in its timing-only form and is committed, untested, to the joint hazard-plus-probe form.

\textbf{Proposed methods:} A safe $2\times2\times2$ deadline-by-exit-by-instruction paradigm in which both available actions are equally effective and only the causal deadline differs; prospective explicit, choice-derived and interval-production horizon estimates probed \emph{repeatedly within a trial}; a mandatory correction-latency module; multimodal psychophysiology; independent transition outcomes; parameter recovery and design-level Fisher information computed under the generative model before recruitment; hierarchical hazard models with subject-level random effects; and preregistered held-out comparison against conditioning/extinction, threat imminence, escape-decision, perceived control, allostatic self-efficacy, interval-timing and active-inference models.

\textbf{Conclusion:} RRHM should survive only if eRRH can be independently measured before defensive transition, remains identifiable after rival variables are modeled, is not a special case of generic timing distortion, predicts transition timing across phobia classes, and prospectively predicts disease-course outcomes. The simulation reported here shows that the first two of those conditions are properties of the measurement design rather than of the theory, and that a deadline-only design cannot decide them at any sample size. Until the revised battery is run, RRHM is an experiment-ready medical hypothesis rather than an established mechanism.
\end{abstract}

\noindent\textbf{Keywords:} specific phobia; psychophysiology; systems neuroscience; avoidance; exposure therapy; threat imminence; escape decisions; controllability; recoverability; computational psychiatry; readout; falsification.

\subsection*{Abbreviations}
\noindent\footnotesize
ADM, anchor--drive model; ASE, allostatic self-efficacy; BII, blood-injection-injury;
$\dRRH$, defensive-termination recoverability horizon; ELPD, expected log predictive
density; $\eRRH$, engagement-preserving regulatory recoverability horizon; FID, flight
initiation distance; IRHAB, Independent Recoverability Horizon Assessment Battery;
PRHC, Premature Regulatory Horizon Contraction; $\RCE$, recoverability calibration
error; RRHM, Regulatory Recoverability Horizon Model; TOST, two one-sided tests.
\normalsize

\begin{table}[!t]
\centering
\caption{Notation used throughout. Every symbol is defined once, in \S\ref{sec:formal} or \S\ref{sec:adm}, and used with that meaning only.}
\label{tab:notation}
\small
\begin{tabular}{@{}lp{0.66\columnwidth}@{}}
\toprule
Symbol & Meaning \\
\midrule
$x_t$ & organism--environment state relevant to continued engagement \\
$\Vset,\ \Hset$ & viability region (staying) and harm-free region (leaving) \\
$\Pi^E_t,\ \Pi^D_t$ & engagement-preserving and defensive-termination policies \\
$\eRRHc,\ \dRRHc$ & causal horizons for staying and for leaving \\
$\eRRHhat,\ \dRRHhat$ & the organism's estimates of those horizons \\
$\RCE_t$ & $\eRRHc-\eRRHhat$: the translation gap \\
$\tcorr$ & estimated latency to execute an effective correction \\
$M^E_t$ & $\eRRHhat-\tcorr$: the transition margin \\
$\Tdt$ & time of defensive transition \\
$D_j$ & externally programmed deadline on trial $j$ \\
$K_{cue},\ \eta_t$ & cue-conditioned gain and additive term of the horizon channel \\
$c$ & internal-clock gain of the rival timing model \\
$d,\ v,\ g,\ T$ & ADM: deviation, intrinsic damping, corrective coupling, drive \\
$\beta_{0i},\ \beta_{1i}$ & subject-level hazard intercept and margin slope \\
\bottomrule
\end{tabular}
\end{table}

\section{Background: the medical problem before the theory}
Specific phobia is common, frequently begins early, and may persist for years despite being tied to circumscribed cues. World Mental Health Survey data estimated lifetime prevalence at approximately 7.4\% and 12-month prevalence at 5.5\%, with substantial impairment and comorbidity \cite{wardenaar2017,wardenaar2018corr}. Exposure-based treatments are effective, including efficient single-session formats \cite{odgers2022}, but treatment efficacy does not by itself identify the upstream variable that governs pathological defensive transition.

Three clinical problems motivate RRHM. First, animal phobia, blood-injection-injury (BII) phobia, height phobia, situational phobia, and procedural or visceral phobias differ in stimulus properties, physiology, true hazard, and sensorimotor constraints yet converge on escape, freezing, avoidance, or distressed endurance. Second, ordinary fear is adaptive, whereas only some cue-linked fears become persistent disorders. Third, remission and relapse are incompletely predicted by conventional fear and extinction measures \cite{trumpf2009remission,devos2025,bohnlein2020}.

The proposed unifying question is therefore not simply ``why is the cue feared?'' but:
\begin{quote}
\textbf{When does the organism predict that continued engagement will cease to remain recoverable, and does that predicted deadline govern defensive transition independently of conventional threat, escape, timing, and control variables?}
\end{quote}
Of the three problems, RRHM addresses the first and third directly. For the second it supplies a maintenance mechanism (avoidance closes the sampling that would recalibrate the estimate; \S\ref{sec:course}) but not a diathesis; that gap is stated as an open problem in \S\ref{sec:limits} rather than implied to be solved.

\section{Clinical boundary and medical mimic gate}
RRHM is not a diagnostic substitute. Palpitations, dyspnea, dizziness, presyncope, syncope, weakness, paresthesia, swallowing difficulty, altered awareness, or chest discomfort can arise from cardiovascular, respiratory, neurological, vestibular, metabolic, autonomic, medication-related, or substance-related conditions. True physiological danger must be evaluated when clinically indicated before a defensive response is interpreted as predictive miscalibration.

The model distinguishes:
\begin{enumerate}
\item appropriate defense to genuine hazard;
\item excessive predicted contraction despite preserved causal recoverability;
\item mixed cases containing real physiological vulnerability plus predictive underestimation.
\end{enumerate}

The differential is not only somatic: the psychiatric differential of DSM-5-TR --- panic disorder, agoraphobia, social anxiety disorder, obsessive-compulsive disorder, posttraumatic stress disorder, illness anxiety disorder, avoidant/restrictive food intake disorder (ARFID), and psychotic disorders --- must be considered before a presentation is classed as specific phobia, and nothing in this manuscript performs that adjudication. This distinction is especially important in BII phobia, where vasovagal susceptibility may shorten actual circulatory recoverability, and in height-related disorders, where postural and vestibular instability may be genuine contributors \cite{ayala2009,ritz2010,huppert2020}. In the generative model of \S\ref{sec:adm}, case 1 is a genuinely short causal horizon, case 2 is a biased estimate of an adequate causal horizon, and case 3 is both.

\section{Epistemic anchoring: the horizon as a readout channel}
\label{sec:anchor}
Every quantitative claim in this paper is stated under one general principle, developed elsewhere by the author as a philosophy and logic for grounding claims \cite{lahtee2026ru,lahtee2026rc} and imported here without any inheritance of formal status: \emph{an estimate is a translation of a state, never the state itself, and the translation gap never closes to zero.} Written for the present construct,
\begin{equation}
\eRRHhat_t \;=\; K_{cue}\,\eRRHc(x_t)\;+\;\eta_t,\qquad \varepsilon_{tot}:=\mathbb{E}\,|\eta_t|>0\ \ \forall t,
\label{eq:recordlaw}
\end{equation}
where $\eta_t$ collects selection noise (which state variables are sampled), mapping noise (how they are converted into a temporal estimate), and self-referential noise (the estimate's dependence on the estimator's own state), and $K_{cue}$ is a cue-conditioned gain. A caution that turns out to be decisive: the gain and the additive term are \emph{not} interchangeable ways of writing the same pathology. Locating the bias in $K_{cue}$ makes it formally indistinguishable from an internal-clock change; locating it in $\mathbb{E}[\eta\mid cue]$ does not (\S\ref{sec:sim}). RRHM therefore places PRHC in $\eta$. Three consequences shape everything that follows.

\begin{enumerate}[leftmargin=1.5em]
\item \textbf{Bias, not error, is the pathology.} Because $\varepsilon_{tot}>0$ always, a nonzero calibration error is the healthy default. RRHM therefore predicts residual miscalibration in every participant, including remitters; what distinguishes phobia is a persistent, cue-conditioned \emph{directional} bias (\S\ref{sec:prhc}).
\item \textbf{Access recalibrates; assumption does not.} When two distinct causal horizons produce the same estimate, no post-processing of the estimate recovers the distinction; only additional access---sampling the state near the boundary---can. Declared assumptions (``this is safe'', ``you can stop at any time'') can sharpen an estimate but cannot restore a distinction the channel never carried \cite{lahtee2026rc}. This yields a precommitted dissociation between told and learned deadlines (\S\ref{sec:decisive}).
\item \textbf{Identifiability is a property of the record set, not of the theory.} Whether $\eRRH$ can be separated from control, timing or correction latency is decided by which records the measurement battery returns. A direction in parameter space that no record informs is structurally unanswerable from that design; the battery must be built so that no theory-critical direction is uninformed (\S\ref{sec:irhab}).
\end{enumerate}

\subsection{Translation table}
Table~\ref{tab:translation} states each RRHM construct in the anchoring vocabulary, together with the status each statement carries. The point of the table is auditability: a reader can see which statements are stipulations, which are declared bridges with a falsifier, and which are open.

\begin{table*}[!t]
\centering
\caption{Translation of RRHM constructs into the readout vocabulary of \cite{lahtee2026ru}, with the status each statement carries in this paper. \emph{definition}: a stipulated object or rule. \emph{Dr}: a declared bridge or reading, always accompanied by a falsifier. \emph{Open}: not established; carries a stance and a falsifier. No RRHM claim carries a machine-checked status.}
\label{tab:translation}
\small
\begin{tabularx}{\textwidth}{p{0.19\textwidth}X p{0.20\textwidth}}
\toprule
\textbf{RRHM construct} & \textbf{Readout-language translation} & \textbf{Status / consequence} \\
\midrule
$\eRRHc$ (programmed deadline) & The world-side record: the state that the estimate is \emph{about}. In the laboratory it is programmed by design and returned to the analyst as a record independent of behavior. & definition; anti-circularity rule \S\ref{sec:irhab} \\
$\eRRHhat$ & The organism's own record of that state: $K_{cue}\,\theta+\eta$ (Eq.~\ref{eq:recordlaw}); latent, not required to be verbal. & definition \\
$\RCE$ & The translation gap between record and state; strictly positive in expectation. & definition; healthy $\neq$ zero error \\
$\PRHC$ & A cue-conditioned \emph{additive bias} of the translation: $\mathbb{E}[\eta\mid cue]\le-\delta_{cue}<0$, persistent across occasions. The multiplicative form $K_{cue}<1$ is \emph{excluded by decision}, not by evidence: it is exactly likelihood-equivalent to a clock-rate change (\S\ref{sec:sim}). & definition (form fixed); Dr (that phobia has it); falsifier: bias not cue-conditioned (Gate K10) \\
$\Pi^E$ versus $\Pi^D$ & Two anchors in competition: an engagement anchor (restore regulation while staying) and an escape anchor (restore regulation by leaving). & Dr; supplies a two-anchor conflict object \\
$\Tdt$ (defensive transition) & A binary readout jump on a smoothly varying state; the location of the cut lives in the readout policy, not in the state. & Dr; hence no universal biological threshold \\
Avoidance & Closing the record channel for the cue: no new records are returned, so the estimate cannot descend toward the state. & Dr; derives cue-locking \S\ref{sec:cuelock} \\
Remission / relapse & Descent of the calibration residual under new records / re-settlement after an evidence flip; the residual floor never reaches zero. & Dr; predicts nonzero residual in remitters \\
``Point of no return'' & The knife-edge past which no policy returns the state to viability; hand-off into a region the estimator cannot read. & Dr \\
PRHC-U (indeterminacy-triggered contraction, \S\ref{sec:prhcu}) & Zero and unresolved are distinct readout states ($0\neq\bot$); the pathology is their conflation on the horizon channel --- an \emph{unresolved} remaining-time readout is retained as \emph{none remaining}. The healthy estimator holds $\bot$ open (reverts to prior); the affected one collapses it to the floor. & Dr; falsifier: no difference between unpredictable and matched predictable endpoints (\S\ref{sec:prhcu}) \\
Situational-geometry locking (\S\ref{sec:prhcu}) & The bias is indexed by the geometry readout (time-to-exit / time-to-help / time-to-correction), not by perceptual cue features; two perceptually unrelated procedures can share one lesion. & Dr; falsifier: eliciting class tracks stimulus modality, not procedural geometry \\
Recalibration of the $\bot$ branch (\S\ref{sec:prhcu}) & Restore the discipline ``never report $\bot$ as $0$'' inside the estimator: bound the next read ($\bot$ with a known re-read time), build a prior by sampling the indeterminate branch (access, not assumption), or borrow a coupled estimator's prior ($k$). & Dr; three falsifiers stated in \S\ref{sec:prhcu} \\
The recalibration law (\S\ref{sec:recal}) & A retained-record ledger over exact counts: censored records are $\bot$-typed and excluded from the outcome tally (Eq.~\ref{eq:ledger}); the lesion is tallying them as outcomes (Eq.~\ref{eq:lesion}); the cure is three ledger operations --- re-type, append completions to $n^{*}$ (Eq.~\ref{eq:nstar}), cover the gap with borrowed counts and a bounded re-read. Toy-exact over $\mathbb Q$; clinical bridge Dr. & Dr; kill criteria in \S\ref{sec:recal} \\
$\sup\tau=+\infty$ & A non-readout: an actual infinity injected into a finite paradigm; must be censored at $\tau_{\max}$. & definition (convention C2, \S\ref{sec:formal}) \\
\bottomrule
\end{tabularx}
\end{table*}

\subsection{Status discipline used in this paper}
Each substantive statement below is tagged where its status could be mistaken: \tier{definition} for stipulated objects and rules; \tier{Dr} for declared bridges or readings, each carrying a falsifier; \tier{Open} for statements not yet established; \tier{finite\_diagnostic} for a number produced by an executed, reproducible run (evidence, not proof). The tag \tier{machine-checked} is not used for any RRHM statement: the only machine-checked objects touched by this paper are borrowed (a sorites-jump core and a discrete energy-balance algebra recorded in the author's companion repositories), and neither transfers its status to a claim about phobia. A tag is never upgraded by proximity to a stronger one.

\section{Layered placement: the informational architecture upstream of RRHM}
\label{sec:layers}
This section corrects a layering error that earlier framings of this program invited: starting from ``fear'', ``belief'', or ``memory'' and retrofitting the readout formalism afterwards. The architecture is imported from the author's companion information-epistemic framework \cite{lahtee2026genesis} and is stated here in medical-science terms; it is a structural placement (\tier{Dr} throughout), inherits no formal status, and adds no free parameter to any gate.

\subsection{The processing order}
The primitive is not an emotion but a \emph{retained distinction}. The organism--environment system evolves as a state update over an ordered retained history,
\begin{equation}
S_{n+1}=F\big(S_n,\,u_n,\,c_n,\,T_n\big),
\label{eq:stateupdate}
\end{equation}
where $u_n$ is current input, $c_n$ context, and $T_n$ the ordered record of what has been retained so far --- so the present state is fixed not by current stimulation alone but by stored information \emph{and the path by which it arrived}. Processing is ordered: retention $\to$ structure $\to$ translation $\to$ readout $\to$ meaning $\to$ experience $\to$ memory $\to$ belief. The medical consequence is immediate: subjective fear, autonomic output, verbal belief, and the defensive policy are \emph{parallel, typed readouts} of one evolving informational state --- not sequential stages of a single fear pathway --- and readout lanes are permitted to dissociate.

RRHM is thereby re-placed. It is not a root theory of phobia; it is a \emph{domain quotient} of this architecture: a projection $q$ of the state onto the regulation-relevant variables, followed by one readout, $S_n \to q(S_n) \to \eRRHhat_n$ --- the estimated interval over which engagement remains recoverable. Felt fear is a different readout lane of the same state, downstream, and is neither the input to $\eRRHhat$ nor certified by it.

\subsection{The disease restated at this layer}
\begin{quote}
\textbf{Specific phobia as a path-dependent disorder of informational readout:} retained history and current state are translated into a prematurely terminating defensive policy; subjective fear is one possible phenomenal readout of that process --- neither necessary nor sufficient for it. \tier{Dr}
\end{quote}

\subsection{Four challenge patterns read at the correct layer}
Each pattern below strains a single-scalar fear account; each receives a structural resolution at this layer; and each resolution's honest remainder is stated. Holdings are declared per citation.

\textbf{(V1) Post-organic swallowing phobia without a premorbid safety baseline.} In current nosology this presentation is the choking-phobia form of ARFID (avoidant/restrictive food intake disorder), a distinct DSM-5 diagnosis with its own treatment pathway and nutritional-risk management; the architectural reading below is a mechanism hypothesis inside that pathway, never a substitute for nutritional assessment. Choking/gagging pathology of organic origin (e.g.\ eosinophilic esophagitis, pediatric dysphagia; food refusal in chronic disease is documented \cite{staiano2003}, abstract holding) can leave persistent phobic avoidance after the organic disorder resolves --- in a child who never possessed a ``solid food is safe'' representation to lose. At this layer no safety belief is required: the original episodes wrote retained distinctions (swallow-completed $\sharp$ swallow-disrupted) onto $T$; resolution of the pathology changes the world-state, not the retained structure, and the reader continues to classify solid-food states by the old records. Treatment by \emph{completed swallows} works not by teaching the proposition ``food is safe'' but by appending new continuations that split state classes the reader had merged --- the same completion-dose logic as Eq.~\ref{eq:nstar}. \emph{Remainder:} why one child weights such records more heavily than another is an empirical retention-parameter question this architecture does not answer.

\textbf{(V2) Response heterogeneity in blood-injury-injection phobia.} The same phobic class can produce tachycardic panic, vasovagal collapse, disgust, or --- documented within a single patient --- combined fainting and psychogenic non-epileptic seizures \cite{vonmucke2022} (abstract holding). A one-scalar ``fear level'' erases these differences. At this layer they are parallel lanes read from one informational trajectory: RRHM reads only the recoverability lane; the autonomic quotient reads a circulatory trajectory (the ADM $v(t)\!\to\!0$ cliff of P3); dissociation between lanes is expected, so vasovagal collapse does not contradict a phobic transition --- it is a different readout of the same state. \emph{Remainder:} the architecture permits PNES as an action-regime outcome but does not predict its occurrence conditions; no mechanistic PNES account is claimed.

\textbf{(V3) Immediate, cue-specific abolition of lifelong phobia after focal resection.} Abolition of lifelong arachnophobia after left mesial temporal lobectomy, with fear responses to other stimuli unchanged, is on record \cite{binks2015} (abstract holding). Extinction-based accounts strain: no new learning occurred. At this layer the case is legible because retained information, translation, and readout are distinct objects: a resection can alter the \emph{reader/translation} while leaving the record intact,
\[
T_{\mathrm{spider}}\ \text{unchanged},\qquad O_{\mathrm{new}}(T_{\mathrm{spider}})\neq O_{\mathrm{old}}(T_{\mathrm{spider}}),
\]
yielding instantaneous, cue-specific loss without extinction --- and predicting, testably, that declarative memory of the feared object and responses in other lanes can survive the abolition. \emph{Remainder:} the architecture does not derive \emph{which} anatomical lesion alters \emph{which} reader; no neuroanatomical prediction is claimed, and retrofitting one post hoc is exactly the move this manuscript forbids.

\textbf{(V4) Therapeutic change from masked exposure without reportable perception.} Fear reduction from exposure the patient cannot report seeing is documented across paradigms \cite{siegel2025}. The puzzle ``what cognition was corrected if consciousness saw nothing?'' dissolves at this layer, because phenomenal experience is downstream of the state update, not its gateway: a masked stimulus can enter Eq.~\ref{eq:stateupdate} and alter $T$ while the phenomenal lane reads approximately nothing, and the patient's report ``I saw no spider'' is itself only another readout --- never a certificate that informational contact was zero (an agent's report is a translation of its state, not the state; Eq.~\ref{eq:recordlaw}). The phenomenon needs no label of ``unconscious cognition''; the exact description is \emph{informational state-update without phenomenal readout}.

\subsection{Pre-phenomenal error classes, and the provenance of felt fear}
\label{sec:preerror}
Because experience sits downstream in the processing order, an error can occur \emph{before} anything is felt --- and the experience that follows is then fully real, but produced from a state that was mistranslated before it reached the phenomenal stage. Felt fear is downstream evidence \emph{of a readout}, not direct evidence of the state being read: ``I feel that I am out of options'' is direct evidence of \emph{having that experience}, and only indirect, provenance-dependent evidence about the body or the world. Two consequences follow from the companion account of readout chains \cite{lahtee2026rc} (full paper held on file), stated here without new formalism. First, \emph{upstream monotonicity}: a distinction collapsed at an upstream stage --- for instance, ``remaining time unknown'' merged with ``no remaining time'', the PRHC-U lesion --- cannot be recovered by any downstream processing, including insight; only new access can restore it, which is why telling oneself ``unknown does not mean dangerous'' can fail even when sincerely believed: the distinction may never reach the phenomenal system at all. Second, pathology can be located in (at least) four pre-phenomenal classes, each with a different repair:
\begin{center}
\footnotesize
\begin{tabularx}{\columnwidth}{@{}p{0.155\columnwidth}Xp{0.30\columnwidth}@{}}
\toprule
\textbf{Error class} & \textbf{What goes wrong} & \textbf{Repair family} \\
\midrule
Retention & history mis-typed (censored $\mapsto$ harm) & ledger operations, T1 \\
Access & needed distinction never sampled & completion engineering, T2 \\
Translation & distinction present, mapping biased & recalibration; reader change (V3) \\
Policy & readout becomes action too fast, too hard & interruptibility, operation 6 \\
\bottomrule
\end{tabularx}
\end{center}
Two people can therefore share an upstream lesion and feel differently, or report identical fear (8/10) from entirely different informational lesions --- which is why lane-decomposed measurement (operation 2 of \S\ref{sec:rgif}) precedes any treatment choice. The same placement also reverses a habitual causal arrow: fear and avoidance need not cause one another; they can be \emph{siblings} of one upstream regulatory misread, which then feed back on each other through action and retention --- rendering ``behavior improved before fear'', ``fear reduced but avoidance remains'', and ``belief corrected but the body still responds'' expected patterns rather than treatment contradictions. A formal error-propagation law (which error class produces which phenomenal signature) is stated later, at the author's direction and with explicit degree-of-freedom accounting, in \S\ref{sec:repl}; nothing is added at this point in the argument.

\subsection{What this placement buys, and does not}
It renders jointly non-paradoxical, under one ordering, the set of dissociations that strain fear-centric accounts: fear without conscious appraisal; behavioral change without reportable perception; avoidance persisting after organic cure; phobia abolished with memory plausibly intact; divergent physiological outputs from one phobic class; and subjective fear and action shifting at different times. It buys no new predictive power by itself --- every prediction in this manuscript still lives in the gates and preregistered cases --- but it adds a discipline: any RRHM statement that cannot be written at its correct layer (record, translation, readout, or policy) is mis-layered and should be rejected on those grounds before any data are consulted.

\section{Formal construct}
\label{sec:formal}
\subsection{Organism--environment state}
Let $x_t$ denote the integrated organism--environment state relevant to continued task engagement, potentially including cardiovascular, respiratory, vestibular, postural, visceral, nociceptive, motor, sensory, spatial, communicative, and contextual variables.

Let $\Vset$ denote a preregistered functional viability region. In experimental medicine, viability means preserved safe task-compatible regulation, not merely survival. The boundary must be defined prospectively and independently of whether the participant later escapes. Let $\Hset\supseteq\Vset$ denote the harm-free region: the set of states from which termination of the encounter restores global regulation without injury. $\Vset$ concerns \emph{staying}; $\Hset$ concerns \emph{leaving}.

\subsection{Two policy classes}
Available policies are partitioned as
\[
\Pi_t=\Pi_t^E\cup\Pi_t^D,
\]
where $\Pi_t^E$ contains engagement-preserving policies and $\Pi_t^D$ contains defensive-termination policies. $\Pi_t^E$ includes actions that permit continued contact while maintaining or restoring functional regulation. $\Pi_t^D$ includes escape, full withdrawal, termination, or avoidance. Either class may be empty at a given time; the conventions below make the horizons well defined in that case.

\subsection{Assumption A1: nested recoverability}
\label{sec:a1}
Let $\tau\ge0$ denote the delay before an engagement-preserving corrective policy begins, and let
\begin{align*}
W_t=\Big\{\tau\ge0:\ &\exists\,\pi^E\in\Pi^E_{t+\tau}\ \text{such that engagement}\\[-0.15em]
&\text{until }t+\tau\text{ followed by }\pi^E\\[-0.15em]
&\text{preserves or restores }\Vset\Big\}
\end{align*}
be the set of effective delays. \textbf{A1 (nested recoverability)} states that $W_t$ is downward closed: if correction started at delay $\tau$ preserves $\Vset$, so does correction started at any $\tau'<\tau$. Under A1, $W_t$ is an interval $[0,\tau^*]$ and its supremum is the latest effective delay. \tier{definition}

A1 is an assumption, not a fact, and it fails in identifiable regimes: a biphasic vasovagal trajectory in which applied tension is effective early, ineffective near presyncope, and effective again after spontaneous recovery; respiratory paradigms in which the efficacy of a corrective breath depends on breathing phase; swallowing and dental paradigms in which corrective windows recur periodically. In such regimes $W_t$ is a union of windows, and the scalar horizon must be replaced by the window set itself, with $\eRRHc$ defined as the supremum of the first window. The window form is strictly more expressive: it predicts ``wait for the next window'' behavior that no scalar horizon can express, and a paradigm must preregister whether A1 is assumed or the window form is used. \tier{Dr}; falsifier: a paradigm in which behavior tracks a later window while the first-window scalar predicts termination.

\subsection{Causal horizons}
Under A1 the causal engagement-preserving horizon is
\begin{equation}
\eRRHc(x_t)=\sup W_t,
\label{eq:errhc}
\end{equation}
with three conventions that make Eq.~\ref{eq:errhc} well posed: \textbf{C1} $\sup\varnothing:=0$ (once no corrective policy preserves $\Vset$, the horizon is zero, not undefined); \textbf{C2} $\eRRHc:=\min\{\sup W_t,\tau_{\max}\}$, where $\tau_{\max}$ is the maximal engagement duration of the paradigm (an unbounded supremum is a non-readout: no finite paradigm returns it, and a hazard model fed an infinite margin predicts that no participant ever terminates a safe procedure, which is false); \textbf{C3} $\eRRHc(x_t)$ is a function of the current state and therefore \emph{flows with time}: along an unperturbed trajectory with an externally fixed deadline it decreases at unit rate. \tier{definition}

The defensive-termination horizon is defined symmetrically:
\begin{equation}
\begin{aligned}
\dRRHc(x_t)=\sup\Big\{\tau\ge0:\ &\exists\,\pi^D\in\Pi^D_{t+\tau}\ \text{such that}\\[-0.15em]
&\text{engagement until }t+\tau\text{ followed}\\[-0.15em]
&\text{by }\pi^D\text{ keeps }x\in\Hset\Big\},
\end{aligned}
\label{eq:drrhc}
\end{equation}
with the same conventions. $\dRRHc$ is the latest delay at which leaving still restores global regulation without harm; $\eRRHc$ is the latest delay at which staying can still be made regulatable. The two horizons share state, policies and conventions and differ only in the region and the policy class. The organism's estimates are $\eRRHhat_t$ and $\dRRHhat_t$; both are latent and neither requires verbal awareness.

\subsection{Calibration error as a translation gap}
Define recoverability calibration error
\[
\RCE_t=\eRRHc(x_t)-\eRRHhat_t .
\]
Positive $\RCE$ denotes underestimation of the causal engagement window. By Eq.~\ref{eq:recordlaw}, $\RCE_t=(1-K_{cue})\,\eRRHc-\eta_t$: it is the translation gap of a channel, expected to be nonzero everywhere, and pathological only when its cue-conditioned expectation is displaced.

\subsection{Correction latency and transition margin}
Let $\tcorr$ denote the estimated latency required to execute an effective engagement-preserving correction. Define the estimated engagement margin
\[
M_t^E=\eRRHhat_t-\tcorr .
\]
RRHM predicts increasing defensive-transition hazard as $M_t^E$ approaches or crosses zero. The preregistered hazard model is hierarchical:
\begin{equation}
\begin{aligned}
h_D(t\mid i)&=\sigma\!\left(\beta_{0i}-\beta_{1i}M_{it}^E+\boldsymbol\gamma^{\top}\mathbf z_{it}\right),\\[-0.1em]
\beta_{0i}&\sim\mathcal N(\beta_0,\sigma_0^2),\quad \beta_{1i}\sim\mathcal N(\beta_1,\sigma_1^2),
\end{aligned}
\label{eq:hazard}
\end{equation}
where $\mathbf z_{it}$ contains rival variables such as threat intensity, physical imminence, uncertainty, perceived control, baseline fear, generic timing distortion, and state burden. The choice among defensive policies once transition occurs (escape versus freezing or distressed endurance) is modeled separately as a function of $\dRRHhat_t$. No universal biological threshold is claimed: the transition is a binary readout of a smoothly varying margin, and the cut lives in the participant's readout policy, which is why $\beta_{0i}$ is a random effect. \tier{Dr}

\subsection{Premature Regulatory Horizon Contraction}
\label{sec:prhc}
PRHC denotes persistent cue-linked underestimation of engagement-preserving recoverability. \textbf{The bias is committed to its additive form.} A multiplicative gain $K_{cue}<1$ on the horizon estimate and a multiplicative gain $c>1$ on the internal clock generate the same transition-hazard family exactly, so a theory stated in the multiplicative form is unfalsifiable against generic timing distortion at any sample size (\S\ref{sec:sim}, Fig.~\ref{fig:rank}). Writing the bias additively costs nothing empirically and restores the distinction --- \emph{given within-trial probes} (\S\ref{sec:k11}); on transition timing alone the additive family is nested in the clock family and no comparison can favor it:
\begin{align*}
\mathbb{E}\big[\RCE_t\mid cue\big]&\ \ge\ \delta_{cue}>0\ \text{ across occasions},\\[-0.1em]
\mathbb{E}\big[\RCE_t\mid \neg cue\big]&\ \approx\ 0 .
\end{align*}
Specific phobia is not equated with low eRRH in general. The disease hypothesis requires circumscribed cue locking, persistent avoidance or distressed endurance, disproportionality, and functional impairment; the second condition above is what makes the bias cue-locked rather than generic, and it is what Gate K10 tests.

\subsection{Case-motivated refinement: indeterminacy-triggered contraction (PRHC-U)}
\label{sec:prhcu}
\textbf{Clinical observation (single adult case, informally observed, hypothesis-generating only; Dr).} Across two procedural settings in one adult --- dental treatment, and three venipuncture episodes at different facilities --- panic onset was time-locked not to nociceptive or perceptual stimulus content but to the moment the procedure's \emph{remaining duration became unpredictable}: an initial needle puncture was tolerated without symptoms, whereas repeated failed venipuncture attempts (endpoint no longer estimable) precipitated panic; a painful but eventually successful dorsal-hand draw with no endpoint communication and no coping aid produced the most severe episode; and a briefed, rapidly completed draw with a coping aid offered in advance produced no symptoms at all. Prolonged mouth-opening beyond a recalled 3--5 minute tolerance showed the same structure. (Provenance: the author's own experience, reported first-person and disclosed by deliberate choice; see Patient consent. Differential note: this presentation sits on the specific-phobia / panic-spectrum / BII-vasovagal boundary, and the psychiatric differential --- panic disorder, agoraphobia, blood-injection-injury phobia with vasovagal syncope --- has not been clinically adjudicated.) Relief and prevention tracked \emph{restored endpoint predictability} (pre-procedural briefing, rapid completion, scheduled pauses, a pre-agreed attempt limit) and, in the venipuncture setting, supine positioning --- the latter flagging a possible mixed vasovagal component (mimic-gate case~3), which is a clinical question, not a modeling one.

\textbf{Hypothesis (PRHC-U).} The additive bias of \S\ref{sec:prhc} is potentiated by loss of temporal predictability. Let $o_t\in\{0,1\}$ mark whether the deadline signal is observable (communicated, signposted, or learned). In unaffected estimators, loss of observability reverts the horizon estimate toward its prior, $\eRRHhat_t\to\mathbb E[D-t\mid \text{prior}]$; in the affected estimator it defaults to the floor:
\begin{equation}
\mathbb E\big[\eta_t \mid \text{class}, o_t=0\big]\;\ll\;\mathbb E\big[\eta_t \mid \text{class}, o_t=1\big]\;\le\;0,
\label{eq:prhcu}
\end{equation}
i.e.\ the pathology lives on the $o=0$ branch: \emph{indeterminate} remaining time is registered as \emph{no} remaining time. Two further commitments follow. First, the eliciting class in such cases is defined by procedural parameters --- postural commitment, externally controlled duration, unsignaled endpoint --- rather than by stimulus modality (\emph{situational-geometry locking}); two perceptually unrelated procedures then share one lesion, which is this manuscript's cross-situation geometry claim observed within a single person. Second, every effective mitigation observed reduces to one operation: converting an unbounded remaining time into a finite, readable one.

\textbf{Predictions and kill criteria.} (i) At fixed geometry and fixed objective duration, unpredictable-endpoint trials should show earlier defensive transition than matched predictable-endpoint trials, more so in affected individuals than controls. (ii) Restoring a bounded branch structure (e.g.\ a pre-agreed attempt limit) should abolish the effect without changing the procedure itself. (iii) If transition timing tracks elapsed duration regardless of endpoint observability, PRHC-U is dead and plain PRHC suffices. (iv) If unaffected controls show a floor-default of equal magnitude, the phenomenon is generic uncertainty aversion, not a disease refinement. The IRHAB learned arm already manipulates deadline knowledge; adding an unobservable-$D$ condition is a one-cell extension of the decisive experiment and is the cheapest test of Eq.~\ref{eq:prhcu}.

\textbf{What would recalibrate the $o=0$ branch (Dr, derived, untested).} The therapeutic question PRHC-U poses is not ``how to believe the situation is safe'' but \emph{how to restore the healthy handling of an indeterminate readout} --- holding \emph{unknown} open as unknown instead of collapsing it to \emph{none}. Three routes follow from the estimator structure, each with a falsifier:
\begin{enumerate}[leftmargin=1.6em]
\item \textbf{Bounded re-read scheduling.} An indeterminate endpoint is tolerable when the time of the \emph{next observation} is determinate: ``unknown how long in total, but the next information/pause arrives within $\Delta$'' replaces the missing countdown-to-end with a countdown-to-next-read. The state stays unknown; only its unboundedness is removed. Predicts that a guaranteed information schedule (staff updates every $\Delta$; the attempt-limit rule) prevents collapse \emph{without any endpoint information at all} --- distinguishing it from simple endpoint disclosure. Falsifier: scheduled updates without endpoint content fail where endpoint disclosure succeeds.
\item \textbf{Prior installation by sampling the indeterminate branch.} Reversion-to-prior requires a prior; an avoided branch has none, because every indeterminate episode was escaped before its outcome was recorded. Graded, safe exposure \emph{to indeterminacy itself} --- procedures or tasks whose endpoint is deliberately unsignaled but whose outcomes are benign and completed --- builds $\mathbb E[D\mid o=0]$ by access, not assumption. This names a specific exposure target distinct from cue exposure: the affected person may tolerate the needle and still need exposure to \emph{not knowing when it ends}. Falsifier: indeterminacy exposure with benign completions produces no shift in the $o=0$ response while matched cue exposure does.
\item \textbf{The anchor as a borrowed prior-holder.} A trusted companion who visibly remains calm is a second estimator whose readout of the same situation is still ``within bounds'': coupling to them lends their prior to the branch the patient's own estimator cannot hold open --- an uncertainty-side reading of $k$ consistent with presence helping precisely when things go wrong (the failed-venipuncture rescue), and with fading as gradual withdrawal of the borrowed prior once one's own is installed (P10). Falsifier: companion benefit is unchanged between determinate and indeterminate phases of a procedure.
\end{enumerate}

\subsubsection{The recalibration law: solving for correct reading}
\label{sec:recal}
The three routes above are one discrete bookkeeping system. Model the $o{=}0$ branch estimator as a retained-record reader over exact counts (everything below is finite and rational; the algebra is toy-exact, the clinical bridge is Dr). After each indeterminate episode the estimator retains one \emph{typed} record: $C$ (ran to benign completion), $H$ (ended in harm), or $E$ (escaped before the outcome --- \emph{censored}). With counts $n_C,n_H,n_E$ and prior pseudo-counts $a,b>0$, the calibrated benign-resolution estimate is the smoothed retained frequency
\begin{equation}
\hat p \;=\; \frac{n_C+a}{\,n_C+n_H+a+b\,},
\qquad
\eRRHhat(o{=}0)\;=\;\hat p\cdot\bar D_C,
\label{eq:ledger}
\end{equation}
where $\bar D_C$ is the exact mean duration of the retained completions --- and where a reader with $n_C=0$ must return $\bot$, not $0$, because $\bar D_C$ is then simply undefined.

\textbf{The lesion, derived.} The floor-default of Eq.~\ref{eq:prhcu} is exactly a \emph{typing error in the ledger}: the affected estimator tallies censored records as outcome records,
\begin{equation}
\hat p'\;=\;\frac{n_C+a}{\,n_C+n_H+\mathbf{n_E}+a+b\,}\ \xrightarrow[\;n_E\gg\;]{}\ 0,
\label{eq:lesion}
\end{equation}
so a lifetime of escapes --- each of which proved nothing about the branch's outcome --- is read as a lifetime of evidence that the branch ends badly. Every escape then lowers $\hat p'$ further, which is the self-confirming persistence loop of \S\ref{sec:course} written as one fraction.

\textbf{The correction, solved.} Reading correctly again is three ledger operations, ordered:
\begin{enumerate}[leftmargin=1.7em]
\item[T1] \textbf{Re-type the censored records} ($E$-records out of the outcome tally): this alone lifts $\hat p'\to\hat p$ with \emph{no new experience required}. Clinically it is the post-episode debrief that re-classifies ``we stopped early'' from failure-evidence to no-evidence --- an epistemic act, not an exposure.
\item[T2] \textbf{Append completions} (the access route): with clean typing and benign outcomes ($n_H{=}0$), the margin condition $\hat p_n\,\bar D_C>\tcorr$ is met after
\begin{equation}
n^{*}\;=\;\max\Big(0,\;\Big\lfloor\,\frac{\rho\,(a+b)-a}{1-\rho}\,\Big\rfloor+1\Big),\qquad \rho=\frac{\tcorr}{\bar D_C}<1,
\label{eq:nstar}
\end{equation}
completed indeterminate episodes --- a finite, per-person, computable dose of ``not knowing when it ends, and it ended fine.'' (The floor-plus-one form enforces the strict margin inequality at integer boundaries --- a correction from independent review; $\bar D_C$ is treated as fixed while completions accrue, a stated simplification.)
\item[T3] \textbf{Cover the gap while $n<n^{*}$}: borrowed pseudo-counts from a coupled estimator, $a\to a+k\,a_{anchor}$ (the companion's prior, weighted by the coupling), and a bounded re-read interval $\Delta$ inside the episode, which caps the unresolved load each read cycle can accumulate at $\Delta$ --- so no collapse occurs provided $\Delta<B$, the person's finite buffer bound. The dental duty cycle (work intervals at $\sim$60--70\% of the recalled 3--5 minute tolerance) is this inequality applied; fading is $k\to0$ as $n\to n^{*}$ with the margin held positive throughout (the record-side statement of P10).
\end{enumerate}

\textbf{Order matters, and that is testable.} T1 predicts that pure reattribution of past escapes (debrief without any new exposure) already shifts the $o{=}0$ response; T2 predicts a dose--response in completed indeterminate episodes with a computable threshold; T3 predicts collapse-free practice below $\Delta<B$ and rung-timed fading failures above it. Kill criteria: if reattribution alone never moves the $o{=}0$ response while completions do, the mis-typing account (Eq.~\ref{eq:lesion}) is wrong and only sampling matters; if completions do not move it either, the ledger model dies entirely; if response tracks total time-in-situation rather than \emph{completed-episode count}, the record-typing structure is superfluous and plain extinction suffices. All constants in Eqs.~\ref{eq:ledger}--\ref{eq:nstar} are exact rationals estimable from the logging protocols already specified; nothing here requires a continuum quantity, and the clinical reading of every symbol is \tier{Dr}.

\subsubsection{Clinical translation (hypothesis-generating, not guidance)}
In medical-science terms, the ledger model specifies a candidate treatment mechanism --- \emph{censoring-aware recalibration} --- as three ordered components, each corresponding to a recognizable clinical operation but carrying a non-standard prediction:
\begin{enumerate}[leftmargin=1.7em]
\item \textbf{Post-episode reattribution of interrupted procedures} (T1). A structured debrief that reclassifies prematurely terminated episodes as \emph{non-informative} about outcome --- targeting the evidential status of past events, not the safety of future ones. This is distinct from generic reassurance, which the model classes as assumption-route and predicts to be non-durable; reattribution of the historical record is predicted to shift the response to unpredictable-duration procedures \emph{before any new exposure occurs} --- a prediction no exposure-based account makes.
\item \textbf{Completion-prioritized graded exposure under unpredictable duration} (T2). The proposed therapeutic unit of dose is the \emph{completed unpredictable-duration episode}, not minutes of within-situation time. This is compatible with inhibitory-learning accounts that emphasize expectancy violation over habituation \cite{craske2014}, but sharpens them: the model predicts dose--response in completed-episode count with a per-patient computable target $n^{*}$ (Eq.~\ref{eq:nstar}), and predicts that time-in-situation without completion (session ends by escape) contributes approximately nothing.
\item \textbf{Scaffolded interim support with planned withdrawal} (T3). Companion-assisted sessions, scheduled intra-procedural information updates, and pause protocols that keep any single unresolved interval below the patient's measured tolerance bound --- with companion fading timed to the completion count rather than to symptom ratings, and with the withdrawal criterion checkable from physiological recovery slopes (C18).
\end{enumerate}
None of the above is clinical guidance: it is a mechanism hypothesis with an unusual property --- its three components make \emph{ordered, separable} predictions (reattribution before exposure; completion count over exposure time; slope-timed fading), so a small factorial trial could refute each component independently. Until such a trial exists, established exposure-based treatment delivered by qualified clinicians remains the standard of care, and nothing in this section modifies it.

\section{A minimal generative model: the anchor--drive model}
\label{sec:adm}
A hypothesis paper that asks for parameter recovery must supply a generative model from which parameters can be recovered. The anchor--drive model (ADM) is the smallest one consistent with the constructs above. Let $d_t\ge0$ be the deviation of a regulated variable from its calm anchor (the subtype decides which variable: end-tidal CO$_2$ error, postural sway, heart-rate excursion, jaw or airway tension), $T_t\ge0$ a cue-locked drive, $v_t\ge0$ the system's intrinsic damping, and $g_t\ge0$ its corrective coupling. The deviation obeys
\begin{equation}
d_{n+1}=d_n+\Delta\big[T_n-(v_n+g_n)\,d_n\big],
\label{eq:adm}
\end{equation}
or in the continuous limit $\dot d=T-(v+g)d$. The energy $E=\tfrac12 d^2$ falls exactly when dissipation dominates injection, $(v+g)\,d\ge T$, so the uncorrected deviation relaxes toward the fixed point $d^{*}=T/(v+g)$. The two policy classes are two moves on this equation: an engagement-preserving policy raises the coupling, $g\to g_0+\Delta g$ (paced breathing, applied tension, gaze anchoring, a corrective button); a defensive-termination policy removes the drive, $T\to0$ (leaving). Viability is $\Vset=\{d\le d_{\max}\}$ and the harm-free region is $\Hset=\{d\le d_{harm}\}$ with $d_{harm}>d_{\max}$. \tier{definition}

\noindent\textbf{Proposition (closed-form horizon).} With constant $T,v,g_0$ and $d(t_0)=d_0<d_{\max}$, writing $d^{*}_0=T/(v+g_0)$ and $d^{*}_{c}=T/(v+g_0+\Delta g)$:
\begin{align}
\eRRHc(t_0)&=0 &&\text{if } d^{*}_{c}>d_{\max},\notag\\[-0.1em]
\eRRHc(t_0)&=\tau_{\max} &&\text{if } d^{*}_{0}\le d_{\max},\notag\\[-0.1em]
\eRRHc(t_0)&=\frac{1}{v+g_0}\ln\frac{d^{*}_0-d_0}{d^{*}_0-d_{\max}} &&\text{otherwise;}
\label{eq:closedform}
\end{align}
the first case is ``no corrective policy preserves $\Vset$'' (C1), the second is ``the state never leaves $\Vset$'', so the horizon is censored at $\tau_{\max}$ (C2).

\emph{Sketch.} The uncorrected trajectory is monotone, so A1 holds; correction started at $t_{on}$ succeeds iff $d(t_{on})\le d_{\max}$, because afterwards $d$ relaxes monotonically toward $d^{*}_c\le d_{\max}$. The latest such $t_{on}$ is the crossing time of $d_{\max}$ on the uncorrected path. $\square$

Five properties follow, each an elementary consequence of Eq.~\ref{eq:closedform}:
\begin{enumerate}[leftmargin=1.5em]
\item[P1] \textbf{Finiteness.} A drive that overwhelms the corrected system yields a finite horizon without an ad hoc cutoff; censoring (C2) is needed only when the drive never exceeds the viability boundary.
\item[P2] \textbf{Flow.} Along the trajectory the horizon decreases at unit rate (C3), which is what Module~A must return trial by trial (\S\ref{sec:irhab}).
\item[P3] \textbf{Collapse boundary.} Rescue exists iff $v\ge v^{*}=T/d_{\max}-g_0-\Delta g$. The horizon therefore does not shrink gracefully as intrinsic damping falls; it drops to zero at $v^{*}$ (Fig.~\ref{fig:adm}). A vasovagal trajectory is a fall of $v(t)$ toward zero, so BII is predicted to show a cliff rather than a slope, and raising $\Delta g$ (applied tension \cite{ost1987}) moves the cliff to lower $v$ or removes it. Applied tension thus reads as a parameter move on Eq.~\ref{eq:adm}, not as a separate mechanism.
\item[P4] \textbf{Correction latency.} Under correction the deviation relaxes with time constant $1/(v+g_0+\Delta g)$, which is the model-side quantity behind $\tcorr$. As $v\to0$ the horizon collapses while the latency rises, so the margin $M^E$ crosses zero before the horizon does: transition is predicted to precede the causal deadline by a latency-sized interval.
\item[P5] \textbf{Exit.} Removing the drive returns $d$ to zero at rate $v+g_0$; $\dRRHc$ is the time until the uncorrected path crosses $d_{harm}$. For a slowly driven system with an open door, $\dRRHc\gg\eRRHc$ whenever $\Delta g$ is small (the spider case: no effective engagement action exists, so the engagement horizon is short while exit remains available). When exit is externally delayed (a scanner bore), $\dRRHc$ is short while $\eRRHc$ may be long.
\end{enumerate}

\begin{figure}[!t]
\centering
\includegraphics[width=0.98\columnwidth]{fig1_adm_horizon.pdf}
\caption{Causal engagement horizon $\eRRHc$ as a function of intrinsic damping $v$ in the anchor--drive model ($T=1$, $g_0=0.2$, $d_{\max}=1$, $d_0=0$, $\tau_{\max}=10$; arbitrary units). Executed run \cite{lahtee2026ru}: $\eRRHc=1.72$ at $v=0.5$, $1.28$ at $v=0.2$, and $0$ for $v<0.2$ when $\Delta g=0.6$; with $\Delta g=1.0$ the collapse boundary $v^{*}=T/d_{\max}-g_0-\Delta g$ moves below zero and the cliff disappears. Where both couplings rescue, the curves coincide, because the horizon is fixed by the uncorrected trajectory; $\Delta g$ decides only whether rescue exists and how fast it acts (P3, P4). \tier{finite\_diagnostic, toy tier}}
\label{fig:adm}
\end{figure}

\subsection{What the model is and is not}
ADM is a toy: linear, single-variable, with piecewise-constant parameters. It is not a claim about circuitry. Its role is threefold: it makes the horizon a computable function of state so that A1, C1--C3 and P1--P5 are consequences rather than stipulations; it is the simulator from which synthetic participants with known parameters are generated for parameter recovery, design-level Fisher information and power (\S\ref{sec:stats}); and it is the positive control of every kill gate (\S\ref{sec:gates}). The discrete energy-balance algebra behind the rescue condition is recorded as machine-checked in the author's companion clinical-context project \cite{lahtee2026birca}; that record was not re-verified for this manuscript, and no RRHM claim inherits its status. Any physiological reading of $d$, $v$, $T$ or $g$ is \tier{Dr}. The respiratory instantiation ($d$ = ventilatory error) has a near neighbor in the suffocation-false-alarm hypothesis of panic \cite{klein1993}: both posit a biased alarm about a regulatory margin. RRHM's distinguishing commitments are that the biased quantity is a \emph{time} to loss of engagement-preserving correction and that the bias is \emph{cue-locked}; Gate K10 is where the two hypotheses part.

\section{The Open-Door dissociation and the rival it names}
\label{sec:opendoor}
A spider-phobic participant may stand near an open exit, fully understand that escape is immediately available, and still experience severe defensive urgency. In the vocabulary above,
\[
\dRRHhat\gg0,\qquad \eRRHhat\rightarrow0:
\]
the participant can leave, but predicts that remaining engaged is no longer regulatable. The reverse case is a calm patient undergoing MRI: immediate exit may be constrained, yet engagement remains tolerable, $\dRRHhat\downarrow$, $\eRRHhat\gg0$. The defining dissociation is
\[
\boxed{\text{escape availability}\neq\text{engagement-preserving recoverability}}.
\]

Earlier editions called this a paradox for ``any theory equating recoverability with the ability to terminate the encounter''. No major theory holds that equation, so the target is now named precisely. The rival is the class of \emph{assumption-route accounts}: any account on which declared safety, reassurance, or the mere availability of exit is a recalibrating input to the quantity that governs transition. By consequence 2 of \S\ref{sec:anchor}, declared safety is assumption augmentation and cannot restore a distinction the horizon channel never carried; only access---sampling the state near the boundary---can. The open door is therefore a prediction rather than a paradox: exit availability and verbal safety should shift $\dRRH$-side behavior and verbal ratings while leaving $\eRRHhat$ and transition timing unchanged, whereas boundary sampling should shift both (\S\ref{sec:decisive}). The clinical literature on behavioral experiments versus verbal disputation is compatible with this reading \cite{bennettlevy2004}, and compatibility is all that is claimed. The MRI case is confounded as an observation---a calm patient has both a long estimated horizon and low fear---and is resolved only by manipulation, not by description.

\section{What neighboring theories already explain}
RRHM should not claim novelty by relabeling established science, and kinship is disclosed wherever it exists. Holdings declaration for this section and Table~\ref{tab:attacks}: the rival characterizations here rest on abstract-level familiarity and secondary readings, not on primary full texts read for this manuscript; under the standing full-text-or-NO-DATA rule every containment claim in this section therefore carries the same provenance limitation as the fenced rows of the companion channel map, and none is scored against RRHM anywhere in the corpus.

\subsection{Conditioning and extinction}
Associative learning explains how cues acquire aversive predictions, how avoidance is learned, and how exposure can produce inhibitory learning \cite{krypotos2015,pittig2018,kausche2025,craske2014}. RRHM adds only a candidate variable for \emph{when engagement is terminated}.

\subsection{Threat imminence}
Threat-imminence models predict changing defensive strategies as danger approaches, with human work showing prefrontal-to-periaqueductal gray shifts under increasing imminence and computational accounts placing avoidance and escape at distinct spatiotemporal scales \cite{mobbs2007,mobbs2020}. RRHM is distinct only if regulatory imminence can vary while physical threat imminence remains fixed.

\subsection{Escape-decision and flight-initiation models}
\label{sec:fid}
Flight-initiation-distance (FID) paradigms model \emph{when} an organism switches to escape from an approaching threat, with Bayesian decision models and neural dissociations between fast reactive and slow cognitive escape \cite{qi2018,fung2019}. This literature is the most direct threat to the novelty claim, and the relationship is stated exactly: the quantity optimized in an FID model is the latest time at which termination remains effective, which is $\dRRHc$. In an FID task there is no engagement-preserving action, $\Pi^E=\varnothing$, so by convention C1 $\eRRHc\equiv0$ and RRHM reduces to the FID switching problem. RRHM contains FID as its $\Pi^E=\varnothing$ special case and adds the engagement horizon. Distinctness therefore requires a paradigm with a non-empty $\Pi^E$ whose deadline is manipulated independently of threat kinematics. The finding that trait anxiety hastens escape only under slow threats \cite{fung2019} locates individual differences in the regime in which a horizon channel has time to be read; RRHM's additional commitment is that in phobia the earliness is cue-locked and tracks the engagement deadline rather than trait anxiety alone. Kill: if adding a genuine $\Pi^E$ to an FID task produces no deadline-tracking of transition timing, the engagement horizon is not a separate variable.

\subsection{Perceived control and learned controllability}
Control strongly modulates defensive behavior and extinction, and the distinction between actual and perceived control is the axis of that literature \cite{maier2016,hartley2014,wood2015}. RRHM does not claim that distinction. Its third dissociation---perceived control versus causally effective engagement preservation---is disclosed as kinship, not novelty. What RRHM claims beyond it is that the \emph{temporal} channel has a bias separately identifiable from the control channel's bias: two conditions can be matched on action number, action efficacy and perceived control and still differ in the causal deadline, and that difference alone must move transition timing.

\subsection{Interval timing under arousal}
Emotion and arousal distort perceived duration in systematic ways, and retrospective ``time slowed down'' reports during fear are partly a memory artifact \cite{droitvolet2007,stetson2007}. The distortion is documented in the target population itself: spider-phobic participants gave longer, more variable estimates of a short interval spent observing a spider than non-phobic controls \cite{watts1984}. $M_{Timing}$ is therefore an empirically grounded rival in specific phobia, not a hypothetical one, and this is part of why the within-trial probe slopes of \S\ref{sec:k11} are mandatory rather than optional. If $\eRRHhat$ contraction is a special case of generic clock distortion under arousal, RRHM adds nothing. A generic timing model $M_{Timing}$ therefore joins the rival set, and Gate K11 requires the contraction to survive adjustment for temporal bisection and interval reproduction measured under matched arousal.

\subsection{Allostatic self-efficacy}
Allostatic self-efficacy concerns estimated regulatory capacity \cite{stephan2016}. RRHM is narrower: it proposes a cue-specific temporal estimate of how long effective engagement-preserving regulation remains available. Failure to discriminate the constructs statistically is a failure of RRHM, not a semantic dispute.

\subsection{Active inference}
Active inference is broad enough to represent interoception, uncertainty, policy selection, expected future states, and an eRRH-like quantity \cite{paulus2019}. RRHM therefore claims no ontological superiority. Its value must come from parsimony, clinical interpretability, and predictive performance on a defined transition problem. A flexible generative model with unconstrained parameters can absorb any attribution after the fact; the comparison in \S\ref{sec:stats} therefore requires preregistered priors on both sides and a negative control.

\subsection{Two-system accounts}
The claim that $\eRRHhat$ is latent and need not be verbal aligns with two-system frameworks separating defensive circuitry from conscious fear \cite{ledoux2016}; RRHM inherits their caution that verbal report is one readout among several, not the state.

\section{Adversarial attacks the theory must survive}
Table~\ref{tab:attacks} precommits the attacks, the RRHM-specific burden, and the kill criterion for each. Three rows are new to this edition (escape-decision models, generic timing, assumption-route recalibration), and the perceived-control row has been corrected: earlier editions proposed to match perceived control while varying action efficacy, which is impossible in principle, because perceived control \emph{is} the belief that one's action affects the outcome. The corrected manipulation keeps both actions fully effective and varies only the deadline.

\begin{table*}[!t]
\centering
\caption{Precommitted adversarial attacks, discriminating predictions, and kill criteria (v5).}
\label{tab:attacks}
\small
\begin{tabularx}{\textwidth}{p{0.15\textwidth}X X p{0.20\textwidth}}
\toprule
\textbf{Attack} & \textbf{Why a rival account may suffice} & \textbf{RRHM-specific burden} & \textbf{Kill criterion} \\
\midrule
Circularity & ``Recoverability'' may merely redescribe persistence or escape. & Estimate $\eRRHc$ and $\eRRHhat$ before $\Tdt$ using independent channels (Modules A--C$'$, F). & If either horizon requires the eventual escape outcome, reject construct (K1). \\
Threat imminence & Earlier defense may reflect closer or more intense threat. & Hold physical imminence and threat constant while manipulating the causal engagement deadline. & No incremental effect beyond threat-imminence variables (K4, K5). \\
Escape-decision models & The FID optimal switching time may already be the variable. & Show deadline-tracking in a task with a non-empty $\Pi^E$ and fixed threat kinematics (\S\ref{sec:fid}). & No effect of the engagement deadline beyond FID variables. \\
Perceived control & A stop button may explain differences through control alone. & Both actions 100\% effective; only the deadline $D$ differs; match action count, control ratings, and exit. & Control ratings absorb the $D$ effect (K3, K4). \\
Generic timing & Arousal-driven clock distortion may produce contraction. & Contraction survives adjustment for temporal bisection and interval reproduction under matched arousal. & Contraction disappears after adjustment (K11). \\
Assumption-route recalibration & Reassurance or exit information recalibrates as well as sampling does. & Told-$D$ versus learned-$D$ dissociation: behavior tracks $D$ before accurate report only in the learned arm; told shifts do not transfer. & Told and learned arms produce equal, durable, transferable shifts. \\
Self-efficacy & eRRH may be a temporal phrasing of regulatory confidence. & Demonstrate discriminant validity and parameter separability from ASE. & Parameters non-identifiable, or the disattenuated $|r|$ demonstrated $\ge0.85$ (K2, K3). \\
Active inference & A sufficiently general generative model can retrofit eRRH. & Win on preregistered held-out prediction with fewer clinically interpretable parameters under preregistered priors. & No predictive gain or parsimony advantage (K5). \\
BII counterexample & Vasovagal physiology violates a universal sympathetic model. & Allow causal eRRH to shorten genuinely ($v\to0$ in ADM) and downstream effectors to differ. & Model requires universal sympathetic arousal. \\
Acrophobia counterexample & Postural and vestibular instability may drive behavior directly. & Treat sensorimotor state as part of $x_t$ and test whether transition timing follows recoverability margin. & eRRH adds nothing beyond postural state. \\
High arousal without phobia & Skydiving and roller coasters show arousal without pathological avoidance. & Predict voluntary engagement when eRRH remains adequate despite arousal. & Model equates arousal with phobia. \\
Unconscious/VR exposure & Therapeutic updating occurs without explicit appraisal or direct hazard contact. & $\eRRHhat$ must be latent and representation-mediated; updating must be access-based. & Model requires conscious report or literal hazard exposure. \\
Natural remission & Recovery occurs outside formal CBT. & Allow informative naturalistic sampling and changes in actual capacity or context. & Formal exposure is required by model. \\
Transdiagnostic spillover & Recoverability may also matter in panic, PTSD, pain, dyspnea, FND. & Establish cue-locked specificity and incremental prediction for specific phobia. & Same construct predicts all disorders equally without cue-locked structure (K10). \\
\bottomrule
\end{tabularx}
\end{table*}

\section{Cross-phobia stress test}
RRHM predicts a shared transition variable but not shared physiology. Table~\ref{tab:cross} states, for each class, which regulated variable plays the role of $d$ in ADM, the candidate contraction, and---new to this edition---the strongest account that needs no horizon at all. RRHM must beat that account on \emph{transition timing}, not on physiology; every row is \tier{Dr}.

\begin{table*}[!t]
\centering
\caption{Cross-phobia mapping under RRHM, with the strongest non-RRHM account per class. All rows are declared bridges (Dr) with the same falsifier: transition timing is fully predicted by the rival account after the shared eRRH parameter is added.}
\label{tab:cross}
\small
\begin{tabularx}{\textwidth}{p{0.11\textwidth}p{0.15\textwidth}X p{0.20\textwidth}p{0.17\textwidth}}
\toprule
\textbf{Class} & \textbf{Regulated deviation $d$} & \textbf{Candidate eRRH contraction} & \textbf{Strongest non-RRHM account} & \textbf{RRHM must win on} \\
\midrule
Animal & distance / sensory-tracking error; injury-avoidance margin & stimulus movement makes continued engagement appear to lose safe corrective options rapidly; $\Delta g\approx0$ & disgust- and imminence-driven responding with no horizon term & timing tracks a programmed engagement deadline at fixed imminence \\
Heights / natural environment & postural sway; visual--vestibular mismatch & estimated balance-correction window narrows despite objective protection & visual--vestibular mismatch drives behavior directly & timing follows recoverability margin beyond sway itself \\
BII & circulatory deviation (BP, HR); $v(t)\to0$ & actual or estimated circulatory recoverability contracts around puncture, blood, injury cues; collapse boundary P3 & vasovagal reflex plus disgust, no estimate involved & estimated horizon predicts transition beyond measured hemodynamics \\
Situational / MRI / flying & respiratory or spatial constraint; ventilatory error & engagement predicted to become non-recoverable while objective danger remains stable & claustrophobic threat appraisal; suffocation false alarm \cite{klein1993} & cue-locked temporal bias, not generic alarm (K10) \\
Procedural / dental & jaw, swallowing, respiratory comfort; communication channel & task-compatible regulatory actions expected to become ineffective during continued procedure & pain expectancy and loss of communication & timing tracks the corrective-window structure (A1 windows) \\
Visceral / choking / vomiting & airway or swallowing error; visceral motor state & predicted point of no effective correction approaches too early & airway-protective reflex; conditioned nausea & timing tracks estimated window beyond reflex latency \\
\bottomrule
\end{tabularx}
\end{table*}

The cross-class claim survives only if a shared eRRH parameter improves prediction of transition timing without requiring unrelated parameter definitions for every subtype. In ADM the permitted heterogeneity is explicit: subtypes may differ in which variable is $d$, in $v$, $T$, $g_0$ and $\Delta g$, but not in the equation or in the meaning of the horizon. If every subtype needs a different construct hidden under one acronym, the unification claim fails (K9).

\section{Treatment observations as stress tests, not proof}\label{sec:txobs}
RRHM is compatible with several treatment observations but must not use compatibility as validation. Single-session exposure can be as effective as multi-session formats in pooled analyses, implying that information gain may matter more than session count \cite{odgers2022}. Virtual-reality exposure produces clinically meaningful improvement across anxiety disorders \cite{carl2019}, indicating that updating can occur through surrogate representations. A systematic review of 39 controlled experiments across ten unconscious-exposure paradigms reports reduced behavioral avoidance and self-reported fear in phobic participants without inducing subjective fear during the intervention, supporting a latent updating process \cite{siegel2025}. Safety behaviors show context-dependent effects, cautioning against a binary rule \cite{meckes2025}. Inhibitory-learning accounts emphasize expectancy violation over habituation \cite{craske2014}, and applied tension in BII is an existing engagement-preserving action with a deadline \cite{ost1987}.

RRHM classifies behavior during exposure by causal function:
\begin{enumerate}
\item \textbf{engagement-expanding}: increases actual task-compatible regulation (raises $\Delta g$ or $v$);
\item \textbf{learning-blocking}: permits contact but prevents independent sampling of the feared boundary (records are not returned; the estimate cannot descend);
\item \textbf{termination}: ends engagement and restores global regulation through escape ($T\to0$);
\item \textbf{assumption-only}: verbal reassurance or safety information without boundary sampling---predicted to sharpen ratings but not to recalibrate $\eRRHhat$ durably.
\end{enumerate}
These categories are hypotheses about function, not treatment prescriptions. Whether a given safety behavior is (1) or (2) is an empirical question about whether it returns boundary records, which may explain the inconsistent literature on safety behaviors.

\section{Independent measurement: IRHAB}
\label{sec:irhab}
Direct validation requires the Independent Recoverability Horizon Assessment Battery (IRHAB), with strict separation of predictor and outcome.

\subsection{Module A: causal ground truth, trial by trial}
A task externally programs a deadline $D_j$ until which an engagement-preserving action remains fully effective; after $D_j$ the action has no effect. For trial onset $t_{0,j}$,
\[
\mathrm{eRRH}_{C,j}(t)=D_j-t,\qquad t\in[t_{0,j},D_j],
\]
and zero afterwards (C1). The deadline is fixed by design and does not depend on participant behavior. Earlier editions defined a per-trial constant $D_j-t_{0,j}$; that quantity is the horizon at onset only, whereas the hazard model of Eq.~\ref{eq:hazard} requires the within-trial flowing value.

\subsection{Module B: prospective explicit estimate, probed repeatedly within a trial}
Before the outcome occurs, participants estimate how much time remains before engagement-preserving correction will cease to work. Fear, threat expectancy, perceived control, and self-efficacy are rated separately.

The estimate must be probed at \emph{two or more times within the same trial}, not once at onset. This is not a refinement: it is what makes the construct identifiable at all. A horizon bias makes the reported remaining interval start lower and fall at unit rate; a clock-rate distortion makes it start correctly and fall faster. Probing once cannot see a rate, so the two accounts collapse into one; probing twice separates them (\S\ref{sec:sim}). Probes are brief (a single dial or key report) and are scheduled at fixed times so that probe timing is not itself informative about $D$. Probe times are specified at \textbf{0.3 and 1.0~s after trial onset}: early probing keeps almost every trial in the sample and so minimises the survivorship attenuation of the probe slope, while retaining enough separation between probes to estimate that slope against report noise (\S\ref{sec:k11}).

\subsection{Module C: choice-derived estimate}
An adaptive staircase estimates the indifference point between correcting now and delaying corrective action. This point is jointly determined by the horizon estimate, the cost of correcting early, risk attitude, and time perception; extracting $\eRRHhat$ from it requires a generative choice model whose utility assumptions are preregistered. Module~C is therefore model-dependent and is retained for its nonverbal character, not as a model-free measure.

\subsection{Module C$'$: interval-production estimate}
Participants press a key at the moment they judge that correction would stop working (interval production). This is a psychophysical estimate rather than a value-based choice: it needs no utility model, provides the nonverbal estimate that Module~C cannot provide model-free, and yields the temporal-bisection and interval-reproduction baselines required for $M_{Timing}$ (Gate K11) when the same task is run on neutral intervals under matched arousal.

\subsection{Module D: multimodal physiology}
Candidate channels include ECG/heart rate, beat-to-beat blood pressure, respiration, end-tidal CO$_2$, electrodermal activity, pupil diameter, EMG, postural sway, and movement kinematics. No single channel is defined as an eRRH biomarker. Recovery slopes of these channels after a perturbation are the measurable proxies for $v$ and $g$ in ADM, and are the only route by which a causal horizon can be bounded outside the laboratory (\S\ref{sec:course}).

\subsection{Module E: independent outcome}
Only after horizon estimation is completed should the experiment score escape latency, termination, persistence, freezing, behavioral approach, or action switching as outcomes.

\subsection{Module F: correction latency (mandatory)}
The margin $M^E=\eRRHhat-\tcorr$ contains a second latent quantity. Module~F returns a prospective estimate of $\tcorr$ (``how long would it take you to make the correction work?'') and a measured proxy: initiation-to-effect latency of the same corrective action on neutral, non-threat trials.

Module~F is mandatory, and the reason is exact rather than prudential. In the hazard model of Eq.~\ref{eq:hazard}, a shift of $\tcorr$ by $\delta$ changes the linear predictor by $\beta_1\delta$, which an offsetting shift of $\beta_{0}$ cancels completely. Correction latency and the hazard intercept are therefore confounded \emph{exactly}, in a direction with ratio $-1/\beta_1$, and no amount of transition-timing data of any kind resolves it; simulated recovery of $\tcorr$ from timing alone returns $r\approx0$ (\S\ref{sec:sim}). Only an independent measurement breaks it.

\subsection{Identifiability gate and anti-circularity}
Before data collection, the design-level Fisher information matrix is computed under ADM for the parameter vector $(K,c,\tcorr,\beta_0,\beta_1)$ --- horizon gain, clock gain, correction latency, and the two hazard parameters. A design whose information matrix is rank-deficient, or ill-conditioned beyond a preregistered bound, in any theory-critical direction is rejected before recruitment; the question it would have asked is structurally unanswerable from the records it returns. \tier{definition}

This gate has been run on the design of \S\ref{sec:decisive} and it rejected the first version of it. A deadline-manipulation design that probes the estimate only at trial onset has numerical rank 3 of 5 --- with two, three, or four deadline levels alike --- so adding deadline levels does not help. Mid-trial probes raise the rank to 4; Module~F raises it to 4 by a different route; only both together reach full rank (Fig.~\ref{fig:rank}). The battery specified above is the smallest one that passes its own gate.

The anti-circularity rule is:
\begin{quote}
\textbf{$\eRRHc$, $\eRRHhat$ and $\tcorr$ must never be computed from the defensive-transition outcome they are intended to predict.}
\end{quote}
This guarantee holds in paradigms in which $D$ is programmed. Outside them it does not transfer automatically; \S\ref{sec:course} states what replaces it.

\subsection{Implicit deadline learning}
Convention C3 requires $\eRRHhat$ to track $D$ within a trial, which is impossible on a naive participant's first trial. Either the deadline is told or it is learned, and the two routes make different predictions (\S\ref{sec:opendoor}). IRHAB therefore includes a learned arm in which $D$ is never stated, deadlines are stable within blocks, and Module~B/C$'$ reports are collected block by block, so that the emergence of deadline-tracking in behavior can be dated relative to the emergence of accurate report.

\section{Decisive experiment}
\label{sec:decisive}
The primary proof-of-principle study is a safe $2\times2\times2$ design:
\begin{itemize}
\item \textbf{eRRH} (within-subject): causal deadline $D=12$~s versus $D=3$~s. Both conditions offer the same corrective action, fully effective if started before $D$ and ineffective afterwards; efficacy, action count, and instructions about efficacy are identical. Only the temporal geometry of recoverability differs.
\item \textbf{dRRH} (within-subject): exit immediately available versus exit delayed by a fixed, disclosed interval. A genuine emergency stop is always retained (\S\ref{sec:safety}).
\item \textbf{Instruction} (between-subject): deadlines told versus learned implicitly.
\end{itemize}
Objective threat, threat probability, physical imminence, uncertainty, sensory intensity, action number, and perceived control are matched by construction rather than by post hoc adjustment. Four within-subject cells follow (Table~\ref{tab:cells}).

\begin{table}[!t]
\centering
\caption{Predicted cells of the decisive experiment.}
\label{tab:cells}
\small
\begin{tabular}{llp{0.47\columnwidth}}
\toprule
Cell & $\eRRH$ / $\dRRH$ & RRHM prediction \\
\midrule
A & long / high & flexible engagement; latest, least frequent transition \\
B & short / high & early transition despite preserved control; escape chosen \\
C & short / low & early transition; freezing or distressed endurance chosen \\
D & long / low & constrained but tolerable engagement; late transition \\
\bottomrule
\end{tabular}
\end{table}

The precommitted predictions are
\[
\boxed{
\begin{aligned}
&\text{(1) same threat, imminence, uncertainty, control,}\\[-0.1em]
&\qquad\text{exit, action count and action efficacy}\\[-0.1em]
&\qquad+\ \textbf{different }D\ \Rightarrow\ \textbf{different transition timing};\\[-0.1em]
&\text{(2) in the learned arm, behavior tracks }D\text{ before}\\[-0.1em]
&\qquad\text{report of }D\text{ is accurate};\\[-0.1em]
&\text{(3) in the told arm, report shifts fully, timing shifts}\\[-0.1em]
&\qquad\text{less and does not transfer across a}\\[-0.1em]
&\qquad\text{context-change block.}
\end{aligned}}
\]
If (1) does not occur reliably, the distinctive mechanism of RRHM is substantially weakened. If (2) and (3) fail while (1) holds, the horizon is real but not latent, and the unconscious-exposure row of Table~\ref{tab:attacks} is lost.

\subsection{Planning sample size}
A simulation-based planning estimate, executed for this manuscript, treats the eRRH main effect on log transition latency as a within-subject contrast with per-participant slope heterogeneity (slope SD 0.15 in residual-SD units). For a smallest effect size of interest of 0.20 residual SD with 16 usable trials per deadline level, $N=48$ participants gives $\approx0.94$ power and $N=40$ gives $\approx0.89$ at two-sided $\alpha=0.05$; with 24 usable trials, $N=32$ gives $\approx0.92$. \tier{finite\_diagnostic, planning tier} Inflating by 25\% for censored trials and exclusions gives a Phase-1 target of $N=60$ per instruction arm ($N=120$ total) in healthy or non-clinical fearful adults. The known-groups comparison (diagnosed specific phobia versus controls on cue-locked $\RCE$, standardized difference 0.5) requires approximately 64 per group at 80\% power or 86 per group at 90\% by the standard two-sample calculation \cite{cohen1988}. The committed $N$ is to be set by full hierarchical simulation from ADM including censoring, and will be preregistered.

\section{Statistical and computational adjudication}
\label{sec:stats}
The primary dependent quantity is transition hazard, $h_D(t)=P(\text{defensive transition at }t\mid\text{no prior transition})$, modeled by Eq.~\ref{eq:hazard}. Rival models are implemented prospectively:
\[
M_{CE},\ M_{TI},\ M_{FID},\ M_{Control},\ M_{ASE},\ M_{Timing},\ M_{AI},\ M_{RRHM}.
\]
\textbf{Primary criterion.} One criterion is primary: held-out expected log predictive density (ELPD) under subject-level cross-validation \cite{vehtari2017}, with differences reported with their standard errors; nested cross-validation, calibration, and information criteria are secondary. A single primary criterion is declared so that the adjudication itself is not a garden of forking paths. RRHM must not be declared superior because it can retrospectively fit the same data.

\textbf{Attribution with a negative control.} Beyond the ELPD difference, the fraction of transition-timing residual attributable to the eRRH posit is reported as
\[
\Phi(\Tdt;\eRRH)=1-\frac{V^{*}(C\setminus\eRRH)}{V^{*}(C)},
\]
where $V^{*}$ is the minimized weighted residual energy of the full model $C$ and of the model with the eRRH term released; this is algebraically a relative of the nested likelihood-ratio statistic and is reported for its attribution semantics, not as new statistics. Two rules bind it: it is always reported against a negative control (the same model with $\eRRHhat$ permuted across trials, which must return $\Phi\approx0$), and every model in the comparison carries preregistered priors, because an unconstrained free parameter returns a spurious $\Phi=1$. The second rule is where the active-inference attack is decided.

\textbf{Discriminant validity.} Equivalence testing \cite{lakens2017} replaces ``statistically indistinguishable'': $\eRRHhat$ is declared collapsed into perceived control, ASE, or timing distortion if the correlation is shown equivalent to or beyond $|r|=0.85$ within the preregistered bounds.

\textbf{Parameter recovery and design information.} Before clinical use, parameter recovery must be demonstrated in synthetic participants generated by ADM with known parameters: refit without access to the generating values, and require recovery correlation $r\ge0.70$ for $\eRRHhat$ and $\tcorr$ with bias below a preregistered bound. The design-level Fisher information matrix is computed on the same synthetic data. If eRRH cannot be recovered independently from threat, control, generic delay discounting, timing distortion, or risk preference, the construct is non-identifiable and the design or the construct must be revised.

\section{Disease-course predictions}
\label{sec:course}
\subsection{Acquisition}
The onset hypothesis is
\begin{align*}
\text{Cue locking}+\text{PRHC}&+\text{Premature termination}\\[-0.1em]
+\ \text{Avoidance}&+\text{Restricted sampling}\\[-0.1em]
&\Rightarrow\ P(\text{Phobia})\uparrow .
\end{align*}
Prospective prediction must exceed established predictors such as baseline fear, trait anxiety, behavioral inhibition, prior adverse experience, family history, and avoidance \cite{trumpf2010incidence}. RRHM supplies a maintenance mechanism, not a diathesis: why one person's estimate becomes biased in the first place is not answered here (\S\ref{sec:limits}).

\subsection{Why the deficit is cue-locked}
\label{sec:cuelock}
The cue-locking of PRHC is derived, not assumed. By Eq.~\ref{eq:recordlaw}, the estimate can only descend toward the causal horizon when boundary records are returned; by consequence 2 of \S\ref{sec:anchor}, only access, not assumption, returns them. Avoidance is precisely the closure of the record channel \emph{for one cue}: the person continues to sample recoverability boundaries in every other context, so those estimates stay calibrated, while for the avoided cue no boundary record ever arrives and the local bias persists indefinitely. The specificity of the deficit is therefore inherited from the specificity of the avoidance, and no additional mechanism is required. \tier{Dr}; falsifier: cue-locked bias that persists in people with documented, repeated boundary sampling of that cue, or generic bias in people who avoid only one cue (the latter is Gate K10).

\subsection{Persistence}
Escape restores immediate regulation but prevents observation that engagement-preserving correction may have remained possible. Persistent early termination therefore predicts maintenance of low $\eRRHhat$, and each escape is self-confirming because the record that would refute it is the one never obtained.

\subsection{Remission}
Remission is hypothesized to involve recalibration of $\eRRHhat$ toward causal eRRH through sufficiently informative boundary sampling. Formal exposure is one route; naturalistic benign encounters, virtual exposure, social modeling, sensorimotor retraining, or physiological change may also update the estimate. Because $\varepsilon_{tot}>0$ always (Eq.~\ref{eq:recordlaw}), the model predicts that remitters retain measurable residual miscalibration rather than reaching zero calibration error---a specific, testable difference from accounts on which recovery is the elimination of a faulty belief.

\subsection{Durability and relapse}
Durable remission requires generalization across contexts, body states, cue intensities, and time. Relapse is predicted when illness, pain, sleep disruption, stress, context change, new adverse experience, or reduced physical capacity again contracts causal or estimated eRRH---in ADM, a fall in $v$ or a rise in $T$ restoring the collapse condition of P3.

\subsection{The horizon outside the laboratory}
Outside programmed-deadline paradigms nobody sets $D$, and the viability region $\Vset$ risks collapsing back into behavior (``task-compatible'' defined as ``still doing the task''). This edition states the scope explicitly rather than letting it leak. The non-circularity guarantee holds only where $D$ is programmed. For clinical claims, the causal horizon is replaced by a declared, measurable bridge: the time until the rate of disruption of a regulated variable would exceed its measured repair rate under the best available corrective action, with the repair rate estimated from Module~D recovery slopes (for example post-perturbation heart-rate-variability recovery, sway recovery, or end-tidal CO$_2$ recovery) and $\Vset$ fixed by preregistered physiological bounds rather than by task persistence. This is a declared bridge \tier{Dr}, and every clinical prediction in this section inherits its status. It also states the weld to the author's broader account of health as coordinated repair capacity \cite{lahtee2026srct}: the engagement horizon is that account's repair capacity projected onto the time axis for a single cue. The bridge is falsified if physiologically estimated horizons show no relationship to programmed horizons in participants measured under both.

\section{Extension (Dr): the coupling term --- companions, bereavement, and the dRRH disorder}
\label{sec:coupling}
\textbf{Scope fence.} Everything in this section is a Dr-tier extension \emph{outside} the specific-phobia scope that the rest of this manuscript is locked to. It inherits no status from the run battery of \S\ref{sec:sim}--\S\ref{sec:simremaining}. Holdings are declared up front: the phenomenology used here is DSM-5-TR criterion text and abstracts (Rachman's safety-signal papers; the cognitive-model literature); \emph{no primary full text of any rival theory has been read}, so no rival comparison is asserted anywhere in this section --- every claim below is a commitment RRHM's own equations make, checkable against documents held in full, with its own kill criterion. The withdrawn-comparison concession and the full-text-or-NO-DATA rule are recorded in the companion preregistration (S7, v2).

\subsection{The hard case}
Agoraphobia's diagnostic text itself contains the two facts this section must earn: fear that ``escape might be difficult or help might not be available'' (the $\dRRH$ side, by definition --- restating it in our notation is a relabel and claims nothing), and the \emph{trusted-companion} phenomenon, strong enough to sit inside the diagnostic criteria: a person who cannot enter a situation can enter it when one specific person is present, the same situation, nothing else changed --- and that person is not substitutable by a stranger. A further clinical observation is onset after bereavement without any conditioning episode or preceding panic --- clinically documented at the holding level available here; whether existing theories explain it is NO-DATA under this corpus's holdings rules (the relevant primary texts are held at title level only).

\subsection{One added term}
Let a companion be an external calm anchor coupled to the regulated variable with strength $k\ge0$:
\begin{equation}
\dot d \;=\; T-(v+g_0)\,d-k\,(d-d_a),\qquad d_a\approx0,
\label{eq:coupling}
\end{equation}
so that with a calm anchor the deviation obeys $\dot d=T-(v+g_0+k)\,d$: \emph{coupling to a calm anchor is additive damping}. A companion who can also act contributes separately to the corrective coupling, $\Delta g\to\Delta g+\Delta g_{ext}$. Four consequences are elementary properties of Eq.~\ref{eq:closedform} with $v\to v+k$ (all verified numerically in S9; the same additive-damping algebra is recorded machine-checked in the author's companion clinical-context project \cite{lahtee2026birca} --- that record was not re-verified here and no status transfers; the derivation below stands on its own):
\begin{enumerate}[leftmargin=1.6em]
\item[P6] \textbf{Horizon extension.} $\eRRHc$ is increasing in $k$ (run: $1.39\to2.01$ at $v=0.3$, $k:0\to0.3$, standard toy parameters).
\item[P7] \textbf{Rescue where self-regulation has failed.} The collapse boundary moves from $v^{*}$ to $v^{*}-k$; in particular at $v=0$ (intrinsic damping gone --- the exhausted, the grieving, the vasovagal) a coupled anchor with $k\ge v^{*}$ still yields a finite positive horizon (run: $0\to1.39$). The anchor helps exactly where self-regulation no longer can.
\item[P8] \textbf{Bereavement onset, derived.} $k$ is a property of the \emph{person}, shared across every situation they accompany. Removing $k$ (death, separation) therefore contracts the horizon in every situation where $T/(v+g_0)>d_{\max}$ \emph{simultaneously} --- with no conditioning event, no panic, no new fear learning (run: $1.72\to0$ at $v=0.1$, $k:0.4\to0$). A non-cue-locked, many-situation contraction is exactly the agoraphobic signature, and it falls out of deleting one shared term.
\item[P9] \textbf{Why fading works and abrupt removal does not.} Gradual companion fading is a path $(k\downarrow,\ \Delta g_{self}\uparrow)$ that can be chosen never to cross the collapse boundary, while boundary records keep returning (engagement-expanding, category 1 of \S\ref{sec:txobs}); abrupt removal crosses it, forces early termination, and closes the record channel (the escape, not the companion, is what blocks learning).
\end{enumerate}

\subsection{The double dissociation this forces}
The extension does not get to share K10's cue-locking. It is \emph{required} to predict the opposite pole: specific phobia is a cue-locked additive bias on the $\eRRH$ channel ($\delta_{cue}$); the coupling-loss disorder is a \emph{generic}, non-cue-locked horizon contraction from a deleted shared term. Same instruments, opposite signatures. S9 simulates both generators (per-situation pre/post latency shifts, four situations, $N=40$, 32 trials, pathology-sized effects $k=0.5$, $\delta=3$): a two-line model comparison (uniform shift versus target-only shift) classifies the generating mechanism with accuracy $1.00/0.99$ across 200 cohorts. \tier{finite\_diagnostic, planning tier} At the small effect sizes of the first (untuned) run the same statistic fails ($\approx0.52$--$0.68$), so the dissociation carries an effect-size floor, stated rather than hidden. If measured agoraphobia turns out cue-locked, this dissociation --- and the extension's reading of Criterion B --- dies.

\subsection{Gate K12: capacity versus matched felt safety}
The companion claim must separate from the safety-signal reading (companion as a signal that changes the \emph{estimate}) without asserting anything about the safety-signal literature itself. ADM has no term for felt safety --- $g$, $k$ are dissipation rates, not beliefs --- so the extension is forced to predict: with felt-safety ratings matched and the observer channel held constant, a companion who \emph{can} act extends engagement more than one who cannot. Gate K12 fires against the extension when the capacity contrast, covaried for felt-safety differences, is null. Run on both generators (S9): fires $0.11$ on the coupling generator and $0.96$ on a felt-safety-only generator at $N=60$, 24 trials, $k=0.5$ --- two-sided; at $k=0.35$ power drops (fires $0.41$ on its own positive control), so the gate is preregistered with an effect floor of $k\ge0.5$ or a correspondingly larger $N$. \tier{finite\_diagnostic, planning tier} The companion straddles a second channel --- a trusted person also changes who is watching --- so the live design is three-factor (capacity $\times$ felt safety $\times$ observer), preregistered as case C16 in S7 v2.

\subsection{The phenomenology fence, and what it forced}
\label{sec:fence}
After the first version of this section was written, the accessible empirical phenomenology was collected deliberately as a fence around the model (holdings declared per item; full-text-or-NO-DATA rule in force; the fence document with items F1--F10 is deposited as S10). Three items changed the section.

\textbf{Presence without capacity is causal at the physiological level.} In the one experimental safe-person study located (CO$_2$ challenge in panic patients with their safe person present or absent \cite{carter1995}, abstract holding), presence lowered not only distress and catastrophic cognitions but \emph{physiological arousal} --- although the safe person could do nothing about the CO$_2$. A capacity-only reading of the companion is therefore too narrow, and the companion term is decomposed:
\begin{equation}
\begin{aligned}
\dot d \;=\;& T + \underbrace{\Delta T_{soc}}_{\substack{\text{observer term}\\ \text{(declared, no form)}}} - \big(v+g_0+\underbrace{\Delta g_{ext}\,\mathbb 1[\text{act}]}_{\text{capacity}}\big)\,d\\[-0.2em]
&-\underbrace{k\,(d-d_a)}_{\substack{\text{co-regulation}\\ \text{(presence, channel-borne)}}},
\end{aligned}
\label{eq:companion}
\end{equation}
three separable contributions with three different readouts: $k$ should appear in \emph{post-perturbation physiological recovery slopes} (Module~D's own proxy for effective damping) independently of what the companion does; $\Delta g_{ext}$ in tolerated engagement under matched felt safety (K12); $\Delta T_{soc}$ has no form yet (below). The recovery-slope prediction is preregistered as case C18 with its locked falsifier: if presence changes reports but never recovery slopes, the co-regulation form of $k$ is dead and only a signal reading survives.

\textbf{The coupling is relationship-graded.} Spousal hand-holding attenuates neural threat responding more than a stranger's hand, graded by marital quality \cite{coan2006} (abstract holding). This is consistent with $k$ as a channel whose gain is built from retained mutual history --- and it does not discriminate $k$ from felt safety, which predicts the same gradient; the discrimination remains with C16/C18, not with this datum.

\textbf{The social hole is central, not peripheral.} In the one phenomenological study read in full text (interpretative phenomenological analysis of adolescents' lived experience of panic disorder \cite{baker2022}), one of the two superordinate themes of the entire lived experience is ``an unacceptable self'' --- shame, being seen failing, the other's judgment. The second-reader gap declared below is therefore not an edge case the model may someday address; it is half of the phenomenon, confirmed at the strongest holding level this manuscript has, and building that term is prerequisite work for any further extension.

\subsection{The portable anchor problem}
\label{sec:panchor}
Once the companion is a term of Eq.~\ref{eq:companion} and not a signal, the question changes from ``why does the companion help?'' to: \emph{what is the minimal record channel that carries the coupling?} If $k$ is a channel property, write $k=k(B,H)$, nondecreasing in the real-time interaction bandwidth $B$ of the channel and in the retained mutual history $H$ of the dyad. Holding $H$ fixed (same person), the benefit must then order by bandwidth:
\begin{align*}
&\text{in person} \;>\; \text{live video} \;>\; \text{live voice}\\[-0.15em]
&\quad>\; \text{text} \;>\; \text{token/talisman} \;>\; \text{none},
\end{align*}
with felt safety covaried --- and must \emph{not} order by symbolic or emotional meaning at fixed bandwidth. A talisman carries maximal meaning and zero bandwidth; a live video call carries bandwidth and modest meaning: the coupling and safety-signal readings of the companion are forced apart by exactly this ladder, and the clinical folklore that ``the telephone helps partially'' is an untested bandwidth data point sitting in every clinic. The ladder also yields a quantitative fading rule extending P9:
\begin{enumerate}[leftmargin=1.6em]
\item[P10] \textbf{Fading as a bandwidth staircase.} Descend one rung ($k_{j+1}<k_j$) only when the anchor-free recovery slope --- the measurable proxy for $v+g_{self}$ --- has risen enough that the rescue condition $v+g_0+\Delta g_{self}+k_{j+1}\ge T/d_{\max}$ (where $\Delta g_{self}$ is the person's own corrective coupling, replacing the companion-supplied $\Delta g$) still holds at the next rung. Fading failures are then predicted to cluster at rung descents taken before the slope criterion is met, and the criterion is checkable per patient from wearable data.
\end{enumerate}
The extracted problem, its falsifiers (benefit ordering by meaning at fixed bandwidth; no gradient at all; reports without slopes), and its observational-only safe test design are registered as OPEN-ANCHOR-BANDWIDTH-001 in the companion registry (S10). This subsection asserts nothing about what safety-signal or cognitive theory would predict here --- their primary texts remain unread --- only what Eq.~\ref{eq:companion} commits RRHM to.

\subsection{Two declared open terms, not absorbed}
\textbf{Anticipatory dread} has no term in this section either: the candidate reading (running the horizon channel offline on a stored record of the situation) adds one degree of freedom and one commitment --- the replayed bias must be the same $\delta_{cue}$ measured in situ (case C17); if they dissociate, dread is its own channel and must be Fisher-counted before use, not slipped in. \textbf{The social channel} (``people will see me fall'') has no term at all: $\Vset$ contains no observer, and writing one in without a declared, separately-measured term would be exactly the silent absorption this manuscript forbids elsewhere. Both gaps sit at the center of the agoraphobic phenotype, not its edge, and are left open by name.

\subsection{Kill criteria for the extension}
(i) K12 fires in adequately powered samples --- the companion is a signal, $k$ is dead. (ii) Agoraphobic horizon bias proves cue-locked --- the double dissociation dies. (iii) A present-but-non-intervening companion blocks recalibration as much as escape does --- the category-1/category-2 taxonomy is wrong. (iv) Anticipatory bias dissociates from in-situ bias --- the replay reading dies and dread requires new structure. (v) Companion presence changes reports but never physiological recovery slopes (C18) --- the co-regulation form of $k$ dies. (vi) Anchor benefit orders by symbolic meaning at fixed channel bandwidth (\S\ref{sec:panchor}) --- $k$ as a channel property dies. Any of these is reportable from the same measurement battery this manuscript already specifies.

\section{From horizon to governance: an intervention framework derived, not validated}
\label{sec:rgif}
\textbf{Fence.} This section assembles the treatment-relevant consequences of the whole manuscript --- the layered architecture (\S\ref{sec:layers}), the recalibration law (\S\ref{sec:recal}), the coupling extension and its fading rule (P9--P10) --- into one intervention framework. It is a \emph{derived treatment hypothesis}, Dr throughout, with no trial evidence of its own; established exposure-based treatment delivered by qualified clinicians remains the standard of care, and nothing here modifies it. The full companion statement is deposited as S11.

\subsection{The unit of treatment}
The implicit unit of most treatment framings is fear reduction; the implicit unit of exposure dosing is time-in-situation. This framework replaces both:
\begin{quote}
\textbf{The proposed unit of treatment is the successful information-bearing continuation}: an episode that runs to a benign completion \emph{while the phobic readout is still active}, with alternative policies still selectable before premature termination. In one line: \emph{treat the readout-to-policy lock, not merely the fear.}
\end{quote}
RRHM supplies the first half of that sentence (when the system predicts recovery will fail); the layered architecture supplies the second (how a readout becomes an action command, and where it can be made interruptible again).

\subsection{Seven ordered operations}
(1)~\textbf{Ground-truth gate}: separate actual hazard and physiological vulnerability from readout error before anything else --- untreated dysphagia is not a phobia, and a genuinely short causal horizon (vasovagal collapse) is repaired, not exposed. (2)~\textbf{Readout decomposition}: measure verbal fear, physiology, action tendency, horizon estimate, and actual termination as separate lanes; no single scalar. (3)~\textbf{Tape audit}: classify the retained history into harm, benign completion, and censored escape, and test for the $E\!\approx\!H$ mis-tally of Eq.~\ref{eq:lesion}. (4)~\textbf{Causal repair}: where causal recoverability is genuinely low, raise it first (applied tension as a $\Delta g$ move on Eq.~\ref{eq:adm}; consistent, at abstract holding, with the BII treatment literature \cite{ayala2009}). (5)~\textbf{Completion engineering}: construct episodes that reach benign completion without deception or bare reassurance; dose by completion count toward $n^{*}$ (Eq.~\ref{eq:nstar}), not by minutes endured. (6)~\textbf{Meta-readout and interruptibility}: train ``a readout is not a command'' --- a second-order readout of the collapse signal, with an alternative-policy set (tension, posture, pause, continue-small, signal clinician) kept available until the commitment point; this is not denial that the symptom is real. (7)~\textbf{Transport test}: remove therapist, context, safety cues, and body-state supports one at a time and measure whether the new readout survives (the generalization side of P10 and case C18).

\subsection{Three hard applications}
\textbf{Post-organic swallowing phobia (choking-phobia ARFID).} All pediatric application of this framework is interpretive only; management belongs to the ARFID clinical pathway with nutritional and growth monitoring. A child whose organic dysphagia has resolved but who has refused solid food for years presents a wall to belief-targeted work: the verbal lane already says ``safe''. The framework starts at the tape instead: hundreds of retained records, nearly all censored escapes tallied as if they were harms; treatment is completion engineering plus record reclassification --- the system does not learn ``I was wrong'', it learns that a danger readout can occur without becoming policy ($R_{\rm danger}\neq S_{\rm danger}$). A 2026 case report of lifelong solid-food refusal with choking phobia treated successfully with multimodal, inhibitory-learning-informed exposure is directionally consistent \cite{finn2026} (abstract holding). \textbf{BII with syncope and functional seizures.} ``Maximal exposure'' can be wrong from the first step: the vasovagal component is a genuinely short causal horizon, so causal repair precedes boundary sampling; and functional seizures are not absorbed into this framework at all --- they are evaluated and managed under their own practice guideline \cite{tolchin2026} (abstract holding), with the framework contributing only the interruptibility target afterwards: the goal is not ``no seizure'' but $R^{(1)}\not\Rightarrow$ a single irreversible policy. \textbf{Reflexive escape with intact explicit belief.} When escape is initiated within a few hundred milliseconds, before the verbal lane reports anything, there is no belief left to restructure; treatment is micro-continuations before the commitment point, moving $P(\text{escape}\mid R_{\rm phobic})$ below one \emph{without requiring} $R_{\rm phobic}=0$. The patient does not have to stop fearing in order to recover.

\subsection{The discriminating prediction, and its measures}
If the unit-of-treatment claim is right, then in a definable subset of patients \textbf{behavioral governability changes before phenomenal fear}: completion counts rise across sessions while fear ratings are initially unchanged, and fear declines only afterwards. A treatment course in which approach is only ever possible after fear has already declined would count heavily against the framework. Four preregistered measures follow (cases C19--C20 in S7~v2, registered before the fear--behavior desynchrony literature was searched): (i)~completion count versus exposure-minutes as outcome predictor (the clinical form of T2); (ii)~the \emph{interruptibility window} --- latency from first collapse readout to committed episode --- with the hypothesis that recovery lengthens or opens this window rather than silencing the warning; (iii)~the \emph{fear-lag}: the temporal dissociation between the governability curve and the fear-rating curve; (iv)~the \emph{transport curve} of operation~7. Kill criteria: approach never precedes fear reduction in any subset; completion count adds nothing over exposure-minutes (which also kills T2); the interruptibility window shows no relationship to clinical course.

\section{External-data stress testing before new recruitment}
\label{sec:external}
Existing public datasets cannot directly validate eRRH because they were not designed to manipulate causal recoverability. They can nevertheless test necessary implications and identify contradictions before new participant recruitment. Candidate resources, with accessions, are: SpiderPhy for multimodal spider-fear psychophysiology (OSF \url{https://osf.io/98cnd}; data descriptor \cite{lor2025}); PsPM-FER02 for return-of-fear dynamics (doi:10.5281/zenodo.5573760); CTX3 for contextual fear retention (doi:10.5281/zenodo.6418992); PsPM-PCF3 for reinforcement-probability competition (\url{https://zenodo.org/records/10849996}); SpiDa-MRI for neural--behavioral relations in spider fear (data descriptor \cite{zhang2025}; behavioral data OSF \url{https://osf.io/w4s2h/}; imaging OpenNeuro ds004630); and NIMH Data Archive/ABCD resources for developmental and longitudinal benchmarks (\url{https://nda.nih.gov/}). Access conditions differ across these resources and secondary analysis requires compliance with each source's own terms; the author has not obtained controlled-access approvals at the time of writing.

The rule is explicit:
\[
\boxed{\text{external compatibility}\neq\text{direct RRHM validation}}.
\]

\subsection{First empirical contact: a preregistered two-source necessary-implication test}
\label{sec:firstcontact}
The program's first contact with real human data was executed on 2 September 2026 under this section's own rule. The analysis plan was frozen as preregistered case C23 (S7~v2) \emph{before any outcome file was opened} --- only repository metadata and READMEs were inspected first --- with locked criteria, locked falsifiers, recorded prior guesses, and zero fitted parameters; full record deposited as S13. Two open datasets from unrelated paradigms, laboratories, and countries were downloaded raw and analyzed end-to-end:

\textbf{Source A (lane dissociation; Law 2 necessary implication).} A 2026 VR spider-exposure dataset (subjective SUDS ratings across eleven in-scene tasks; salivary cortisol at seven timepoints; PPG/smartwatch heart rate and HRV; $N{=}28$) \cite{vrspider2026}. Frozen criteria: group-level $|r|$ between mean subjective fear and physiological reactivity $\le0.70$, and both discordant quadrants (high-subjective/low-physiological and the converse) occupied by $\ge15\%$ of participants. Result: $r(\mathrm{SUDS},\ \Delta\mathrm{cortisol})=0.119$ ($n{=}27$), $r(\mathrm{SUDS},\ \Delta\mathrm{HR})=0.005$ ($n{=}28$); discordant quadrants 26--32\% each, with full-range subjective variance (means 0--8.6/10). Both criteria are met --- but independent review established that neither could realistically have failed: the recorded prior guess ($r\approx0.3$--$0.6$) already sits inside the pass region of criterion 1, and 25\% per discordant quadrant is the \emph{independence null}, so the 15\% floor of criterion 2 is passed automatically for any weak association; the low-reliability physiological indices (a two-timepoint cortisol difference; task-pooled SUDS) further push $r$ toward zero mechanically. The result stands as recorded, but it is confirmatory in form only and carries little evidential weight. \tier{finite\_diagnostic for the numbers}

\textbf{Source B (temporal-readout necessary implication) --- criterion failed; headline result withdrawn.} PsPM-HRA1: discriminant delay fear conditioning, visual CS, electric US, fixed CS--US interval of 3.5~s, trial-by-trial SCR, $N{=}20$ \cite{hra1data}. Frozen criteria: CS$+$~$>$~CS$-$ anticipatory amplitude, and temporally structured differential activity within the CS--US window. Under the frozen binary rule the amplitude criterion \textbf{failed} (mean $\Delta$SCR $=+0.013$, $t=1.51$, direction correct in 12/20; the frozen pipeline pools habituation and acquisition trials); an earlier edition scored this ``marginal'' and aggregated ``3.5 of 4 criteria'' --- both labels were invented after unblinding and are \textbf{retracted}. An earlier edition also reported the differential curve peaking at a median 3.48~s against the 3.5~s deadline as a striking pass of the timing criterion. Independent review identified this as a candidate window-edge artifact --- the analysis window's right edge \emph{is} the deadline, and SCR effector dynamics are still rising there --- and three of the four controls it ordered were then run (deposited with the companion code): extending the window to 6~s moved the median peak to 5.99~s, the new edge, in 16/20 participants; label-permutation nulls also piled at the edge (0.24); and 7/20 differentials were monotone or edge-pinned. \textbf{The timing claim is therefore withdrawn as a window-edge argmax artifact}; the fourth control (a different-SOA replication) remains open, and any future timing claim requires deconvolution of the SCR response function before peak-latency inference.

\textbf{What this does and does not establish.} Honestly scored, C23 yields no informative preregistered pass: Source A's criteria were met but could not have failed; Source B's amplitude criterion failed and its timing result was an artifact caught by independent review after initial over-reporting. What survives is methodological: the full loop --- freeze, test, over-read, independent challenge, control experiments, withdrawal --- ran to completion and is reported here in full, and the lesson (criteria must be failable, windows must not contain their own conclusion, effector latency must be modeled) is folded into every later freeze. Nothing here manipulated a causal deadline; the boxed rule above is unchanged, and the decisive experiment remains the only path to direct validation. \tier{finite\_diagnostic for the numbers; Dr for every bridge}
A direct test requires purpose-built independent measurement of $\eRRHc$, $\eRRHhat$, $\tcorr$, and transition outcome. A full external protocol, including per-dataset preregistered questions and a compatible-signature-versus-evidence ledger, is available as a companion document \cite{lahtee2026extprotocol}.

\subsection{Second round: four further preregistered contacts (C25--C28)}
\label{sec:secondcontact}
The lessons folded in from C23 --- criteria must be able to fail, no analysis window may contain its own conclusion, effector latency must be modeled or designed out, binary scoring only --- were applied to a second round on 2 September 2026. For each case the criteria, falsifiers, and prior guesses were frozen and committed to the public repository (with the commit hash as timestamp) \emph{before any outcome file was opened}; where a dataset's internal variable structure was unknown at freeze time, a declared two-stage procedure was used (conceptual criteria frozen first; the column mapping recorded from structure and codebook inspection only, as a separate commit, before any outcome value was computed). The question under test is the architecture's typed-lane commitment: whether experiential, physiological, and defensive-motor channels behave as one scalar fear state read out through different effectors (the rival M0) or as typed, dynamically coupled but non-identical lanes (M1). All numbers below are reproducible from the deposited scripts. \tier{finite\_diagnostic for the numbers; Dr for every bridge}

\textbf{C25 (lane typing within one paradigm).} A two-day differential fear-conditioning dataset with per-trial subjective ratings, skin-conductance responses, and fear-potentiated startle ($N{\approx}64$) \cite{iuhamburg2021}. Per subject, phase, and lane, the CS$+$/CS$-$ discrimination score was computed; the frozen primary criterion is a multitrait--multimethod inequality: mean within-lane, between-phase rank stability must exceed mean between-lane, same-phase correlation. Under M0 the ordering reverses (a shared momentary state should bind different lanes on the same day more tightly than one lane binds itself across a state-changing phase transition). Result: within-lane $+0.127$ versus between-lane $+0.081$ (PASS); end-of-extinction rating--SCR dissociation $r=-0.080$ with discordant quadrants 27\%/16\% (PASS). Honest annotation, recorded at scoring: the within-lane advantage is carried almost entirely by the rating lane ($\rho=+0.503$; SCR $-0.171$, startle $+0.048$), so a reliability-floor explanation could not be excluded from this dataset alone --- the point C28 was designed to settle.

\textbf{C26 (transport out of the fear domain).} A within-subject crossover comparing three psychosocial stress paradigms (Trier Social Stress Test, socially evaluated cold pressor, mental arithmetic) with profile-of-mood-states, salivary cortisol, and heart rate \cite{giles2015}. Frozen criteria transported verbatim from C23-A ($|r|\le0.70$; both discordant quadrants $\ge15\%$), with the variable-precedence list frozen before the spreadsheet was opened and the cortisol effector lag handled by using the post-task peak. Result: mood--cortisol $r=+0.064$ (quadrants 21\%/21\%), mood--heart-rate $r=+0.037$ (26\%/26\%); both PASS. The three stress-task observations per subject are not independent (disclosed at freeze); the case shows the dissociation pattern is not a phobia-specific curiosity.

\textbf{C27 (temporal predictability; the closest open analogue of the unpredictability construct in PRHC-U).} A large threat-of-shock action-preparation paradigm ($N{\approx}182$--$260$ analyzable) crossing threat with temporal predictability --- a cue whose blink rate either signals the moment of avatar onset or carries no temporal information --- with postural sway, cardiac interbeat intervals, and skin-conductance level \cite{koppold2026}. Analysis windows were fixed by design (1-s bins 1--8 of the 9-s anticipation phase; neither edge coincides with the threatened moment), honoring the C23 window rule. Three frozen criteria: (a) tonic freezing larger under unpredictable than predictable threat --- PASS, but small and carried by the cardiac channel (57\% of subjects; pre-declared as weakly diagnostic, since imminence rivals predict the same direction); (b) a steeper anticipatory temporal gradient under predictable threat --- \textbf{formally PASS but annotated FRAGILE at scoring}: it passes only through the frozen ``at least one marker, other not significantly reversed'' clause, the cardiac group mean trends in the opposite direction, and this criterion is recorded as carrying no citable support for the model beyond ``not falsified''; (c) lane dissociation of the per-subject unpredictability effect between cardiac and electrodermal channels --- $r=-0.029$, quadrants 24\%/24\%, $n=169$; PASS.

\textbf{C28 (the six-month multitrait--multimethod test; the principal result of the round).} A longitudinal dataset in which 120 individuals performed fear conditioning twice, six months apart, with per-trial ratings and skin conductance \cite{rankstab2022}. The frozen primary criterion repeats C25's inequality across a six-month gap: within-lane cross-timepoint stability must exceed between-lane same-day coupling. The comparison is asymmetric \emph{against} the prediction: two different lanes measured on the same day share context, state, and environment, yet the prediction is that one lane separated from itself by half a year will agree with itself more. Result: both the subjective-rating and the skin-conductance lane showed moderate rank-order stability across six months ($\rho=+0.448$, $n{=}92$, and $+0.389$, $n{=}97$), whereas their contemporaneous cross-lane association remained weak (mean $\rho=+0.156$; $+0.179$ at baseline, $+0.133$ at follow-up; $n{=}95$--$115$); difference $+0.262$; PASS. The low cross-lane convergence therefore cannot be attributed \emph{solely} to an absence of stable individual structure in the SCR lane --- the direct answer to C25's reliability caveat, stated no more strongly than that.

Immediately adjacent, because it changes the reading of that result: a preregistered secondary prediction \textbf{failed} and is conceded. The C25-derived reading that the rating lane intrinsically carries the stable individual structure (frozen as $\rho_{\mathrm{RAT}}-\rho_{\mathrm{SCR}}\ge+0.10$) was falsified at $+0.059$: both lanes carry temporal structure at a comparable level, and the privileged-rating-lane interpretation is withdrawn. C25's earlier asymmetry is \emph{consistent with} dataset- or measurement-specific differences; C28 does not identify its causal source. The concession leaves the architecture more symmetric than the earlier reading --- two lanes, each with its own persistence, weakly coupled --- and no lane is privileged as the carrier of the ``real'' state. The dissociation criterion also passed at both waves ($r=+0.225$ and $+0.182$; quadrants 20--29\%): the pattern of individual discordance recurs rather than appearing once.

Two ordered sensitivity checks (exploratory, deposited with the scripts; the frozen result is untouched): on the complete-case subsample with all four cells present ($n{=}92$), the inequality survives with sample composition equalized (within $+0.448$/$+0.373$ versus between $+0.171$/$+0.167$; $\Delta r = +0.242$); a participant bootstrap of $\Delta r$ (10{,}000 resamples, fixed seed) gives a 95\% interval of $[+0.03, +0.45]$ --- positive and excluding zero, but wide, so the claim is worded as ``substantially larger under the preregistered criterion,'' not as a significance statement. Two cautions bind every use of these numbers: the six-month correlations are test--retest stabilities that include true change and must not be used as reliability coefficients to disattenuate the cross-lane correlations; and a single-latent-factor model with unequal loadings, measurement error, and state effects is \emph{not} formally excluded by this pattern. That comparison was therefore run as its own preregistered case (C29), with its contamination declared: because C28 had already revealed the acquisition-phase correlations, the acquisition comparison was frozen as \emph{quantification-only}, and a genuinely semi-blind arm was added on the extinction phase, whose discrimination scores had never been computed in this workspace. Both a one-common-factor model (M0, $k{=}8$) and a two-lane-factor model with free inter-factor coupling (M1, $k{=}9$) were fitted by maximum likelihood under a frozen BIC decision rule. The two arms carry different weights and are reported so: the round yields \emph{one contaminated confirmation and one genuinely semi-blind model comparison}, and the empirical bite is in the semi-blind arm. Under the preregistered BIC bands, both contrasts constituted weak preferences for the two-lane model ($\Delta\mathrm{BIC}=+2.50$ acquisition, contaminated; $+7.83$ extinction, semi-blind). Separately, parametric calibration under the fitted single-factor model placed both observed contrasts beyond the simulated null range: across 500 simulations with M0 true, the median $\Delta\mathrm{BIC}$ was $\approx-4.5$ with 95th percentile near $-1.6$ to $-1.9$; 2 of 500 simulations reached the acquisition value and 0 of 500 the extinction value, giving conservative parametric-bootstrap tail probabilities under fitted M0 of $0.006$ and $0.002$ --- tail probabilities of that specific rival, not p-values proving the two-lane model. The two readouts are reported together and neither replaces the other. The author's recorded near-certain prior that the contaminated arm would reach the strong band was \emph{wrong} --- a forecast failure of the investigator, distinct from any failure of the model, kept in the score file as a calibration lesson about BIC's $\ln n$ penalty at $n{\approx}92$. An initially reported ``sign-unstable'' inter-lane coupling was corrected on review-ordered audit: factor models carry sign indeterminacy, and after canonicalization (reference loadings positive) the coupling is a \emph{consistent} weak positive $\hat\varphi=+0.22$ (acquisition) and $+0.31$ (extinction) --- the apparent flip was a parameterization artifact, caught before any mechanism story was built on it. The claim ceiling after C29: the preregistered two-lane architecture received model-level support over the \emph{specified} single-factor rival, including in a previously unexamined phase, while evidence remained below the study's preregistered strong-BIC threshold. ``A single latent factor is falsified'' is still \emph{not} said --- the rival that lost is one specified M0, not every conceivable scalar architecture.

\textbf{C30: the external transport test, and the boundary it drew.} The comparison was then taken to an external dataset with everything fixed in advance, deliberately giving the single-factor model a chance to win in a way that could not be repaired afterward. The preregistration (public commit, before any C30 computation) declared cell-by-cell which correlations C25 had already revealed, extended the architecture to all three lanes (ratings, SCR, startle; one common factor, $k{=}12$, versus three lane factors with free couplings, $k{=}15$), fixed the factor allocation, carried the frozen BIC bands, and stated the falsifier: $\Delta\mathrm{BIC}\le-2$ means the single-factor model wins externally, recorded as a failure of typed-lane transport with no post-hoc repair permitted. A companion parameter-transport arm (fit on one dataset, score frozen parameters on the other) was ruled \emph{not identifiable} at freeze time --- the two datasets differ in instruments, transforms, and above all in occasion semantics (a six-month retest versus a within-visit phase transition) --- and was declared unresolved rather than forced. Result: $\Delta\mathrm{BIC}=-10.03$ ($n{=}61$). \textbf{The falsifier triggered; the single-factor model wins externally, and the frozen consequence is recorded verbatim: typed-lane structure does not transport universally; the theory takes a claim boundary here.} The recorded prior had given this branch $p=0.2$; it happened. A diagnostic reading, labeled post-hoc and not a rescue: the three-lane fit is degenerate on this dataset because two of its three lanes carry near-zero within-lane cross-phase stability (known since C25), so lane factors have nothing to bind and their extra parameters (penalty $3\ln 61\approx12.3$) buy almost no likelihood. The boundary this draws --- stated for the record, not as a preregistered finding --- is that lane-factor structure is model-preferred where lanes are individually stable (C29) and BIC-rejected where they are not (C30-A); whether that boundary tracks measurement quality or true architecture is open, and it is the next preregisterable question. The evidence chain therefore closes at its honest edge: cross-sectional dissociation $\rightarrow$ replication $\rightarrow$ six-month within-lane stability $\rightarrow$ formal model-level support $\rightarrow$ an external transport test that the architecture \emph{lost}, on terms it set for itself in advance.

\textbf{Standing of the round: the challenge ledger.} The round is reported as a ledger of frozen challenges and their consequences, not as a score:

\begin{table*}[!t]
\centering
\caption{Challenge ledger for the second-round open-data contacts. Every prediction was frozen (public commit hash) before the outcome file was opened; no failed entry was rescued. \tier{finite\_diagnostic for the numbers; Dr for the consequence column}}
\label{tab:ledger}
\small
\begin{tabularx}{\textwidth}{p{0.05\textwidth}p{0.29\textwidth}p{0.20\textwidth}p{0.07\textwidth}X}
\toprule
\textbf{Case} & \textbf{Frozen prediction} & \textbf{Result} & \textbf{Verdict} & \textbf{Consequence} \\
\midrule
C25.1 & within-lane cross-phase stability $>$ between-lane same-phase & $+0.127$ vs $+0.081$ & PASS & initial typed-lane signature; reliability caveat recorded (rating-lane-driven) \\
C25.2 & end-of-extinction rating--SCR dissociation & $r=-0.080$; quadrants 27/16\% & PASS & dissociation inside one paradigm \\
C26 & dissociation transports out of the fear domain & $r=+0.064$, $+0.037$ & PASS & pattern is not phobia-specific \\
C27a & tonic freezing larger under unpredictable threat & cardiac $57\%$ of subjects, small & PASS & direction only; weakly diagnostic (pre-declared) \\
C27b & steeper anticipatory gradient under predictable threat & passes only via the one-marker clause; cardiac trends opposite & PASS, \textsc{frag\-ile} & recorded as carrying no citable support; model-specific signature undemonstrated \\
C27c & unpredictability effect dissociates across lanes & $r=-0.029$; quadrants 24/24\% & PASS & lane typing extends to the manipulation's effect \\
C28.1 & within-lane six-month stability $>$ same-day cross-lane convergence & $+0.418$ vs $+0.156$ & PASS & stable, typed, weakly coupled lanes; C25 caveat answered \\
C28.2 & dissociation recurs at both waves & $r=+0.225$, $+0.182$; quadrants occupied & PASS & longitudinal replication, not a snapshot \\
C28.3 & rating stability exceeds SCR stability by $\ge0.10$ & $+0.059$ & \textbf{FAIL} & privileged-rating-lane interpretation withdrawn; architecture read as symmetric \\
\bottomrule
\end{tabularx}
\end{table*}

Across C23-A, C25, C26, C27, and C28 the same empirical object has now appeared in five independent datasets from three contexts (phobic exposure, fear conditioning, psychosocial stress) and, in C28, across a six-month retest --- and after C28 it is a stronger object than a low correlation: \emph{lane-specific individual structure is longitudinally stable while contemporaneous cross-lane convergence remains weak}, with recurrent individual discordance. This is a necessary implication of the typed-lane architecture, not of RRHM's horizon mechanism specifically, and it neither proves the readout architecture nor formally falsifies a one-factor model; the boxed rule stands --- compatibility, not validation. \tier{finite\_diagnostic for the numbers; Dr for the M0/M1 reading}

\section{Compatibility audit of preregistered cases C6--C15}
\label{sec:audit}
Ten compatibility cases were preregistered before their literature was searched, with four fixed scoring labels (SUPPORT-D / SUPPORT-N / TENSION / NO-DATA), a named rival per case, a prior guess recorded per case, and the binding rule that anything not accessible in full text is scored NO-DATA, never inferred from an abstract. The preregistration is deposited verbatim as S7; cases C1--C5 were scored after looking, are recorded there as post-hoc, and are not counted. The audit was executed at v9 against PubMed (E-utilities) and PubMed~Central open-access full text; that access boundary is itself part of the result and is stated in the table.

\begin{table*}[!t]
\centering
\caption{Executed compatibility audit, C6--C15. Scored under the preregistered rules of S7: full-text access required for any label above NO-DATA. ``Prior'' is the guess recorded in S7 before searching. Labels: S-D = SUPPORT-D, S-N = SUPPORT-N, T = TENSION, ND = NO-DATA. \tier{finite\_diagnostic for the search results; the scores are rule applications, not evidence for RRHM}}
\label{tab:audit}
\small
\begin{tabularx}{\textwidth}{p{0.04\textwidth}p{0.24\textwidth}p{0.09\textwidth}p{0.09\textwidth}X}
\toprule
\textbf{\#} & \textbf{Case} & \textbf{Prior} & \textbf{Score} & \textbf{Basis} \\
\midrule
C6 & exit constraint changes policy, not timing & S-N & ND & The discriminating sources (multi-centre MRI termination rates against matched distress) are not full-text accessible from this environment; an accessible cohort reports 75\% discomfort prevalence against 6.5\% negative overall experience (unscored observation), without the matched-fear termination comparison the rule requires. \\
C7 & applied tension is BII-specific & S-D & ND & The trial literature returned applied-tension trials only in BII (\"Ost \& Sterner 1987; \"Ost 1989, 1991); no non-BII applied-tension trial surfaced, i.e.\ the weakest of the preregistered ``consistent if'' conditions held at search level only. Rescored on independent review: search-level absence cannot license a label above NO-DATA under rule 4; the absence is recorded as an unscored observation. \\
C8 & interoceptive load shifts transition at matched fear & ND & ND & No study covarying fear and reporting residual effects on termination timing was found; prior guess correct. \\
C9 & remitters retain residual cue-specific miscalibration & S-N & ND & Return-of-fear and follow-up studies exist but none accessible in full text reports a residual cue-specific index predicting relapse. \\
C10 & forced non-avoidant sampling recalibrates without therapy & S-N & ND & Occupational-epidemiology searches returned off-target results; the named source set is silent here. \\
C11 & action count inert, deadline geometry not & T & ND & The expected choice-availability counter-evidence did not surface in the named source set; the case cannot be scored against inaccessible controllability full texts. \\
C12 & onset follows failed correction, not high fear & ND & ND & Onset-event studies (e.g.\ NCS life-event analyses) do not code coping-action availability at the index event; prior guess correct. \\
C13 & acrophobia timing tracks correction window beyond sway & ND/T & ND & Sway studies exist (e.g.\ acrophobic visual-exploration work) but their full texts were not accessible; no study separates a corrective-window estimate from sway. \\
C14 & dental/choking show periodic corrective windows & ND & ND & No study of phase-locking of stop requests to swallow or respiratory cycles was found in any cell searched; prior guess correct, and the gap is independently corroborated by a parallel structured search this year that found no clamp-duration--distress timing study at all. \\
C15 & panic shows the construct without the cue lock & T & ND & The closest accessible comparisons are panic versus social phobia and mixed childhood anxiety under CO$_2$; no accessible panic-versus-specific-phobia comparison addresses cue-locked versus generic bias structure. \\
\bottomrule
\end{tabularx}
\end{table*}

Three things the executed audit establishes, none of them support for RRHM. First, the score distribution after review-rescoring is ten of ten NO-DATA (C7's initial SUPPORT-N could not be licensed by search-level absence under rule 4): on what is accessible under the preregistered rules, the existing literature is almost entirely silent on the model's distinctive predictions, which is consistent with the manuscript's own claim that no existing study manipulates a causal engagement deadline---and is exactly what the audit's ``expected label for several cases'' clause anticipated, though not at this rate. Second, the prior guesses were correct in three of ten cases (C8, C12, C14), half-right in one (C13, whose frozen prior was ``NO-DATA or TENSION''), and wrong in six --- with the wrong guesses sharing one direction: the author over-predicted that the literature would carry usable signal. Under S7's own framing this is the informative outcome---the model's author could not predict where the theory would meet resistance, mostly because the resistance has never been measured. Third, the NO-DATA rate is partly a property of the audit environment (PubMed abstracts plus PMC open access, a bounded query set per case), not only of the literature; a full-library pass could raise several cases above NO-DATA in either direction, and the preregistration binds any such future pass to the same labels and priors.

\section{Neural implementation: temporal predictions, not a region list}
RRHM does not posit a single neural locus. Candidate circuitry includes interoceptive and sensory integration, valuation and action selection, controllability-related prefrontal systems, striatal policy selection, amygdala/BNST threat processing, periaqueductal gray defensive organization, brainstem autonomic control, and vestibular/postural networks. Such a list is not evidence for the theory.

The falsifiable neural claim is temporal: trialwise neural and physiological state estimates associated with horizon contraction should change \emph{before} defensive transition and should improve prediction beyond conventional threat and control variables. EEG is suitable for temporal precedence; fMRI can test distributed network dynamics after the behavioral phenomenon is established. A second, sharper claim follows from \S\ref{sec:fid}: because escape-decision work dissociates fast reactive from slow cognitive escape circuitry \cite{qi2018,fung2019}, a manipulation of the engagement deadline at fixed threat kinematics should modulate the slow route selectively. If deadline effects appear only in fast reactive responding, the engagement horizon is not the variable proposed.

\section{Safety and ethics}
\label{sec:safety}
Initial experiments should use healthy adults or non-clinical fearful participants and low-risk reversible paradigms. All paradigms proposed here are adult-only in their current form: any pediatric extension requires separate ethical design, developmentally validated temporal-estimation measures (every measurement module here is a time estimate, and children's interval timing differs systematically), and excludes the delayed-exit manipulation; pediatric case material (V1) is interpretive only. Investigators should not intentionally induce hypoxia, airway obstruction, syncope, arrhythmia, seizure, or uncontrolled panic. Participants must retain a genuine emergency stop; the manipulation concerns the geometry of engagement-preserving action, not removal of safety protection. The delayed-exit condition manipulates a disclosed, bounded waiting interval for the ordinary exit, never the emergency stop, and the distinction must be explained during consent.

BII paradigms require syncope precautions. Height paradigms require objective physical protection. Choking or vomiting paradigms should use validated representational cues rather than genuine airway or aspiration risk. Dental paradigms should avoid prolonged maximal mouth opening and intentional respiratory restriction. Clinical safety overrides theoretical discrimination. Every paradigm proposed in this manuscript requires review and approval by a competent research ethics committee before any human data collection; this section is a design constraint, not an approval, and the present manuscript reports no human or animal experimentation.

\section{Hard falsification gates}
\label{sec:gates}
Earlier editions listed ten kill criteria in words. A criterion that cannot fail is not evidence, and a criterion with no number cannot be shown to fail; each gate below therefore carries a numeric pass predicate, a positive control (a synthetic dataset generated by ADM with a true eRRH effect, which the gate must pass) and a negative control (a synthetic dataset generated by a rival model with no horizon channel, or by permutation, which the gate must fail). A gate that passes both controls, or fails both, is a convention rather than a test and must be redesigned before the study runs. Thresholds are stated as defaults to be fixed at preregistration; they are deliberately concrete so that a reviewer can dispute a number rather than a mood.

\begin{table*}[!t]
\centering
\caption{The eleven numeric kill gates K1--K11. Each is a two-sided test carrying a positive and a negative control; the theory is rejected, reduced, or substantially revised if a gate fires against it in repeated studies. K11$'$ is not a twelfth gate but the transition-timing-only variant of K11, retained in the table because \S\ref{sec:sim} shows the two variants behave differently. Defaults shown; final values preregistered.}
\label{tab:gates}
\small
\begin{tabularx}{\textwidth}{p{0.05\textwidth}p{0.26\textwidth}X p{0.24\textwidth}}
\toprule
\textbf{Gate} & \textbf{Claim under test} & \textbf{Fires against RRHM when} & \textbf{Controls} \\
\midrule
K1 & eRRH is measurable before the behavior it predicts & any horizon estimate requires the transition outcome, or Module A/B/C$'$ timing order is violated in $>1\%$ of trials & pos: ADM data with outcome-blind estimates; neg: estimate constructed from outcome \\
K2 & Parameter recovery & recovery $r<0.70$ for $\eRRHhat$ or $\tcorr$ on ADM synthetic data at the planned $N$ and trial count, \emph{with within-trial probes and Module~F in place} & pos: ADM with separable parameters; neg: design with deliberately collinear $\eRRHhat$ and control. Run: $r=0.70$ without probes, $0.88$ with (\S\ref{sec:sim}) \\
K3 & Discriminant validity & a one-sided \emph{minimum-effect} test (not TOST): the gate fires when the lower confidence bound of the reliability-disattenuated $|r|$ between $\eRRHhat$ and perceived control, ASE, expectancy or timing distortion is \emph{demonstrated to exceed} 0.85. Run: void on observed $r$ at any $N$; two-sided disattenuated at $N\ge220$ (\S\ref{sec:simremaining}) & pos: ADM with independent control channel; neg: simulated single-channel data \\
K4 & Causal deadline effect & no effect of $D$ on transition timing with threat, exit, control, action efficacy and count matched; $\ge2$ independent samples & pos: ADM deadline manipulation; neg: identical task with $D$ non-binding \\
K5 & Held-out predictive value & $\Delta\mathrm{ELPD}(M_{RRHM}-\text{best rival})<4$ with SE such that the difference is not $\ge2$~SE; and $\Phi(\Tdt;\eRRH)<0.10$ under preregistered priors. Models must be fitted \emph{hierarchically}: per-subject fits inflate the SE past the threshold and make the gate uninformative in both directions (\S\ref{sec:sim}) & pos: ADM data; neg: permuted-deadline dataset, which must return $\Phi\approx0$ \\
K6 & Cross-phobia generalization & a shared eRRH parameter fails measurement invariance across $\ge3$ classes, or class-specific parameters are required for fit, \emph{on the joint hazard-plus-probe likelihood}. Run: uninformative in the timing-only form at feasible $N$; the joint form is committed but not yet shown two-sided (\S\ref{sec:simremaining}) & pos: multi-class ADM with shared horizon; neg: multi-class data with class-specific latent variables \\
K7 & Prospective onset & baseline cue-locked $\RCE$ carries an incremental effect below $\beta=0.30$ log-odds per SD over established risk factors, with the CI excluding $0.30$. Stated as an effect-size floor rather than as $\Delta$AUC: at a 2\% incidence, $\Delta$AUC $\ge0.02$ is unreachable below that $\beta$ at any cohort size (\S\ref{sec:sim}) & pos: simulated cohort with horizon-driven onset; neg: cohort with onset driven by baseline fear alone \\
K8 & Remission and relapse & change in $\eRRHhat$ adds no prediction of functional remission or relapse beyond fear and avoidance ($\Delta$ELPD $<4$ at the preregistered $N{=}80$; threshold and $N$ fixed jointly, and an $N$-sweep remains owed before ``usable as written'' generalizes) & pos: ADM recalibration cohort; neg: cohort with fear-only recovery \\
K9 & One construct, not many & subtype fits require different latent definitions rather than different ADM parameters, under the same hierarchical 2-SE rule as K5. Run: two-sided at $N=60$ (fires 0.02/0.43); the fixed threshold leaks at $N=120$, so threshold and $N$ are preregistered jointly (\S\ref{sec:simremaining}) & pos: shared-equation multi-subtype data; neg: data generated with distinct per-subtype constructs \\
K10 & Cue-locked specificity & cue-locked and non-cue $\RCE$ are equivalent within TOST bounds ($d=0.40$, or $d=0.30$ with $N\ge220$) in phobia, or the same bias magnitude appears across unrelated disorders with no cue structure & pos: ADM with cue-conditioned bias; neg: ADM with generic bias. At $d=0.30$ and $N=90$ the gate fires on its own negative control only 0.27 of the time (\S\ref{sec:sim}) \\
K11 & Not generic timing & the additive eRRH bias is absorbed by $M_{Timing}$ ($\Delta$ELPD $<4$), where $M_{Timing}$ is constrained by within-trial probe slopes and by temporal bisection and interval reproduction measured under matched arousal. \emph{The gate is void unless probe slopes are collected}: on transition timing alone the two models are exactly likelihood-equivalent and the gate fires whether or not RRHM is true. With probes it is two-sided, and has been shown to be so by simulation (\S\ref{sec:k11}) & pos: ADM with additive horizon bias and normal clock; neg: data with clock distortion only \\
\bottomrule
\end{tabularx}
\end{table*}

A theory that can absorb every outcome after the fact is not a useful medical theory. Two rules make that discipline enforceable rather than declarative: no gate may be counted as evidence unless both of its controls have been run and behaved as specified, and no gate outcome may be reported as a headline verdict if only one side of the control pair exists. Three outcomes are permitted per gate---fires, passes, or uninformative (the controls did not separate)---so that ``no evidence'' and ``evidence against'' cannot be silently merged.

Those rules are cheap to state and were expensive to obey. The next section reports what happened when the controls were actually run.

\section{Pre-study simulation of the falsification battery}
\label{sec:sim}
A battery of kill gates is a claim about a measurement design, and a claim about a design can be tested before any participant is recruited. Ten of the eleven gates of \S\ref{sec:gates} are testable this way; K1 is a procedural ordering rule (predictors must precede the outcome) and is enforced by protocol, not by simulation. All ten have now been run on synthetic cohorts generated by ADM, with both of the controls the gate-typing rule requires. Six were run first (K2, K4, K5, K7, K10, K11): of those, one behaved as written, one needed only its fitting method specified, and four failed---two on their threshold, one on the design's information content, and one on the algebraic form of the construct itself. The remaining four (K3, K6, K8, K9) were run at v9 and are reported in \S\ref{sec:simremaining}. The manuscript above is the revised version; this section reports what the revision cost and why, because a battery whose failures are hidden is worth less than one that reports them.

\subsection{Calibration and what the cohorts are not}
Cohorts are synthetic. No public dataset was analyzed: the resources named in \S\ref{sec:external} could not be obtained within this work, and none of them manipulates a causal engagement deadline in any case. What is anchored to published data is the \emph{parameters}. The hazard intercept and slope were tuned until a simulated cohort reproduced escape-decision benchmarks from intracranial and fMRI FID studies: transition on 85\% of trials against a published escape rate of $80.6\pm10.9$\% (fast attack) and $81.2\pm12.5$\% (slow attack), and premature transition on 13\% of trials against a published premature-escape rate of $14.4\pm13.3$\% in the slow-attack context \cite{zhang2025escape}. Cohort size follows the same literature ($N\approx28$ in \cite{qi2018}). Trials are simulated in 0.1~s steps and censored at 15~s. \tier{finite\_diagnostic, planning tier}

\subsection{The finding that forced the revision}
Write the multiplicative rival explicitly as $\eRRHhat_t=K\,D-t$ --- a gain on the \emph{onset} readout with unit-rate internal countdown. (Independent review noted this is a rival model specification distinct from Eq.~\ref{eq:recordlaw}'s flowing form, whose literal multiplicative reading $K\,(D-t)$ instead satisfies $a_1=-a_2$ and nests in the clock family with $K$ confounded into $\beta_1$; the RCE decomposition of \S\ref{sec:anchor} holds at onset. The practical conclusions below survive under either parametrization.) The hazard's linear predictor for this rival and for a clock-rate distortion:
\begin{align*}
\text{horizon gain } K:&\quad [\beta_0+\beta_1\tau] + [-\beta_1 K]\,D + [\beta_1]\,t,\\
\text{clock gain } c:&\quad [\beta_0+\beta_1\tau] + [-\beta_1]\,D + [\beta_1 c]\,t.
\end{align*}
Both are $a_0+a_1D+a_2t$, and each reaches every admissible $(a_0,a_1,a_2)$. They are the same likelihood family. Fitted to the same synthetic cohorts the two models return identical maxima to four decimals---$-8331.294$ on horizon-generated data, $-8087.079$ on clock-generated data---with the parameters simply exchanged ($\hat K\hat\beta_1$ of one equalling $\hat c\hat\beta_1$ of the other). Held-out $\Delta$ELPD is $0.00$ in both directions, so Gate~K11 fires against RRHM in every run, including every run in which RRHM is true.

The additive form is not affected: it constrains $a_1=-a_2$, which the clock model does not, so the two remain distinguishable. This is why \S\ref{sec:prhc} now commits PRHC to the additive form. Nothing empirical was given up; an unfalsifiable way of writing the same claim was.

Commitment to the additive form restores identifiability but does not by itself win the comparison. Fitted hierarchically to transition timing alone, the additive model still loses to the timing model on held-out prediction even when the additive model is the one that generated the data ($\Delta$ELPD $-39\pm10$), because the timing model is strictly the more flexible of the two and absorbs between-subject heterogeneity the constrained model cannot. Identifiability is necessary and not sufficient: what decides the comparison is the within-trial probe slope, which is $-1$ under a horizon bias and $-c$ under a clock distortion. Gate~K11 is therefore conditional on those slopes being collected, and is void without them.

\subsection{What the design can and cannot identify}
The design-level Fisher information over $(K,c,\tcorr,\beta_0,\beta_1)$ has numerical rank 3 of 5 for a deadline-only design---with two, three, or four deadline levels alike. Two directions carry no information at all. One is the horizon-versus-clock trade-off above. The other is exact and independent of it: shifting $\tcorr$ by $\delta$ and $\beta_0$ by $-\beta_1\delta$ leaves the likelihood unchanged, so correction latency is unidentifiable from transition timing at any sample size, and simulated recovery of $\tcorr$ returns $r\approx0.01$--$0.08$. Mid-trial probes raise the rank to 4 by fixing the first direction; Module~F raises it to 4 by fixing the second; only both together give full rank (Fig.~\ref{fig:rank}). Recovery of the horizon parameter tracks the same boundary: $r=0.70$ without probes---exactly at Gate~K2's threshold, with a $+0.11$ upward bias---and $r=0.88$ with them.

\begin{figure}[!t]
\centering
\includegraphics[width=0.98\columnwidth]{fig2_rank.pdf}
\caption{Rank of the design-level Fisher information over $(K,c,\tcorr,\beta_0,\beta_1)$. Adding deadline levels does not help; only within-trial probes together with an independent correction-latency measure reach full rank. Red bars are designs the identifiability gate of \S\ref{sec:irhab} rejects. \tier{finite\_diagnostic}}
\label{fig:rank}
\end{figure}

\subsection{Gate-by-gate outcome}
Table~\ref{tab:sim} reports every simulable gate against both controls, in the order the outcomes are discussed. Counting by outcome: \textbf{K4} discriminated cleanly as written and needed no change; \textbf{K5} discriminated once its fitting method was specified as hierarchical; \textbf{K2} and \textbf{K10} failed at the planned sample size and were repaired by making two measurement modules mandatory and by widening one equivalence bound; \textbf{K7} failed on a threshold property demonstrated in the simulated cohort (the $\beta\to\Delta$AUC mapping depends on the baseline model's discrimination, not on arithmetic alone) and was restated as an effect-size floor; and \textbf{K11} was void in both directions and is the subject of \S\ref{sec:k11}.

\begin{table*}[!t]
\centering
\caption{Pre-study run of the falsification battery on ADM-generated cohorts calibrated to published escape-decision benchmarks. ``Fires'' is the proportion of runs in which the gate would reject RRHM. A usable gate fires rarely on its positive control and often on its negative control; a gate that behaves alike on both carries no evidence.}
\label{tab:sim}
\small
\begin{tabularx}{\textwidth}{p{0.05\textwidth}p{0.20\textwidth}p{0.15\textwidth}p{0.15\textwidth}X}
\toprule
\textbf{Gate} & \textbf{As written} & \textbf{Positive control} & \textbf{Negative control} & \textbf{Verdict and action taken} \\
\midrule
K4 & deadline effect on transition timing & fires 0.00 ($N=24$, 32, 48; detection 1.00) & fires 0.98 & Usable as written, and at a quarter of the planned $N$. Unchanged. \\
K2 & recovery $r\ge0.70$ & $r=0.70$ (bias $+0.11$) without probes; $r=0.88$ with & $r(\tcorr)\approx0$ in both & Fails at the planned $N$ for a design reason, not a theory reason. Probes and Module~F made mandatory. \\
K5 & $\Delta$ELPD $\ge4$ over the no-deadline null & fires 0.00, $\Delta$ELPD $754\pm28$ & fires 0.80, $\Delta$ELPD $-45\pm13$ & Usable, but only under a hierarchical fit: with per-subject fits the SE reached 185 and the gate fired 0.36 on its own positive control. Fitting method now specified. \\
K11 & $\Delta$ELPD $\ge4$ over $M_{Timing}$, multiplicative PRHC & fires 1.00 ($\Delta$ELPD $-0.03$) & fires 0.50 & Void as written: exact likelihood equivalence (above). PRHC committed to the additive form. \\
K11$'$ & same, additive PRHC, transition timing only & fires 1.00 ($\Delta$ELPD $-39\pm10$) & fires 0.88 & The additive form is identifiable in principle but still loses on held-out prediction to a strictly more flexible timing model. Gate made conditional on within-trial probe slopes, which is what breaks the tie. \\
K10 & TOST equivalence, $d=0.30$ & fires 0.00 & fires 0.00 ($N=40$), 0.05 ($N=60$), 0.27 ($N=90$) & Cannot fire at the planned $N$. Bound widened to $d=0.40$, or $N\ge220$ retained at $d=0.30$. \\
K11$''$ & K11 re-run: additive PRHC \emph{plus} within-trial probe slopes & fires 0.00 ($\Delta$ELPD $+565\pm40$) & fires 1.00 ($\Delta$ELPD $-81\pm21$) & Two-sided for the first time. The repair works; \S\ref{sec:k11}. \\
K7 & $\Delta$AUC $\ge0.02$ over risk factors & detection $\le0.21$ at $N=10{,}000$ & --- & Threshold, not cohort, is the obstacle: at 2\% incidence $\Delta$AUC $\ge0.02$ requires $\beta\ge0.30$ log-odds per SD ($\beta=0.22\Rightarrow\Delta$AUC $0.012$). Restated as an effect-size floor. \\
\bottomrule
\end{tabularx}
\end{table*}

\subsection{Gate K11 re-run with probe slopes}
\label{sec:k11}
The repair proposed above---commit PRHC to its additive form and probe the reported horizon twice inside a trial---was, in the previous version of this manuscript, argued from the design's information content and not demonstrated. It is demonstrated here, because an argument that a gate \emph{would} work is not evidence that it does.

\textbf{The model-free statistic.} A horizon bias shifts the reported remaining interval down but leaves it falling at unit rate; a clock distortion leaves it starting correctly and makes it fall at rate $c$. A per-subject regression of probe reports on probe time therefore separates the two without fitting anything. It does: cohort-mean slopes of $-0.31$ (SD $0.09$) under an additive horizon bias and $-1.13$ (SD $0.04$) under a clock gain of $1.30$, giving complete cohort-level separation ($d'=11.3$ under the pooled-SD convention on unrounded inputs; classification accuracy $1.00$ over 200 simulated cohorts of $N=40$) --- under a correctly specified generative model with no measurement mismatch: an upper bound, never a power estimate.

\textbf{A caution the numbers force.} Those slopes are not $-1.00$ and $-1.30$. They are attenuated, and the reason is survivorship: a probe is only observed on a trial still running at the probe time, and the surviving trials are those with the larger reported horizon. Attenuation falls as the probes move earlier---at probe times of 1.0 and 2.5~s the additive-bias slope reads $-0.32$ (reported as $-0.31$ above; the discrepancy is rounding of unrounded run outputs and is noted rather than silently unified, since the generating engine is not in this deposit), at 0.5 and 1.5~s it reads $-0.71$, and at 0.3 and 1.0~s it reads $-0.95$ against a true $-1.00$---but very short spacing reintroduces noise, because the lever arm shrinks relative to report error. Probes are therefore specified at 0.3 and 1.0~s after trial onset (\S\ref{sec:irhab}), and the slope is treated as a \emph{discriminating statistic}, never as an estimate of $c$. Any published estimate of the clock gain must come from a model that represents the censoring, not from the raw regression.

\textbf{The gate.} With probes in the design and both models fitted hierarchically on the joint hazard-plus-probe likelihood, Gate~K11 becomes two-sided for the first time: held-out $\Delta$ELPD $=+565\pm40$ when the additive horizon bias generated the data, with the gate firing in $0.00$ of runs, and $-81\pm21$ when the clock distortion generated it, with the gate firing in $1.00$ of runs. The same gate, on the same theory, at the same sample size, fired in every run in both directions before the repair. What changed was not the theory but two probe reports per trial.

\textbf{What is still assumed.} The probe term in the likelihood treats a report as conditionally independent of survival given the parameters, which is the same simplification that produces the attenuation above; it is shared by both models in the comparison and so does not favor either, but it means the fitted $\sigma$ of the probe measurement should not be interpreted as report precision. Modeling the selection explicitly is the obvious next refinement and is listed in \S\ref{sec:limits}.

\subsection{The remaining four simulable gates: K3, K6, K8, K9}
\label{sec:simremaining}
Earlier versions of this section claimed, wrongly, that only six gates were simulable; the four below always were, and a battery that reports six of ten runs is not a run battery. They were run at v9 on an independent re-implementation of ADM and the hazard of Eq.~\ref{eq:hazard} written from this manuscript's own specification (the S1 engine was not used; code and run logs are deposited as S8), with simplified likelihoods declared per gate: Gaussian log predictive density on log transition latency for the timing components; leave-one-subject-out cross-validation for K3/K6/K9 and grouped 10-fold (eight subjects per fold, $N{=}80$) for K8; probes at 0.3 and 1.0~s; censoring at 15~s; 200 cohorts for K3, 100 for K6/K8/K9 unless noted. The ``2-SE rule'' used below means: the summed held-out ELPD difference must exceed 4 \emph{and} exceed twice its standard error estimated as $\mathrm{sd}(\text{per-fold differences})\times\sqrt{n_{\text{folds}}}$. Independent review noted that the K6/K9 fits are pooled OLS over hierarchically generated data --- the artifact class recorded at \S\ref{sec:sim} error \#2 --- so both verdicts are \textbf{provisional}; a post-review two-stage subject-level refit (per-subject slopes, then subject-level tests) did not rescue either gate (K6 fires 0.38/0.44; K9 probe-slope contrast 0.72/1.0), and a full mixed-effects implementation remains owed. Review also noted K6's negative control varies hazard slopes rather than horizon parameters, so the ``timing-only designs cannot carry this'' conclusion is specific to that parameterization; and K3$'$'s Fisher interval understates the sampling error of a disattenuated ratio estimator (false-kill 0.18 on its own positive control at $N=40$), so the $N\ge220$ requirement may be an artifact of the split-half correction and needs a bootstrap interval before it constrains any design. \tier{finite\_diagnostic, planning tier}

\begin{table*}[!t]
\centering
\caption{Pre-study run of the four remaining simulable gates, on an independent re-implementation of ADM (S8). Conventions as in Table~\ref{tab:sim}: a usable gate fires rarely on its positive control and often on its negative control.}
\label{tab:sim2}
\small
\begin{tabularx}{\textwidth}{p{0.05\textwidth}p{0.22\textwidth}p{0.15\textwidth}p{0.15\textwidth}X}
\toprule
\textbf{Gate} & \textbf{As written / as repaired} & \textbf{Positive control} & \textbf{Negative control} & \textbf{Verdict and action taken} \\
\midrule
K8 & $\Delta$ELPD $<4$ for adding $\Delta\eRRHhat$ to fear+avoidance & fires 0.01 ($\Delta$ELPD $+10.4\pm3.8$) & fires 1.00 ($\Delta$ELPD $-0.8\pm1.1$) & Usable as written. Unchanged. \\
K3 & TOST at $|r|=0.85$ on \emph{observed} correlations, $N=40$ & fires 0.00 & fires 0.00 (mean $|r|=0.76$) & Void as written: measurement attenuation puts even a same-latent relabel below the bound, so the gate cannot fire on its own negative control at any feasible $N$. \\
K3$'$ & same, on split-half-disattenuated $r$ & fires 0.18 ($N=40$); 0.09 ($N=120$); 0.03 ($N=220$) & fires 0.59 ($N=40$); 0.82 ($N=120$); 0.89 ($N=220$) & Two-sided only with reliability correction and $N\ge220$ --- the same $N$ at which K10 retained its $d=0.30$ bound. Disattenuation made part of the gate; $N\ge220$ or a latent-variable formulation required. \\
K6 & measurement invariance on transition timing, $N=60$, 16 trials & fires 0.05 & fires 0.21 & Cannot discriminate. \\
K6$'$ & same, 2-SE rule, $N=120$, 32 trials & fires 0.10 & fires 0.32 & Still uninformative at double the planned density: timing-only invariance testing does not separate shared from class-specific horizons in this design. Committed, untested, to the joint hazard-plus-probe form; flagged as such for preregistration. \\
K9 & per-subtype construct wins $\Delta$ELPD $\ge4$, $N=60$ & fires 0.23 & fires 0.89 & As written it fires spuriously on its own positive control: the more flexible family wins on flexibility, the artifact already seen at K5/K11$'$. \\
K9$'$ & same, with the K5-style 2-SE rule & fires 0.02 ($N=60$); 0.17 ($N=120$) & fires 0.43 ($N=60$); 0.95 ($N=120$) & The 2-SE rule controls the false fire at $N=60$ but leaks again at $N=120$ because a fixed $\Delta$ELPD $\ge4$ does not scale with $N$. Threshold and $N$ must be preregistered jointly. \\
\bottomrule
\end{tabularx}
\end{table*}

Three of the four therefore repeat, on fresh gates, exactly the failure classes the first six produced: an equivalence bound unreachable under realistic measurement noise (K3, as K10), a flexible family beating the true model unless an SE rule constrains it (K9, as K5 and K11$'$), and a timing-only design that cannot carry the discrimination at all (K6, as K11). One lesson is new: the K9 leak at $N=120$ shows that a fixed $\Delta$ELPD threshold is not sample-size invariant, so every ELPD-thresholded gate in Table~\ref{tab:gates} must have its threshold fixed jointly with $N$ at preregistration. K6 is the one gate in the battery that has not been shown two-sided in any form; until the joint-form run exists, cross-phobia invariance (and with it part of K9's subject matter) rests on a commitment, not a demonstrated test.

\subsection{Three errors in the simulation itself}
Recorded because a battery that hides its own false starts teaches the wrong lesson. First, the no-deadline rival was initially parameterised with a free offset whose fit pinned against a bound, which produced a spurious $+26$ ELPD advantage for RRHM on its own negative control; it is re-specified here as the properly nested null. Second, the first pass compared pooled fits on hierarchically generated data, so the more flexible family won regardless of which was true---the artifact that made the additive-versus-timing comparison look decisive in the wrong direction. Third, the no-deadline negative control was not in fact free of the deadline: trial-level estimate noise had been scaled by $D$, so the deadline still modulated the spread of the margin even where it did not shift its mean, and Gate~K5 accordingly failed to fire on it in 62\% of runs. With the scale made deadline-independent the same gate fires in 80\%. The first two were found by inspecting fitted parameters against their bounds; the third by asking what information the negative control still carried. Both checks are now part of the preregistered analysis plan.

\subsection{What this section does not establish}
Nothing here is evidence about phobia. Every number is a property of a design evaluated against a toy generative model, and a design that passes its own gate can still be wrong about the organism. What the exercise does establish is narrower and, for a hypothesis paper, more useful: two of this manuscript's central quantities were not identifiable in the design that v5 proposed, one kill criterion condemned the theory unconditionally, and two others could not have fired at the sample sizes named. Those are the kinds of failure that a study discovers after 120 participants and a year, or that a simulation discovers in an afternoon.

\section{Three unsolved attacks, and a narrowed scope}
\label{sec:attacks}
Five hard problems were deliberately drawn from the world literature \emph{outside} the four challenge patterns of \S\ref{sec:layers}, and run against this manuscript under one discipline: no new equation may be added; each problem must either be derived from the existing structure or recorded as unsolved. Two were substantially handled; three were not, and they are recorded here as standing attacks --- because a battery whose failures are hidden is worth less than one that reports them. All external holdings below are abstract-level and PubMed-verified unless noted.

\subsection{Attack 1: the reconsolidation bypass problem}
A single $\sim$2-minute tarantula exposure followed by 40~mg propranolol produced a near-binary switch from avoidance to approach, lasting at least a year --- with self-reported fear \emph{not} dropping immediately: behavior changed first and reported fear followed months later \cite{soeter2015}. Two honesty notes. First, that ordering matches this manuscript's governability-before-fear commitment (C19), but as a \emph{postdiction}: the 2015 result existed before C19 was written, and no prediction credit is claimed. Second, and against this manuscript: a massive, durable policy change occurred \emph{without any accumulated benign-completion history}, which the recalibration law (\S\ref{sec:recal}) says should carry the dose. Later preregistered work sharpens rather than dissolves the problem: at $N{=}69$, propranolol showed no short-term advantage over placebo, and a long-term advantage \emph{only} when the phobic cue was re-encountered within days \cite{elsey2025}; substantial placebo and non-specific effects have also been reported in reconsolidation-based treatment \cite{steinfurth2021}. The re-encounter dependency is directionally consistent with the access route (some new record is still required), but RRHM cannot currently derive \emph{why molecular destabilization of a retained record substitutes for record accumulation}, and it is forbidden here to absorb the result by declaring that propranolol ``recalibrates the horizon'' post hoc. \textbf{Status: unsolved.}

\subsection{Attack 2: the disgust problem, now with its evidence}
The limitation of \S\ref{sec:limits} is promoted to a standing attack. Spider-phobic participants show elevated disgust sensitivity \cite{merckelbach1996}; spider imagery elicits both fear and disgust with contamination ideation, measurable in facial EMG \cite{devries2002}; and state disgust relates to avoidant behavior in clinical and nonclinical samples \cite{woody2002}. The structural problem is exact: a person can hold a \emph{long} horizon estimate (``nothing will happen if I stay'') and still terminate engagement out of revulsion --- there is no deadline anywhere in that decision. The layered architecture licenses a separate disgust lane feeding the policy selector, but writing that lane in without an independently measured, preregistered, Fisher-counted term would be exactly the silent absorption this manuscript forbids. \textbf{Status: unsolved; the kill consequence for K6/K9 stated in \S\ref{sec:limits} stands.}

\subsection{Attack 3: the emetophobia recursive loop}
In specific phobia of vomiting, the eliciting cue need not be external at all: anxiety produces gastrointestinal distress, which is monitored, catastrophically appraised, and thereby amplified \cite{boschen2007}; nausea is frequent in the population and marks a more severe, longer-lasting phenotype with lower BMI \cite{vandenbergh2013}; and the treatment evidence base remains thin (37 studies, 68\% single-case, two RCTs) \cite{keyes2026}. The architecture reads this cleanly as a \emph{self-generated informational loop}: the state update is recursive, so a readout can alter the body state that becomes its own next input --- no external stimulus required. The horizon reading (visceral deviation as $d$; the estimate ``how long before correction can no longer prevent vomiting'') fits the visceral row of Table~\ref{tab:cross}, but disgust, contamination, and the social/second-reader component are all live here and uncovered. The decisive test is stated now: if the horizon estimate adds nothing to prediction of defensive transition beyond nausea, disgust, and vomiting expectancy, RRHM loses in emetophobia. If it does add prediction, the case may force a scope amendment from \emph{cue-locked} to \emph{body-state-locked} bias --- a transdiagnostic pressure that is flagged and deliberately not enacted here. \textbf{Status: structurally strong, mechanistically incomplete.}

\subsection{Two supporting results, absorbed without new structure}
\textbf{Readout-to-reality feedback (acrophobia).} At height, fearful individuals genuinely alter postural control --- restricted gaze and head movement, co-contraction, altered sway \cite{kugler2015,cleworth2021}; a case is on record in which behavioral/VR therapy improved avoidance while dizziness and sway persisted until vestibular physical therapy \cite{whitney2005}; VR exposure itself has RCT support \cite{donker2025}. The existing mixed-case machinery covers this, but the loop deserves its name: a readout error selects a stiffening policy that \emph{genuinely degrades} the body state, so the next cycle's readout receives real confirming evidence. The clinical order this forces is already the framework's own: repair $v,g$ (sensorimotor/vestibular capacity) before boundary sampling. \textbf{Absorbed.} \quad \textbf{Etiology $\neq$ maintenance (CO$_2$).} At 35\% CO$_2$ challenge, situational/natural-phenomena phobics over-respond relative to controls while animal phobics do not \cite{verburg1994}, and CO$_2$ hypersensitivity is established in panic \cite{vollmer2013}: a candidate \emph{constitutional} respiratory diathesis for the situational subtype only. This does not fill the diathesis gap --- it shows the gap has subtype structure, and it fixes a scope law:
\[
\boxed{\text{shared maintenance mechanism}\ \neq\ \text{shared etiology}.}
\]

\subsection{The narrowed claim}
The evidence assembled in this section does not support any universal statement. The defensible claim, stated once and replacing broader readings of earlier versions:
\begin{quote}
\textbf{RRHM may describe one convergent transition mechanism across etiologically and physiologically heterogeneous phobias --- not a universal cause of all phobic responding.} Different origins (respiratory vulnerability, vestibular instability, organic illness, learning history, disgust propensity) produce different state and readout disturbances, of which \emph{some} converge on a collapsing engagement margin and premature defensive transition. Where a subtype's transition timing is fully carried by another lane --- disgust being the named candidate --- RRHM does not apply there, and saying so is part of the model.
\end{quote}

\section{The Readout Error Propagation and Sedimentation Principle (REPL v0.1)}
\label{sec:repl}
\textbf{Status and provenance split.} This section formalizes what \S\ref{sec:preerror} stated in prose, at the author's direction, and reverses \S\ref{sec:attacks}'s no-new-structure discipline \emph{explicitly and with accounting}: Laws~1 and~3 are liftings of held results (the upstream-monotonicity and provenance principles of the companion Readout Condition \cite{lahtee2026rc}, full paper held; the state-update and lane structure of \cite{lahtee2026genesis}); Laws~2, 4, 5, 6 are new Dr-tier structure, and they introduce exactly \textbf{two new degrees of freedom} --- the lane-coupling matrix $C$ and the loop gain $G_{\rm loop}$ --- which are declared here, must be independently measured, and must be Fisher-counted before any fit that uses them. Nothing in this section is machine-checked or calibrated on data.

\subsection{Setup}
Let $Z_n=(S_n,T_n,c_n)$ be the informational state (current organism--environment state, ordered retained history, context), read out into typed lanes $\mathbf R_n=(R^{reg}_n,R^{phen}_n,R^{aut}_n,R^{act}_n,R^{verb}_n,\dots)$, feeding a policy and an action that alter the next state, whose outcome is retained:
\begin{equation}
Z_n\to\mathbf R_n\to\Pi_n\to a_n\to S_{n+1}\to T_{n+1}\to Z_{n+1}.
\label{eq:replloop}
\end{equation}
The system is a loop, not a line; every law below is a property of Eq.~\ref{eq:replloop}.

\subsection{The six laws}
\textbf{Law 1 (upstream error inheritance; lifted).} If the actual reader collapses a real distinction, $\tilde R(x)=\tilde R(x')$ for $x\neq x'$, then absent new access every downstream stage inherits the collapse: $\ker\tilde R\subseteq\ker(D_k\circ\tilde R)$. Applied to $x=\bot$ (remaining time unknown) versus $x'=0$ (none remaining), this yields \emph{phenomenological collapse}: $P(\bot)=P(0)$ --- the person cannot \emph{feel} the difference between unknown and doomed, because the distinction was deleted before the phenomenal stage. This is why insight (``unknown does not mean dangerous'') can fail while being sincerely believed, and it is the PRHC-U lesion (\S\ref{sec:prhcu}) restated as an inheritance theorem.

\textbf{Law 2 (non-uniform lane coupling; new).} An error at stage $j$ propagates lane-wise, $\Delta R^{(k)}_n\approx C_{jk}\,\epsilon_{j,n}$, with no requirement that the couplings be equal. $C_{j,act}\gg0$ with $C_{j,phen}\approx0$ gives behavior change without felt fear; the converse gives strong fear without escape. Hence \emph{same upstream lesion $\not\Rightarrow$ same phenomenology}, and two patients reporting identical fear can carry different lesions --- the architectural ground for lane-decomposed measurement (operation~2, \S\ref{sec:rgif}). $C$ is a declared new object: it earns its place only if lane-decomposed data can estimate it identifiably.

\textbf{Law 3 (phenomenology is real; its attribution can be wrong; lifted).} $\tilde P_n=\Phi(\tilde R^{phen}_n,c_n)\neq0$ is a fact about the phenomenal state --- the fear, air hunger, and trappedness are fully real. What fails is the inference $P_{\text{out of options}}\Rightarrow S_{\text{out of options}}$: experience certifies the readout, not the state read. Clinically: validate the experience, audit its provenance.

\textbf{Law 4 (readout-to-reality conversion; new).} Because $\partial\Pi/\partial\epsilon\neq0$ and $\partial S_{n+1}/\partial\Pi\neq0$ compose, $\partial S_{n+1}/\partial\epsilon\neq0$: an informational error manufactures part of the later reality it predicted. The acrophobic stiffening loop of \S\ref{sec:attacks} is the worked instance --- a misread that degrades real postural dynamics, which then supply genuine confirming evidence. The imaginary/real dichotomy is not available here.

\textbf{Law 5 (sedimentation; new).} Retention converts transient error into path dependence: if outcome typing errs, $\mathcal E_T(E)=H$, the biased tape $T^{biased}_{n+1}$ biases the next translation, so $\epsilon_{T,n}\to\epsilon_{R,n+1}$ --- \emph{error has memory}. This is the ledger lesion (Eq.~\ref{eq:lesion}) as a dynamical statement, and it is why tape audit precedes exposure in \S\ref{sec:rgif}.

\textbf{Law 6 (closed-loop error gain; new).} Around Eq.~\ref{eq:replloop}, define locally
\begin{equation}
G_{\rm loop}=\Big|\tfrac{\partial R}{\partial T}\tfrac{\partial T}{\partial S'}\tfrac{\partial S'}{\partial a}\tfrac{\partial a}{\partial\Pi}\tfrac{\partial\Pi}{\partial R}\Big|.
\label{eq:gloop}
\end{equation}
$G_{\rm loop}<1$: errors decay; $G_{\rm loop}>1$: a small initial error self-amplifies across cycles. This is a candidate \emph{formal definition of pathological persistence}, broader than phobia: no large founding trauma is required, only a nonzero seed and a loop gain above one. It is a local-linear idealization; no $G_{\rm loop}$ has ever been measured, and the quantity is meaningful only jointly with the measurement program below.

\subsection{Diagnosis by error locus}
The taxonomy of \S\ref{sec:preerror} extends to seven loci: retention error (mistyped past biases the present), access loss (phenomenological blindness --- distinct states feel identical), translation bias (miscalibrated urgency), readout-lane bias (lane dissociation), policy lock ($P(a\mid R)\to1$: readout becomes command), feedback error (action degrades the state --- Law 4), retention feedback (mis-recorded outcomes --- Law 5). The diagnostic question becomes: \emph{which distinction is lost, at which layer, propagating through which lanes, and does it sediment?} --- rather than ``what is this person afraid of?''

\subsection{Preregistered predictions (P-EP1--P-EP4)}
\textbf{P-EP1} (layer-specific intervention; registered as C21): if the founding defect is a retention error, phenomenal-fear-targeted work alone is less durable than record retyping; if the defect is a genuine physiological state, retyping alone must fail. \textbf{P-EP2}: upstream repair can improve policy before subjective fear (identical to C19's commitment via Law~2). \textbf{P-EP3}: a collapsed distinction cannot be restored by rehearsing the same readout (insight, reassurance); only new access restores it --- the told/learned dissociation of \S\ref{sec:decisive} is its laboratory form. \textbf{P-EP4} (registered as C22): where readout-to-reality feedback is claimed (acrophobic stiffening), changing policy first must measurably improve the physiological state afterwards; a phenotype in which large policy change leaves the state untouched falsifies Law~4 for that phenotype. Kill for the whole section: if lane couplings prove unmeasurable or non-identifiable in lane-decomposed data, Law~2 and with it the section's added structure must be withdrawn, not defended.

\section{Beyond phobia: four severe-case probes of the layer taxonomy}
\label{sec:probes}
\textbf{Fence.} REPL (\S\ref{sec:repl}) claims its error-locus taxonomy is broader than phobia. That claim invites a deliberately non-self-serving test: cases chosen \emph{not} because they resemble the model, but because they jointly satisfy four criteria --- severe psychiatric phenomenology, a measurable biological or circuit-level finding, an intervention under which patients genuinely improve, and published phenomenological data to compare against. Four such cases are probed below. Everything here is a \emph{structural} compatibility reading of the layer taxonomy, Dr throughout; all holdings are PubMed-verified abstracts; no identity claim is made anywhere ($\Delta v\to\Delta E$ places node $v$ in the causal ancestry of $E$; it never licenses $v=E$). The organizing inequality adds no degree of freedom --- it is the taxonomy in one line:
\[
\boxed{\,L_{\mathrm{bio}}\ \neq\ L_{\mathrm{readout}}\ \neq\ E_{\mathrm{phen}}\ \neq\ \Pi_{\mathrm{action}}\,}
\]

\begin{table*}[!t]
\centering
\caption{Four severe-case probes of the pre-phenomenal layer taxonomy. All evidential holdings are abstract-level, PubMed-verified; every layer diagnosis is Dr. No case validates the architecture; each tests whether its distinctions are \emph{needed} to state what is known without contradiction.}
\label{tab:probes}
\small
\begin{tabularx}{\textwidth}{p{0.10\textwidth}p{0.20\textwidth}p{0.20\textwidth}p{0.15\textwidth}X}
\toprule
\textbf{Case} & \textbf{Measured upstream finding} & \textbf{Phenomenology} & \textbf{Working intervention} & \textbf{Layer diagnosis and honest remainder} \\
\midrule
Anti-NMDAR encephalitis & antibody-driven NMDAR internalization; surface receptors and currents reduced, selectively and reversibly, titre-dependent \cite{hughes2010} & in 164 autoimmune-psychosis patients: hallucinations 83\%, delusions 77\%, disorganized speech 64\%, severe cognitive dysfunction 75\%, catatonia 64\% \cite{guasp2025} & immunotherapy (96\% first-line in a 702-patient cohort; 38\% improving within two weeks; $\sim$80\% good one-year functional outcome; delay predicts worse outcome) \cite{neos2026}; antipsychotics alone help aggression but poorly treat the psychosis/catatonia \cite{sarkis2021} & \emph{Upstream substrate failure}: one lesion class fans out into many lanes at once --- the strongest available natural experiment that repairing an upstream node outperforms treating the phenomenal readout (the biological-scale form of P-EP1). Remainder: nothing derives the \emph{content} of a given hallucination from receptor counts. \\
Capgras delusion & absent autonomic discrimination of familiar faces with intact orienting \cite{ellis1997}; modality-specific autonomic loss with right-temporal lesion \cite{bobes2020}; a second belief-evaluation deficit posited \cite{coltheart2022} & ``the face is right, the person is not'': visual identity preserved, felt familiarity absent & no established causal treatment & \emph{Cross-lane mismatch plus downstream evaluation failure}: the readout-level lesion matches the \emph{shape} of the experience more specifically than in any other case --- and, per the provenance principle, the ``impostor'' distinction is not carried by the visual source but added downstream. Remainder: explanation--phenomenology concordance is highest here while causal-treatment proof is lowest; not solved. \\
Severe OCD under responsive DBS & ventral-striatal electrophysiology tracks obsessional state; responsive stimulation durably improves obsessions and compulsions \cite{provenza2024}; signatures are strongly patient-specific (cross-patient models near chance) \cite{fridgeirsson2023} & intrusion, urge, felt loss of control; after DBS, patients describe restored cognitive-affective control and capacity to \emph{manage} the OCD --- agency over the obsession, not merely fewer thoughts \cite{bosanac2023} & responsive VC/VS DBS (investigational; severe refractory cases) & \emph{Readout-to-policy over-binding}: the pathology reads as $P(\Pi_{comp}\mid R_{intr})\to1$, and the intervention as restored interruptibility --- the qualitative reports match the governability construct (C20) more closely than a thought-deletion account. Remainder: no universal lesion; only that some circuit states track and are manipulable. \\
Catatonia & no single lesion; a multi-cause motor-regulatory syndrome & in 1,456 episodes mutism, immobility/stupor, and withdrawal dominate; among the 68 patients with phenomenological accounts, only $\sim$35\% express fear while 72\% give a narrative account \cite{dawkins2022}; painful inner experience persists and is predicted by immobility \cite{subjcat2026,northoff1996} & lorazepam (catatonic signs resolved in 16/21 completed trials; non-responders responded rapidly to ECT) \cite{ungvari1996}; single-dose response measurable in emergency settings \cite{lorchallenge2021} & \emph{Action-channel lock with phenomenology partly preserved}: felt experience continues while speech and movement channels are locked --- the strongest structural evidence in this table that what the person feels $\neq$ what the body can currently execute, and a direct counterexample to any single-scalar (fear-driven) reading. Remainder: why benzodiazepine tone-shifting unlocks the channel is not derived. \\
\bottomrule
\end{tabularx}
\end{table*}

\subsection{The many-to-many law}
Read jointly, the four probes force a generalization of \S\ref{sec:attacks}'s scope law to the psychiatric scale: \textbf{same phenomenology $\not\Rightarrow$ same lesion} (immobility can be policy lock, fear, delusion, or physiological failure; psychosis can be autoimmune-synaptic or primary-psychiatric with near-identical surface), and \textbf{same lesion class $\not\Rightarrow$ same phenomenology} (one antibody produces psychosis, catatonia, amnesia, dyskinesia, and autonomic failure at once). The mapping between lesion layers and phenomenal outcomes is many-to-many, which is precisely what a lane architecture with non-uniform couplings (Law 2) permits and what any one-symptom--one-lesion scheme forbids. The corresponding diagnostic discipline is the one already stated in \S\ref{sec:repl}: locate the failing distinction and its layer before naming the disorder by its feeling.

\subsection{The next falsifiable step, registered}
The probes suggest the program's next empirical object: a phenomenology--lesion--repair matrix across psychiatric syndromes, whose load-bearing prediction is registered now (case C24, S7~v2, before any recovery-order literature is searched): \emph{if layer diagnoses are real, the order in which readout lanes recover under a working intervention is predictable in advance from the layer of the lesion} --- upstream substrate repair should restore lanes in dependency order; policy-unlock interventions should restore action before phenomenal content shifts; cross-lane mismatch repair, if ever achieved, should restore felt familiarity before belief revision. If recovery order proves idiosyncratic and unpredictable from layer, the taxonomy is descriptive decoration and should be said to be so. Building that matrix is future work and belongs in a companion paper, not here.

\section{Discussion}
RRHM changes the proposed unit of explanation from fear intensity to the estimated time remaining before effective engagement-preserving regulation ceases to be available. Its novelty claim in this edition is deliberately narrower than in v4 and is stated once:
\[
\boxed{
\begin{aligned}
&\text{The engagement horizon is a distinct readout}\\[-0.1em]
&\text{channel with its own gain and bias, and specific}\\[-0.1em]
&\text{phobia is a persistent cue-locked bias of that}\\[-0.1em]
&\text{channel, separately identifiable from the escape,}\\[-0.1em]
&\text{imminence, control and generic-timing channels.}
\end{aligned}}
\]
The three dissociations of earlier editions are now three ways of separating this channel from a neighbor: escape availability versus engagement recoverability ($\eRRH$ versus $\dRRH$), physical threat imminence versus regulatory imminence, and perceived control versus causally effective engagement preservation. The third is kinship rather than novelty---the actual-versus-perceived-control distinction is the axis of the controllability literature \cite{maier2016}---and it is retained only as a separability requirement between two biases, not as a new distinction. A narrow claim that is fully defensible is stronger than a broad claim with one weak leg.

Three things the anchoring buys, stated plainly. First, a nonzero calibration error is the healthy default, so the pathology is a cue-conditioned bias and the model predicts residual miscalibration even after successful treatment. Second, telling someone they are safe and letting them sample the boundary are different operations, one of which cannot durably move the horizon; this is a prediction, and it is where the Open-Door observation stops being a paradox. Third, whether the construct is identifiable is a fact about the battery, which is why Module~F exists and why the design's Fisher information is checked before recruitment rather than after.

A fourth thing the anchoring bought, unplanned. Because the constructs were written as channels with gains and biases, the whole battery could be simulated before recruitment---and doing so found that two of the model's own quantities were unidentifiable in the design that was about to be run, and that one kill criterion condemned the theory whether or not it was true (\S\ref{sec:sim}). The correction was cheap at this stage and would have been expensive after data collection. A hypothesis paper that specifies numeric gates should be expected to show that it ran them.

The model remains vulnerable in precisely the way a medical theory should be vulnerable. If eRRH is a new label for control, imminence, expected value, or clock distortion, it should disappear under head-to-head modeling with preregistered priors. If it cannot be measured independently, the central construct fails. If it does not predict onset, remission, and relapse, the disease-course claim fails. If cross-phobia generalization requires incompatible definitions, the unification claim fails. If the engagement horizon is only the FID switching time under another name, \S\ref{sec:fid}'s $\Pi^E\neq\varnothing$ test will show it.

\section{Limitations}
\label{sec:limits}
RRHM currently lacks direct empirical validation. eRRH is a proposed latent variable, not a biomarker. The viability boundary $\Vset$ must be operationalized separately for each experimental paradigm. The correct transition function is unknown; Eq.~\ref{eq:hazard} is a preregistered convenience, not a derived form. ADM is a linear toy whose parameters are piecewise constant and whose physiological reading is a declared bridge; the closed-form horizon of Eq.~\ref{eq:closedform} is a property of that toy, not of any organism. Active inference may subsume eRRH mathematically. Existing exposure, safety-behavior, nonconscious-learning, epidemiological, and neuroimaging findings are compatible with portions of the model but are not specific evidence for RRHM.

The probe likelihood used in \S\ref{sec:k11} treats a horizon report as conditionally independent of trial survival, which is why the raw probe slope is attenuated; representing that selection explicitly is a known and unperformed refinement, and until it is done no clock-gain estimate should be reported from these data. The simulation of \S\ref{sec:sim} inherits every limitation of ADM and adds its own: synthetic cohorts calibrated to published rates are not data, the negative controls are only as adversarial as the rival models written for them, and three defects in the simulation itself were found and corrected during the work. Its conclusions are about what a design can identify, never about what is true of phobia.

A third gap is the widest one this edition leaves open: \textbf{RRHM has no account of disgust.} Disgust is a primary, well-documented driver of responding in exactly two of the subtypes the cross-phobia table claims to cover---spider phobia and blood-injection-injury phobia---and disgust-motivated avoidance is about contamination and revulsion, not about the time remaining before corrective action stops working. Nothing in the horizon vocabulary translates it: a disgusting stimulus can be perfectly recoverable on every regulatory axis and still be avoided absolutely. The cross-phobia table names disgust in the rival column of both rows, but naming a rival is not answering it. The commitment is stated plainly: if disgust-based measures fully account for defensive-transition timing in the animal and BII subtypes, the shared-parameter claims of Gates K6 and K9 fail in those subtypes, and the cross-phobia unification claim dies with them regardless of how the situational subtypes behave. Any future edition that quietly absorbs disgust as ``another drive term $T$'' without an independent measurement of it should be read as the post hoc absorption this paper promises not to perform.

Two further gaps are named rather than implied to be closed. First, RRHM has no diathesis: it explains why a bias persists once present, not why it arises in one person and not another, and the second problem raised in the Background is therefore an open question in this edition rather than a solved one. Second, the clinical bridge of \S\ref{sec:course} carries the non-circularity guarantee only by declaration; until physiologically estimated horizons are shown to track programmed horizons in the same participants, every disease-course claim rests on an untested transfer.

The construct may prove transdiagnostic rather than specific to specific phobia. That possibility is a scientific result, not a reason to redefine the theory post hoc. Until prospective validation exists, no RRHM-derived quantity should be used for diagnosis, prognosis, treatment selection, or clinical risk stratification.

\section{Conclusion}
RRHM proposes that specific phobia may involve a cue-locked premature contraction of the estimated engagement-preserving regulatory horizon: the predicted interval during which continued engagement can still be maintained or restored by effective corrective action. This edition states that estimate as a translation of a state with an irreducible gap, supplies a generative model from which the causal horizon is computed rather than stipulated, requires the measurement battery to inform every theory-critical direction before recruitment, converts ten verbal kill criteria into eleven numeric gates each of which must be able to fail---and then runs all ten that simulation can reach. Of the first six, four failed; the design, one construct's algebraic form, and two thresholds were changed in consequence, and the repaired battery was re-run to confirm the change worked. Of the remaining four, run at v9, one worked as written, two repeated known failure classes and carry their repairs, and one---cross-phobia invariance---remains a commitment rather than a demonstrated test.

The strongest test remains deliberately unforgiving. When objective threat, physical threat imminence, uncertainty, perceived control, escape availability, action number, and action efficacy are all held constant, a manipulation that changes only the causal duration of engagement-preserving effectiveness must change defensive-transition timing. If it does not, the distinctive RRHM mechanism should not survive.

If the dissociation replicates across healthy experiments, clinically diagnosed specific-phobia groups, physiologically heterogeneous subtypes, treatment studies, and longitudinal cohorts---and if eRRH outperforms preregistered rival models, including escape-decision and interval-timing models, in held-out data---then RRHM would qualify as a candidate new explanatory variable in human defensive regulation. Until then it remains an explicit, measurable, and falsifiable medical-science hypothesis.

\section*{Declarations}
\subsection*{Ethics statement}
This manuscript presents a theoretical medical-science framework and reports no new human or animal experimentation. Every paradigm proposed here, including the blood-injection-injury, height, dental, and choking-cue paradigms, requires prospective approval by a competent research ethics committee before any human data collection. Section~\ref{sec:safety} is a set of design constraints and is not an ethics approval.

\subsection*{Patient consent}
The single-case material in \S\ref{sec:prhcu} and \S\ref{sec:rgif} is the author's own first-person experience, disclosed here by the author's deliberate, informed choice (corrected at v21 from an earlier third-person framing, at the author's direction). No third party's clinical information is reported, so no external consent is required; the CARE single-case reporting standard will still be followed in structure for any journal submission, with the author as the consenting subject. The methodological position is stated plainly: this research program is motivated by a personal experience the author wants to understand, and it converses with the world's knowledge about that experience honestly --- which is precisely why the first-person material is fenced as hypothesis-generating narrative only, carries zero evidential weight, and every claim built near it is subjected to preregistered, falsifiable, adversarially reviewed tests. First-person limitations (self-observation bias, memory reconstruction, single case) are acknowledged and are the reason the fences exist.

\subsection*{Data availability}
No new dataset was generated for this theoretical manuscript. Reproducibility status, stated exactly (corrected at v20 after independent review found the previous wording over-claimed): the K3/K6/K8/K9 gate battery, the coupling-extension simulations, the C23 open-data pipelines (including the artifact-control scripts), the recalibration-law calculators, the figure-generation scripts, the citation-verification tool, and the frozen preregistrations (v1 and v2, with scored outcomes kept in a separate file) are deposited in the public companion repository \textbf{rrhm-open-lab} (github.com/morrocwi/rrhm-open-lab; release follows this workspace's pre-publication review gate; the reviewed snapshot is git commit \texttt{db0fd9d009bbe9998fc9783131b6f3c32862596b}, which also serves as the immutable timestamp for the preregistration chain). The earlier engine and logs S1--S6 (the v6-era battery behind Table~\ref{tab:sim} and the original figures) are \emph{not} in that deposit and remain available from the author on request only; every number that depends on them is tagged accordingly in the text. The two open datasets analyzed in \S\ref{sec:firstcontact} are third-party Zenodo deposits used under their own terms. The second-round pipelines of \S\ref{sec:secondcontact} (cases C25--C28, with the frozen criteria, the two-stage structure-only mapping records, the scored outcomes including the C28.3 concession, and the run outputs) are in the same repository; the freeze commits are \texttt{6923c0a} (C25--C27 criteria), \texttt{b203bbd} (C27 mapping), \texttt{84f55a7}/\texttt{1fa0d4d} (C28), with results at \texttt{0d2acd8}; the four datasets analyzed there are third-party Zenodo/Dryad deposits used under their own terms and are re-downloaded by the scripts rather than redistributed.

\subsection*{Funding}
This work received no specific grant from any funding agency in the public, commercial, or not-for-profit sectors.

\subsection*{Competing interests}
The author is the founder of ARAYA Nikah Social Enterprise Co., Ltd., which provides marriage and family services and premarital education, and is the author of the companion projects cited as \cite{lahtee2026ru,lahtee2026rc,lahtee2026birca,lahtee2026srct}. Neither the enterprise nor those projects sells any diagnostic or therapeutic product derived from RRHM. The author declares no other competing interests.

\subsection*{Author contributions (CRediT)}
Y.L. is the sole author: \emph{Conceptualization}, \emph{Methodology}, \emph{Formal analysis}, \emph{Software} (the generative model, the gate battery, and the figures), \emph{Investigation}, \emph{Visualization}, \emph{Writing --- original draft}, \emph{Writing --- review and editing}, and \emph{Project administration}. No funding was administered.

\subsection*{Preregistration}
This is a hypothesis and design paper; no study has been run. No analysis reported here was preregistered, and none of it is confirmatory. Every threshold in Table~\ref{tab:gates} is a default proposed for preregistration, not a result. The simulations of \S\ref{sec:sim} and \S\ref{sec:k11} were run to evaluate a design before recruitment and are reported in full, including the runs that failed and the three defects corrected during the work; they carry no confirmatory status and no hypothesis was tested against data.

\subsection*{Reporting standards}
No reporting checklist (CONSORT, STROBE, PRISMA, ARRIVE) applies to a theoretical manuscript with no empirical sample. Should the paradigm of \S\ref{sec:decisive} be run, it would be reported under CONSORT for the randomised comparison, with the analysis plan, thresholds, and simulation code deposited before the first participant.

\subsection*{Supplementary material}
Listed for provenance; items S1--S6 are available from the author on request and are NOT in the public deposit (see Data availability); S7--S14 are in the public repository: (S1) \texttt{rrhm\_engine.py}, the hierarchical generative model of the decisive experiment and the discrete-time hazard likelihood; (S2) the gate battery, rounds 1--4, including the corrected negative controls; (S3) the design-level Fisher-information code behind Fig.~\ref{fig:rank}; (S4) the anchor--drive closed-form horizon and Fig.~\ref{fig:adm}; (S5) the planning power estimate of \S\ref{sec:decisive}; (S6) a run log of every figure reported in this manuscript; (S7) the compatibility-audit preregistration of \S\ref{sec:audit}, frozen verbatim before the C6--C15 literature was searched; (S8) the independent re-implementation and run logs of the K3/K6/K8/K9 battery of \S\ref{sec:simremaining}, including the repaired variants and the sample-size sweeps; (S9) the coupling-extension simulations of \S\ref{sec:coupling}: closed-form checks P6--P8, the K12 capacity gate with both generators and the effect-floor sweep, and the double-dissociation classifier with its failure at small effects; (S10) the phenomenology fence F1--F10 with per-item holding status, the C18 preregistration, and the OPEN-ANCHOR-BANDWIDTH-001 registry entry; (S11) the readout-governance intervention framework of \S\ref{sec:rgif} in its full companion statement; (S12) the tier-1 ordinal-consistency simulations; (S13) the preregistered two-source open-data test of \S\ref{sec:firstcontact}, with the frozen plan (C23), analysis code, and scored outcomes including the wrong prior guess. Software: Python 3.12 with NumPy 2.4, SciPy 1.17; random seeds are fixed in each script.

\subsection*{Version history}
v4 (adversarial edition, verbal kill criteria) $\rightarrow$ v5 (formal repairs; generative model; escape-decision and interval-timing rivals added; kill criteria made numeric) $\rightarrow$ v6 (battery simulated before recruitment; four gates failed; PRHC committed to its additive form; probes and Module~F made mandatory) $\rightarrow$ v7 (reference list verified field-by-field against PubMed and Crossref; three citation errors corrected, one of them an entry inherited from v4 that pointed at a different article) $\rightarrow$ v8 (Gate~K11 re-run with probe slopes and shown two-sided; internal counts reconciled; preprint apparatus completed) $\rightarrow$ v9 (the four remaining simulable gates K3/K6/K8/K9 run on an independent re-implementation, with K3 repaired by disattenuation, K9 bound to the 2-SE rule, and K6 honestly left conditional; the ``six of eleven simulable'' overcount of earlier versions corrected; the preregistered compatibility audit C6--C15 executed and reported with its prior-guess scorecard; the disgust gap stated as a named limitation with its kill consequence; the interval-timing rival grounded in the phobic population via Watts \& Sharrock 1984) $\rightarrow$ v10 (the Dr-tier coupling extension of \S\ref{sec:coupling} --- companion as additive damping $k$, bereavement onset as deletion of a shared term, the phobia/agoraphobia double dissociation with its simulated separability and declared effect floor, Gate K12 shown two-sided by simulation, anticipatory dread and the social-observer channel named as open terms; the withdrawn rival-comparison table and the full-text-or-NO-DATA rule recorded in S7 v2's concession ledger) $\rightarrow$ v11 (the phenomenology fence executed and deposited as S10; the companion term decomposed into co-regulation $k$ / capacity $\Delta g_{ext}$ / declared observer term on the strength of the safe-person CO$_2$ result; the portable-anchor bandwidth problem stated with its ladder, falsifiers, and quantitative fading rule P10; case C18 preregistered; the second-reader gap upgraded from declared-open to prerequisite work after full-text confirmation of its centrality) $\rightarrow$ v12 (the case-motivated refinement PRHC-U --- indeterminacy-triggered horizon contraction, Eq.~\ref{eq:prhcu} --- stated in medical-science terms with situational-geometry locking, its predictions and kill criteria, the one-cell unobservable-$D$ extension of the decisive experiment, three derived recalibration routes for the indeterminate branch, and three new translation-table rows keeping the construct invertible to the readout vocabulary ($0\neq\bot$)) $\rightarrow$ v13 (the recalibration law solved as a discrete retained-record ledger --- the floor-default derived as a typing error on censored records, the correction as three ordered ledger operations with a computable completion dose $n^{*}$, order-of-operations predictions and kill criteria, and the corresponding translation-table row; clinical translation added as censoring-aware recalibration with ordered, separately refutable components) $\rightarrow$ v14 (the layered informational architecture placed upstream of RRHM --- state update over ordered retained history, RRHM re-placed as a domain quotient, fear as one typed readout lane --- with four challenge patterns read at the correct layer: post-organic swallowing phobia without a safety baseline, BII response heterogeneity including combined fainting and PNES, cue-specific abolition of lifelong arachnophobia after mesial temporal lobectomy, and masked-exposure change without phenomenal readout; each with its honest remainder stated) $\rightarrow$ v15 (the derived intervention framework of \S\ref{sec:rgif} --- unit of treatment as the information-bearing continuation, seven ordered operations, three hard applications, and the governability-before-fear prediction with the interruptibility window, fear-lag, and transport-curve measures preregistered as C19--C20 before the desynchrony literature was searched) $\rightarrow$ v16 (under a no-new-equations discipline: three standing attacks recorded --- the reconsolidation bypass problem, the disgust problem with its evidence, the emetophobia recursive loop --- with the propranolol behavior-before-fear result marked as postdiction, not prediction credit; readout-to-reality feedback and the CO$_2$ subtype diathesis absorbed without new structure; the scope claim narrowed to one convergent transition mechanism rather than a universal cause; the pre-phenomenal error-class taxonomy added with upstream monotonicity from the Readout Condition, now held in full) $\rightarrow$ v17 (at the author's direction and reversing v16's no-new-structure discipline with explicit accounting: REPL v0.1 --- six laws of readout error propagation and sedimentation, of which two are liftings of held companion results and four are new Dr-tier structure introducing exactly two declared degrees of freedom ($C$, $G_{\rm loop}$); the seven-locus diagnostic taxonomy; predictions P-EP1--P-EP4 with P-EP1/P-EP4 registered as C21--C22; and the section-level kill criterion that non-identifiable lane couplings withdraw the added structure) $\rightarrow$ v18 (first empirical contact --- preregistered case C23 frozen before data access, two unrelated open datasets downloaded raw and analyzed end-to-end; lane dissociation confirmed emphatically ($r\le0.12$ between subjective and physiological reactivity, discordant quadrants 26--32\%), the conditioning differential found to peak at median 3.48~s against a 3.5~s programmed deadline, the amplitude criterion reported marginal without post-hoc rescue; compatibility, not validation, claimed; full typographic pass: overfull columns in the error-class table and the bandwidth ladder repaired, audit-table labels compacted with a legend, all nine tables visually verified) $\rightarrow$ v19 (four severe-case probes of the layer taxonomy chosen by non-self-serving criteria --- anti-NMDAR encephalitis, Capgras delusion, responsive-DBS OCD, catatonia --- each with a verified evidential anchor, a layer diagnosis, and an honest remainder; the many-to-many law stated at psychiatric scale; the phenomenology--lesion--repair matrix named as future work with its load-bearing recovery-order prediction registered as C24 before any recovery-order literature was searched) $\rightarrow$ v20 (this version: the corpus's first \emph{independent} adversarial review --- six parallel reviewers over mathematics, statistics/simulations, clinical accuracy, epistemic discipline, citations, and reproducibility/language --- integrated in full. Dispositions: the C23 Source-B timing result withdrawn as a window-edge argmax artifact after the review-ordered controls were run (extended-window peak follows the new edge to 5.99~s; permutation nulls pile at the edge); the ``3.5 of 4'' fractional scoring and the invented ``marginal'' label retracted, C23 rescored under its frozen binary rules; Source-A criteria annotated as unfailable-by-design; C7 rescored to NO-DATA (audit now ten of ten); four bibliography author attributions corrected (Ramirez-Bermudez, Nuara, Nho, Bush--Francis); the patient-consent declaration corrected and the single-case provenance stated; psychiatric differential added to the mimic gate; ARFID named and fenced with a pediatric-only-interpretive rule; the $n^{*}$ formula's integer-boundary off-by-one fixed (floor-plus-one); the multiplicative-PRHC parametrization clarified as an onset-gain rival distinct from Eq.~1's flowing form; Eq.~\ref{eq:companion} rewritten in dimensionally coherent dynamics form; P10's rescue inequality corrected; K3 renamed a minimum-effect test; the 2-SE rule defined; K6/K9 verdicts marked provisional (pooled-fit artifact; a post-review two-stage refit did not rescue them); default holdings declared for the legacy rival sections; the public reproduction repository cited with its reviewed commit hash; and the review's confirmed findings archived in that repository) $\rightarrow$ v21 (this version: second round of preregistered open-data contacts, C25--C28, run under the C23 lessons --- failable criteria, design-fixed windows, two-stage structure-only mapping, binary scoring, freeze-before-access enforced by public commit hashes; reported as a challenge ledger (Table~\ref{tab:ledger}) rather than a score; the principal result is C28's six-month multitrait--multimethod inequality (within-lane $+0.418$ versus same-day between-lane $+0.156$, surviving a complete-case sensitivity check and a participant bootstrap whose wide positive interval is reported as such), with the C27 unpredictability-effect lane dissociation and the out-of-domain transport of C26 alongside; the C27 temporal-gradient criterion recorded as formally passed but FRAGILE and non-citable as support; the C28.3 stability-ordering prior falsified and conceded adjacent to the primary result, withdrawing the ``rating lane carries the stable structure'' reading of C25 and leaving the lane architecture symmetric; the empirical claim upgraded from replicated dissociation to \emph{longitudinally stable lane-specific structure with weak contemporaneous cross-lane convergence}, with the explicit fences that a one-factor model with unequal loadings is not formally excluded and that the named next step is preregistered model comparison; the rankStab analysis previously recorded as unresolved ($\bot$, tooling) executed after the tooling block was removed; the named model comparison then executed as preregistered case C29 with contamination declared and a semi-blind extinction arm --- two-lane preferred at $\Delta\mathrm{BIC}=+7.83$ semi-blind and $+2.50$ contaminated, both in the frozen weak band, with the M0-true bootstrap calibration and its tail probabilities ($0.002$, $0.006$) reported alongside and one recorded prior guess wrong; an apparent inter-lane coupling sign flip corrected on audit as a factor-orientation artifact (canonical $\hat\varphi$ $+0.22$/$+0.31$, consistent); and the external transport test C30-A run on IU-Hamburg with cell-level contamination declared, the parameter-transport arm ruled not identifiable at freeze, and the falsifier TRIGGERED --- $\Delta\mathrm{BIC}=-10.03$, the single-factor model wins externally, the frozen consequence recorded verbatim as a claim boundary on typed-lane transport, and no factor allocation edited after the result).

\subsection*{Use of generative artificial intelligence}
Generative AI assistants were used for literature mapping, adversarial review of successive drafts, implementation of the simulation code in \S\ref{sec:sim} and \S\ref{sec:k11}, the independent re-implementation and execution of the K3/K6/K8/K9 battery of \S\ref{sec:simremaining}, the execution of the preregistered compatibility-audit searches of \S\ref{sec:audit}, LaTeX preparation, and editorial restructuring. They were not used to generate data, to select results, or to decide what the manuscript concludes. Every formal claim, equation, control, threshold, and cited source was checked by the author; the reference list was additionally verified field-by-field against PubMed and Crossref, which found and corrected three errors (see Version history). The author takes full responsibility for all content. No AI system is listed as an author, in line with current authorship guidance.

\begin{thebibliography}{99}
\footnotesize
\setlength{\itemsep}{0.10em}
\setlength{\parskip}{0pt}

\bibitem{wardenaar2017}
Wardenaar KJ, Lim CCW, Al-Hamzawi AO, et al. The cross-national epidemiology of specific phobia in the World Mental Health Surveys. \emph{Psychol Med}. 2017;47(10):1744--1760. PMID: 28222820. doi:10.1017/S0033291717000174.

\bibitem{wardenaar2018corr}
Wardenaar KJ, Lim CCW, Al-Hamzawi AO, et al. The cross-national epidemiology of specific phobia in the World Mental Health Surveys --- CORRIGENDUM. \emph{Psychol Med}. 2018;48(5):878. PMID: 28994357. doi:10.1017/S0033291717002975.

\bibitem{odgers2022}
Odgers K, Kershaw KA, Li SH, Graham BM. The relative efficacy and efficiency of single- and multi-session exposure therapies for specific phobia: A meta-analysis. \emph{Behav Res Ther}. 2022;159:104203. PMID: 36323055. doi:10.1016/j.brat.2022.104203.

\bibitem{trumpf2010incidence}
Trumpf J, Becker ES, Vriends N, Meyer AH, Margraf J. Incidence and predictors of specific phobia in young women: a prospective community study. \emph{J Anxiety Disord}. 2010;24(1):87--93. PMID: 19815371. doi:10.1016/j.janxdis.2009.09.002.

\bibitem{trumpf2009remission}
Trumpf J, Becker ES, Vriends N, Meyer AH, Margraf J. Rates and predictors of remission in young women with specific phobia: a prospective community study. \emph{J Anxiety Disord}. 2009;23(7):958--964. PMID: 19604666. doi:10.1016/j.janxdis.2009.06.005.

\bibitem{bohnlein2020}
B\"ohnlein J, Altegoer L, Muck NK, et al. Factors influencing the success of exposure therapy for specific phobia: A systematic review. \emph{Neurosci Biobehav Rev}. 2020;108:796--820. PMID: 31830494.

\bibitem{devos2025}
de Vos JH, Lange I, Goossens L, et al. Long-term exposure therapy outcome in phobia and the link with behavioral and neural indices of extinction learning. \emph{J Affect Disord}. 2025;375:324--330. PMID: 39889926. doi:10.1016/j.jad.2025.01.133.

\bibitem{krypotos2015}
Krypotos AM, Effting M, Kindt M, Beckers T. Avoidance learning: a review of theoretical models and recent developments. \emph{Front Behav Neurosci}. 2015;9:189. PMID: 26257618.

\bibitem{pittig2018}
Pittig A, Treanor M, LeBeau RT, Craske MG. The role of associative fear and avoidance learning in anxiety disorders: gaps and directions for future research. \emph{Neurosci Biobehav Rev}. 2018;88:117--140. PMID: 29550209.

\bibitem{craske2014}
Craske MG, Treanor M, Conway CC, Zbozinek T, Vervliet B. Maximizing exposure therapy: an inhibitory learning approach. \emph{Behav Res Ther}. 2014;58:10--23. PMID: 24864005.

\bibitem{kausche2025}
Kausche FM, Carsten HP, Sobania KM, Riesel A. Fear and safety learning in anxiety- and stress-related disorders: an updated meta-analysis. \emph{Neurosci Biobehav Rev}. 2025;169:105983. PMID: 39706234. doi:10.1016/j.neubiorev.2024.105983. Corrigendum: \emph{Neurosci Biobehav Rev}. 2025;172:106123.

\bibitem{mobbs2007}
Mobbs D, Petrovic P, Marchant JL, et al. When fear is near: threat imminence elicits prefrontal-periaqueductal gray shifts in humans. \emph{Science}. 2007;317:1079--1083. PMID: 17717184.

\bibitem{mobbs2020}
Mobbs D, Headley DB, Ding W, Dayan P. Space, time, and fear: survival computations along defensive circuits. \emph{Trends Cogn Sci}. 2020;24(3):228--241. doi:10.1016/j.tics.2019.12.016.

\bibitem{qi2018}
Qi S, Hassabis D, Sun J, Guo F, Daw N, Mobbs D. How cognitive and reactive fear circuits optimize escape decisions in humans. \emph{Proc Natl Acad Sci USA}. 2018;115(12):3186--3191. doi:10.1073/pnas.1712314115.

\bibitem{zhang2025escape}
Zhang H, Cheng J, Hu K, Wang F, Qi S, Liu Q, Yao Y, Mobbs D, Wu H. An intracranial dissection of human escape circuits. \emph{Nat Commun}. 2025;16:5520. doi:10.1038/s41467-025-60666-9.

\bibitem{fung2019}
Fung BJ, Qi S, Hassabis D, Daw N, Mobbs D. Slow escape decisions are swayed by trait anxiety. \emph{Nat Hum Behav}. 2019;3(7):702--708. PMID: 31110337. doi:10.1038/s41562-019-0595-5.

\bibitem{maier2016}
Maier SF, Seligman MEP. Learned helplessness at fifty: insights from neuroscience. \emph{Psychol Rev}. 2016;123(4):349--367. PMID: 27337390. doi:10.1037/rev0000033.

\bibitem{hartley2014}
Hartley CA, Gorun A, Reddan MC, Ramirez F, Phelps EA. Stressor controllability modulates fear extinction in humans. \emph{Neurobiol Learn Mem}. 2014;113:149--156. PMID: 24333646.

\bibitem{wood2015}
Wood KH, Wheelock MD, Shumen JR, Bowen KH, Ver Hoef LW, Knight DC. Controllability modulates the neural response to predictable but not unpredictable threat in humans. \emph{NeuroImage}. 2015;119:371--381. PMID: 26149610.

\bibitem{droitvolet2007}
Droit-Volet S, Meck WH. How emotions colour our perception of time. \emph{Trends Cogn Sci}. 2007;11(12):504--513. PMID: 18023604.

\bibitem{stetson2007}
Stetson C, Fiesta MP, Eagleman DM. Does time really slow down during a frightening event? \emph{PLoS One}. 2007;2(12):e1295. PMID: 18074019.

\bibitem{binks2015}
Binks S, Chan D, Medford N. Abolition of lifelong specific phobia: a novel therapeutic consequence of left mesial temporal lobectomy. \emph{Neurocase}. 2015;21(1):79--84. PMID: 24460482. doi:10.1080/13554794.2013.873056. (Abstract holding only.)

\bibitem{vonmucke2022}
von M\"ucke-Heim IA, Walter I, Nischwitz S, Erhardt A. Combined fainting and psychogenic non-epileptic seizures as significant therapy hurdles in blood-injury-injection phobia: a case report and mini-review. \emph{Front Psychiatry}. 2022;13:915058. PMID: 35903630. doi:10.3389/fpsyt.2022.915058. (Abstract holding only.)

\bibitem{staiano2003}
Staiano A. Food refusal in toddlers with chronic diseases. \emph{J Pediatr Gastroenterol Nutr}. 2003;37(3):225--227. PMID: 12960640. doi:10.1097/00005176-200309000-00004. (Abstract holding only.)

\bibitem{finn2026}
Finn DM, Giambrone CN, Ciobanu CS, et al. Lifelong solid food refusal in an 8-year-old girl with a choking phobia: a case report. \emph{J Eat Disord}. 2026. PMID: 42243909. doi:10.1186/s40337-026-01673-z. (Abstract holding only.)

\bibitem{tolchin2026}
Tolchin B, Goldstein LH, Reuber M, et al. Management of Functional Seizures Practice Guideline Executive Summary: Report of the AAN Guidelines Subcommittee. \emph{Neurology}. 2026. PMID: 41370742. doi:10.1212/WNL.0000000000214466. (Abstract holding only.)

\bibitem{soeter2015}
Soeter M, Kindt M. An abrupt transformation of phobic behavior after a post-retrieval amnesic agent. \emph{Biol Psychiatry}. 2015. PMID: 25980916. (Abstract holding only.)

\bibitem{elsey2025}
(Authors per PubMed record.) Re-encountering the phobic cue within days after a reconsolidation intervention is crucial to observe a lasting fear reduction in spider phobia. \emph{Mol Psychiatry}. 2025. PMID: 39789312. (Abstract holding only.)

\bibitem{steinfurth2021}
(Authors per PubMed record.) Placebo and non-specific effects in reconsolidation-based treatment for arachnophobia. \emph{Front Psychiatry}. 2021. PMID: 34867558. (Abstract holding only.)

\bibitem{boschen2007}
Boschen MJ. Reconceptualizing emetophobia: a cognitive-behavioral formulation and research agenda. \emph{J Anxiety Disord}. 2007. PMID: 16890398. (Abstract holding only.)

\bibitem{vandenbergh2013}
(Authors per PubMed record.) Nausea in specific phobia of vomiting. \emph{Behav Sci (Basel)}. 2013. PMID: 25379248. (Abstract holding only.)

\bibitem{keyes2026}
(Authors per PubMed record.) Treatment interventions for emetophobia: an extensive scoping review. \emph{J Psychiatr Res}. 2026. PMID: 41905115. (Abstract holding only.)

\bibitem{kugler2015}
(Authors per PubMed record.) Acrophobia impairs visual exploration and balance during standing and walking. \emph{Ann N Y Acad Sci}. 2015. PMID: 25722015. (Abstract holding only.)

\bibitem{cleworth2021}
(Authors per PubMed record.) The effects of virtual height exposure on postural control and psychophysiological stress are moderated by individual height intolerance. \emph{Front Hum Neurosci}. 2021. PMID: 35095450. (Abstract holding only.)

\bibitem{whitney2005}
Whitney SL, et al. Acrophobia and pathological height vertigo: indications for vestibular physical therapy? \emph{Phys Ther}. 2005. PMID: 15842192. (Abstract holding only.)

\bibitem{donker2025}
(Authors per PubMed record.) Efficacy of exposure scenario in virtual reality for the treatment of acrophobia: a randomized controlled trial. \emph{J Behav Ther Exp Psychiatry}. 2025. PMID: 40186999. (Abstract holding only.)

\bibitem{merckelbach1996}
(Authors per PubMed record.) Disgust and spider phobia. \emph{J Abnorm Psychol}. 1996. PMID: 8772018. (Abstract holding only.)

\bibitem{devries2002}
(Authors per PubMed record.) Disgust and disgust sensitivity in spider phobia: facial EMG in response to spider and oral disgust imagery. \emph{J Anxiety Disord}. 2002. PMID: 12396207. (Abstract holding only.)

\bibitem{woody2002}
(Authors per PubMed record.) The relationship between disgust sensitivity and avoidant behavior: studies of clinical and nonclinical samples. \emph{J Anxiety Disord}. 2002. PMID: 12396211. (Abstract holding only.)

\bibitem{verburg1994}
(Authors per PubMed record.) A 35\% carbon dioxide challenge in simple phobias. \emph{Acta Psychiatr Scand}. 1994. PMID: 7892773. (Abstract holding only.)

\bibitem{vollmer2013}
(Authors per PubMed record.) The carbon dioxide challenge test in panic disorder: a systematic review of preclinical and clinical research. \emph{Braz J Psychiatry}. 2013. PMID: 24142095. (Abstract holding only.)

\bibitem{guasp2025}
Ramirez-Bermudez J, Espinola-Nadurille M, et al. Autoimmune psychosis: psychopathological patterns and outcome after immunotherapy. \emph{Schizophr Res}. 2025. PMID: 40306141. (Abstract holding only.)

\bibitem{hughes2010}
Hughes EG, et al. Cellular and synaptic mechanisms of anti-NMDA receptor encephalitis. \emph{J Neurosci}. 2010. PMID: 20427647. (Abstract holding only.)

\bibitem{neos2026}
(Authors per PubMed record.) Development and validation of the NEOS2 score for prediction of long-term outcomes and improvement after first-line immunotherapy in patients with anti-NMDAR encephalitis: an international cohort study. \emph{Lancet Reg Health Eur}. 2026. PMID: 41488792. (Abstract holding only.)

\bibitem{sarkis2021}
(Authors per PubMed record.) Psychiatric management of anti-NMDAR encephalitis: a cohort analysis. \emph{Psychol Med}. 2021. PMID: 31739809. (Abstract holding only.)

\bibitem{ellis1997}
Ellis HD, Young AW, et al. Reduced autonomic responses to faces in Capgras delusion. \emph{Proc Biol Sci}. 1997. PMID: 9263474. (Abstract holding only.)

\bibitem{bobes2020}
Nuara A, et al. Catching the imposter in the brain: the case of Capgras delusion. \emph{Cortex}. 2020. PMID: 32540160. (Abstract holding only.)

\bibitem{coltheart2022}
Coltheart M, Davies M. What is Capgras delusion? \emph{Cogn Neuropsychiatry}. 2022. PMID: 34890309. (Abstract holding only.)

\bibitem{provenza2024}
Nho YH, Rolle CE, et al. Responsive deep brain stimulation guided by ventral striatal electrophysiology of obsession durably ameliorates compulsion. \emph{Neuron}. 2024. PMID: 37865084. (Abstract holding only.)

\bibitem{bosanac2023}
(Authors per PubMed record.) Phenomenological changes associated with deep brain stimulation for obsessive-compulsive disorder: a cognitive appraisal model of recovery. \emph{Brain Sci}. 2023. PMID: 37891812. (Abstract holding only.)

\bibitem{fridgeirsson2023}
(Authors per PubMed record.) Patient-specific intracranial neural signatures of obsessions and compulsions in the ventral striatum. \emph{J Neural Eng}. 2023. PMID: 36827705. (Abstract holding only.)

\bibitem{dawkins2022}
(Authors per PubMed record.) Catatonia psychopathology and phenomenology in a large dataset. \emph{Front Psychiatry}. 2022. PMID: 35677876. (Abstract holding only.)

\bibitem{northoff1996}
(Authors per PubMed record.) [The subjective experience in catatonia: systematic study of 24 catatonic patients]. \emph{Psychiatr Prax}. 1996. PMID: 8657812. (Abstract holding only.)

\bibitem{subjcat2026}
(Authors per PubMed record.) Subjective experiences of catatonia in acute psychiatric inpatients: an exploratory study. \emph{Psychiatry Res}. 2026. PMID: 42617301. (Abstract holding only.)

\bibitem{ungvari1996}
Bush G, Fink M, Petrides G, Dowling F, Francis A. Catatonia. II. Treatment with lorazepam and electroconvulsive therapy. \emph{Acta Psychiatr Scand}. 1996. PMID: 8686484. (Abstract holding only.)

\bibitem{lorchallenge2021}
(Authors per PubMed record.) Revisiting lorazepam challenge test: clinical response with dose variations and utility for catatonia in a psychiatric emergency setting. \emph{Aust N Z J Psychiatry}. 2021. PMID: 33124447. (Abstract holding only.)

\bibitem{vrspider2026}
Dataset for ``Do Virtual Spiders Feel Real? Dissociation Between Subjective and Physiological Fear During Virtual Reality Exposure.'' Zenodo record 20796399; 2026. \url{https://zenodo.org/records/20796399}. (Raw dataset downloaded and analyzed 2026-09-02; summary-level variables only; per-participant streams not used.)

\bibitem{hra1data}
PsPM-HRA1: Skin conductance responses in fear conditioning with visual CS and electrical US. Zenodo record 321641; PsPM team. \url{https://zenodo.org/records/321641}. Described in Bach et al.\ 2010 (doi:10.1016/j.biopsycho.2010.06.007) and Staib et al.\ 2015 (doi:10.1016/j.jneumeth.2015.08.009), per the dataset README. (Raw dataset downloaded and analyzed 2026-09-02.)

\bibitem{iuhamburg2021}
Klingelh\"ofer-Jens M, Morriss J, Lonsdorf TB. Effects of intolerance of uncertainty on subjective and psychophysiological measures during threat acquisition and delayed threat extinction. Zenodo record 5648055; 2021. doi:10.5281/zenodo.5648055. (Creators per the Zenodo deposit metadata and dataset README; raw CSVs downloaded and analyzed 2026-09-02 under preregistered case C25.)

\bibitem{giles2015}
Giles GE, Mahoney CR, Bruny\'e TT, Taylor HA, Kanarek RB. Data from: Stress effects on mood, HPA axis, and autonomic response: comparison of three psychosocial stress paradigms. Dryad/Zenodo deposit, record 4973029; doi:10.5061/dryad.64313. (Creators per the deposit metadata; single spreadsheet downloaded and analyzed 2026-09-02 under preregistered case C26.)

\bibitem{koppold2026}
Koppold A. How temporal predictability of threat and action preparation affect defensive freezing responses. Zenodo record 18713991; 2026. doi:10.5281/zenodo.18713991. (Creator per the Zenodo deposit metadata; the deposit's methods file names further study authors. Raw archive downloaded and analyzed 2026-09-02 under preregistered case C27.)

\bibitem{rankstab2022}
Klingelh\"ofer-Jens M, Ehlers MR, Kuhn M, Keyaniyan V, Lonsdorf TB. Robust group- but limited individual-level (longitudinal) reliability and insights into cross-phases response prediction of conditioned fear. Zenodo record 7323547; 2022. doi:10.5281/zenodo.7323547. Published article: doi:10.7554/eLife.78717 (title and author set verified against Crossref). (Authors per the deposit's own codebook; raw data downloaded and analyzed 2026-09-02 under preregistered case C28.)

\bibitem{lahtee2026genesis}
Lahtee Y. \emph{Readout Genesis: Information-Epistemic Foundation Journal}. Companion document, Open Civil Science Initiative; 2026. (Held on file; source of the layered architecture of \S\ref{sec:layers}.)

\bibitem{carter1995}
Carter MM, Hollon SD, Carson R, Shelton RC. Effects of a safe person on induced distress following a biological challenge in panic disorder with agoraphobia. \emph{J Abnorm Psychol}. 1995;104(1):156--163. PMID: 7897039. (Abstract holding only; cited for the abstract-stated result that safe-person presence attenuated distress, catastrophic cognitions, and physiological arousal under CO$_2$.)

\bibitem{coan2006}
Coan JA, Schaefer HS, Davidson RJ. Lending a hand: social regulation of the neural response to threat. \emph{Psychol Sci}. 2006;17(12):1032--1039. PMID: 17201784. (Abstract holding only; cited for the spouse-versus-stranger gradient and its moderation by marital quality.)

\bibitem{baker2022}
Baker HJ, Hollywood A, Waite P. Adolescents' lived experience of panic disorder: an interpretative phenomenological analysis. \emph{BMC Psychol}. 2022;10:143. PMID: 35668509. PMC9167912. doi:10.1186/s40359-022-00849-x. (Full text read; cited for the two superordinate themes, including ``an unacceptable self''.)

\bibitem{watts1984}
Watts FN, Sharrock R. Fear and time estimation. \emph{Percept Mot Skills}. 1984;59(2):597--598. PMID: 6514503. doi:10.2466/pms.1984.59.2.597. (Metadata and abstract verified against PubMed; cited only for the abstract-stated result that phobic participants overestimated a feared-stimulus interval.)

\bibitem{stephan2016}
Stephan KE, Manjaly ZM, Mathys CD, et al. Allostatic self-efficacy: a metacognitive theory of dyshomeostasis-induced fatigue and depression. \emph{Front Hum Neurosci}. 2016;10:550. PMID: 27895566.

\bibitem{paulus2019}
Paulus MP, Feinstein JS, Khalsa SS. An active inference approach to interoceptive psychopathology. \emph{Annu Rev Clin Psychol}. 2019;15:97--122. PMID: 31067416.

\bibitem{ledoux2016}
LeDoux JE, Pine DS. Using neuroscience to help understand fear and anxiety: a two-system framework. \emph{Am J Psychiatry}. 2016;173(11):1083--1093. PMID: 27609244.

\bibitem{klein1993}
Klein DF. False suffocation alarms, spontaneous panics, and related conditions: an integrative hypothesis. \emph{Arch Gen Psychiatry}. 1993;50(4):306--317. PMID: 8466392.

\bibitem{ayala2009}
Ayala ES, Meuret AE, Ritz T. Treatments for blood-injury-injection phobia: a critical review of current evidence. \emph{J Psychiatr Res}. 2009;43:1235--1242. PMID: 19464700.

\bibitem{ost1987}
\"Ost LG, Sterner U. Applied tension: a specific behavioral method for treatment of blood phobia. \emph{Behav Res Ther}. 1987;25(1):25--29. PMID: 3593159.

\bibitem{ritz2010}
Ritz T, Meuret AE, Ayala ES. The psychophysiology of blood-injection-injury phobia: looking beyond the diphasic response paradigm. \emph{Int J Psychophysiol}. 2010;78(1):50--67. PMID: 20576505.

\bibitem{huppert2020}
Huppert D, Wuehr M, Brandt T. Acrophobia and visual height intolerance: advances in epidemiology and mechanisms. \emph{J Neurol}. 2020;267(Suppl 1):231--240. PMID: 32444982.

\bibitem{siegel2025}
Siegel P, Peterson BS. ``All we have to fear is fear itself'': paradigms for reducing fear by preventing awareness of it. \emph{Psychol Bull}. 2024;150(9):1118--1154. PMID: 39088015. doi:10.1037/bul0000437.

\bibitem{carl2019}
Carl E, Stein AT, Levihn-Coon A, et al. Virtual reality exposure therapy for anxiety and related disorders: a meta-analysis of randomized controlled trials. \emph{J Anxiety Disord}. 2019;61:27--36. PMID: 30287083.

\bibitem{meckes2025}
Meckes S, Cole A, Gee K, Lancaster C. Testing judicious use of safety behaviors during exposure: a randomized controlled trial examining when and for whom safety behaviors improve fear reduction. \emph{Cogn Ther Res}. 2026;50(2):341--359. doi:10.1007/s10608-025-10667-1.

\bibitem{bennettlevy2004}
Bennett-Levy J, Butler G, Fennell M, Hackmann A, Mueller M, Westbrook D, eds. \emph{Oxford Guide to Behavioural Experiments in Cognitive Therapy}. Oxford: Oxford University Press; 2004.

\bibitem{vehtari2017}
Vehtari A, Gelman A, Gabry J. Practical Bayesian model evaluation using leave-one-out cross-validation and WAIC. \emph{Stat Comput}. 2017;27:1413--1432. doi:10.1007/s11222-016-9696-4.

\bibitem{lakens2017}
Lakens D. Equivalence tests: a practical primer for t tests, correlations, and meta-analyses. \emph{Soc Psychol Personal Sci}. 2017;8(4):355--362. doi:10.1177/1948550617697177.

\bibitem{cohen1988}
Cohen J. \emph{Statistical Power Analysis for the Behavioral Sciences}. 2nd ed. Hillsdale, NJ: Lawrence Erlbaum; 1988.

\bibitem{lor2025}
Lor CS, Steyrl D, Karner A, et al. SpiderPhy dataset: a multimodal dataset of physiological, psychometric and behavioral responses to fear stimuli. \emph{Sci Data}. 2025;12(1):599. PMID: 40210881. doi:10.1038/s41597-025-04908-x.

\bibitem{zhang2025}
Zhang M, Karner A, Kostorz K, et al. SpiDa-MRI: behavioral and (f)MRI data of adults with fear of spiders. \emph{Sci Data}. 2025;12(1):284. PMID: 39962218. doi:10.1038/s41597-025-04569-w.

\bibitem{lahtee2026ru}
Lahtee Y. \emph{Readout Universe: A Philosophy and Logic for Grounding Claims}. Open Civil Science Initiative; 2026. ORCID: 0009-0005-3861-0626. Companion materials include the anchor--drive simulation code behind Fig.~\ref{fig:adm}.

\bibitem{lahtee2026rc}
Lahtee Y. The Readout Condition: distinguishability, access, and epistemic warrant. Open Civil Science Initiative; 2026. (Full paper held on file.)

\bibitem{lahtee2026birca}
Lahtee Y. birca: a safety-gated clinical-context specification (educational and research use only, CC BY-NC-SA 4.0); 2026. Cited for the discrete energy-balance rescue algebra only; no clinical content is imported and no status transfers to RRHM.

\bibitem{lahtee2026srct}
Lahtee Y. Systemic Repair Capacity Theory: health as coordinated repair capacity across layers. Preprint, Open Civil Science Initiative; 2026.

\bibitem{lahtee2026extprotocol}
Lahtee Y. Why we phobia? External-data readout protocol for RRHM. Companion protocol document; 2026.

\end{thebibliography}

\balance
\end{document}
